% La population mondiale — Géographie 1re Tech — Cours signé OnBuch+
% Programme officiel MINESEC. Compiler avec : tectonic N01-population-mondiale.tex
\newcommand{\DOCMATIERE}{Géographie}
\newcommand{\DOCNIVEAU}{1re Tech}
\newcommand{\DOCMODULE}{Géographie – classe de Première technique}
\newcommand{\DOCLECON}{N01}
\newcommand{\DOCTITRE}{La population mondiale}
% OnBuch+ — préambule « cours signés » ENRICHI (compilé avec tectonic / XeLaTeX)
% Système visuel « Pop » d'OnBuch+ : fond crème (jamais blanc), bordures encre
% épaisses, ombres « dures » décalées, titres Archivo Black en capitales.
% Version MATHÉMATIQUES : graphes (pgfplots/tikz), boîtes pédagogiques
% (définition, propriété, méthode, attention, le savais-tu, exercice résolu,
% exercice, TP, bilan), page de garde et signature OnBuch+ ancrée en bas.
%
% Macros attendues dans main.tex AVANT \input{preamble} :
%   \DOCMATIERE  \DOCNIVEAU  \DOCMODULE  \DOCLECON (numéro)  \DOCTITRE
\documentclass[11pt]{article}

\usepackage{fontspec}
\usepackage[french]{babel}
\usepackage[a4paper,margin=2cm,top=3.4cm,bottom=2.8cm,headheight=1.4cm,footskip=1.3cm]{geometry}
\usepackage{amsmath,amssymb,amsfonts}
\usepackage{mathtools} % \xmapsto, \coloneqq... souvent utilisés par les modèles
\usepackage{cancel} % simplifications barrées (bilans de forces)
% \abs{z} (module, valeur absolue) et \norm{u} (norme), utilisés par les modèles
\providecommand{\abs}{}\let\abs\relax\DeclarePairedDelimiter{\abs}{\lvert}{\rvert}
\providecommand{\norm}{}\let\norm\relax\DeclarePairedDelimiter{\norm}{\lVert}{\rVert}
\usepackage[table]{xcolor} % [table] : fournit aussi \rowcolors (alternance de lignes)
\usepackage{pagecolor}
\usepackage{tikz}
\usetikzlibrary{arrows.meta,positioning,calc,shapes.geometric,shapes.misc,shapes.arrows,decorations.pathmorphing,decorations.pathreplacing,decorations.markings,patterns,shadows,mindmap,trees,fit,backgrounds,matrix,intersections,angles,quotes}
\usepackage{pgfplots}
\usepackage[version=4]{mhchem} % équations nucléaires \ce{^{235}_{92}U}
\usepackage[european,siunitx]{circuitikz} % schémas électriques
\pgfplotsset{compat=1.18}
\usepgfplotslibrary{fillbetween,groupplots} % aires entre courbes (calcul intégral)
\usepackage{enumitem}
\usepackage{fancyhdr}
% [most] charge toutes les bibliothèques tcolorbox (listings, minted, raster,
% documentation...) et gonfle inutilement la mémoire TeX sur un document long
% (~300 boîtes) : on ne charge que ce qui est réellement utilisé.
\usepackage[skins,breakable]{tcolorbox}
\usepackage{graphicx}
\usepackage{colortbl}
\usepackage{booktabs}
\usepackage{tabularx}
\usepackage{multirow}
\usepackage{diagbox} % case d'en-tête diagonale des tableaux à double entrée
% Colonnes extensibles centrée / gauche / droite (C, L, R), souvent utilisées par les modèles
\newcolumntype{C}{>{\centering\arraybackslash}X}
\newcolumntype{L}{>{\raggedright\arraybackslash}X}
\newcolumntype{R}{>{\raggedleft\arraybackslash}X}
\usepackage{array}
\usepackage{multicol}
\usepackage{siunitx}
% parse-numbers=false : accepte aussi \SI{2,5 \times 10^{-3}}{...}
\sisetup{locale=FR,output-decimal-marker={,},group-minimum-digits=4,parse-numbers=false}
\DeclareSIUnit\ohm{\ensuremath{\symup{\Omega}}} % U+2126 absent de la police
\DeclareSIUnit\division{div} % réglages d'oscilloscope (ms/div, V/div)
\DeclareSIUnit\tr{tr} % tours (vitesse de rotation, ex. \SI{1500}{\tr\per\minute})
\DeclareSIUnit\molar{M} % concentration molaire (ex. \SI{3}{\molar}, \SI{1}{\micro\molar})
\DeclareSIUnit\an{an} % année (le modèle confond parfois avec \year, un compteur TeX) — ex. \SI{5}{\kilo\meter\per\an}
\DeclareSIUnit\FCFA{FCFA} % franc CFA (ex. \SI{500}{\FCFA\per\kilo\gram})
\DeclareSIUnit\degreeBrix{\textdegree Brix} % teneur en sucre d'un sirop/jus (réfractométrie)
\DeclareSIUnit\month{mois} % le modèle utilise parfois \month, un compteur TeX, comme unité de durée
\DeclareSIUnit\microliter{\micro\liter} % le modèle écrit parfois \microliter en un seul token
\DeclareSIUnit\million{millions} % dénombrement (ex. \SI{15}{\million\per\mL} spermatozoïdes)
\usepackage{qrcode}
% \degree : commande fréquemment utilisée par les modèles pour un angle ou une
% température en dehors de \SI{}{} (non fournie par les paquets chargés).
\providecommand{\degree}{^{\circ}}
% \qed (fin de démonstration), utilisé par les modèles sans amsthm
\providecommand{\qed}{\hfill$\square$}
% \barre : double barre des valeurs interdites dans un tableau de variations (tkz-tab)
\providecommand{\barre}{\ensuremath{\|}}
% environnement proof (amsthm non chargé) utilisé par les modèles
\ifdefined\proof\else\newenvironment{proof}[1][Démonstration]{\par\noindent\textit{#1.}\ }{\qed\par}\fi
\usepackage{lastpage}
\usepackage{hyperref}
\hypersetup{hidelinks}

% ============================================================
% Palette « Pop » OnBuch+ — mêmes hex que la classe P (plus_ui.dart)
% ============================================================
\definecolor{poporange}{HTML}{F59321}
\definecolor{popdark}{HTML}{E07A0C}
\definecolor{popcream}{HTML}{FFF6E8}
\definecolor{popink}{HTML}{1C1714}
\definecolor{poppurple}{HTML}{7A5AE0}
\definecolor{popgreen}{HTML}{2E9E6A}
\definecolor{poppink}{HTML}{E0457B}
\definecolor{popblue}{HTML}{2F7DE1}
\definecolor{popgold}{HTML}{C98A12}
\definecolor{popmuted}{HTML}{8A7F76}
\definecolor{popline}{HTML}{EADFD2}
\definecolor{popgray}{HTML}{8A7F76}
\colorlet{popgrey}{popgray}
% Alias tolérants (le modèle nomme parfois les couleurs différemment)
\colorlet{popwhite}{white}
\colorlet{popblack}{popink}
\colorlet{popred}{poppink}
\colorlet{popyellow}{popgold}
% Teintes claires (fonds de tableaux/encadrés) — le modèle nomme parfois
% les couleurs avec un suffixe L (ex. popblueL) sans les avoir définies.
\colorlet{poporangeL}{poporange!20}
\colorlet{popdarkL}{popdark!20}
\colorlet{poppurpleL}{poppurple!20}
\colorlet{popgreenL}{popgreen!20}
\colorlet{poppinkL}{poppink!20}
\colorlet{popblueL}{popblue!20}
\colorlet{popgoldL}{popgold!20}
\colorlet{popredL}{popred!20}
\colorlet{popgrayL}{popgray!20}
% Teintes claires pour fonds de tableaux / zones
\colorlet{popblueL}{popblue!10!white}
\colorlet{poporangeL}{poporange!14!white}
\colorlet{popgreenL}{popgreen!12!white}
\colorlet{poppurpleL}{poppurple!12!white}
\colorlet{poppinkL}{poppink!12!white}
\colorlet{popgoldL}{popgold!14!white}

\pagecolor{popcream}

% ============================================================
% Typographie
% ============================================================
\defaultfontfeatures{Ligatures=TeX}
\setmainfont{PlusJakartaSans-Medium.ttf}[Path=../../../fonts/,BoldFont=PlusJakartaSans-Bold.ttf]
% Maths en sans-serif (Fira Math) avec chiffres et lettres droites de Plus
% Jakarta, pour que les formules se fondent dans le texte.
\usepackage[math-style=ISO,bold-style=ISO]{unicode-math}
\setmathfont{FiraMath-Regular.otf}[Path=../../../fonts/]
\setmathfont{PlusJakartaSans-Medium.ttf}[Path=../../../fonts/,range={up/num,up/{Latin,latin},\mathup}]
\usepackage{icomma} % virgule décimale sans espace en mode math
% Caractères mathématiques Unicode absents des polices de texte (Plus Jakarta,
% Archivo Black) : rendus par leur équivalent mathématique, en texte comme en
% formule — sinon ils s'affichent en carré vide (ex. « Arithmétique dans ℤ »).
\usepackage{newunicodechar}
\newunicodechar{ℝ}{\ensuremath{\mathbb{R}}}
\newunicodechar{ℤ}{\ensuremath{\mathbb{Z}}}
\newunicodechar{ℂ}{\ensuremath{\mathbb{C}}}
\newunicodechar{ℕ}{\ensuremath{\mathbb{N}}}
\newunicodechar{ℚ}{\ensuremath{\mathbb{Q}}}
\newunicodechar{𝔻}{\ensuremath{\mathbb{D}}}
\newunicodechar{α}{\ensuremath{\alpha}}
\newunicodechar{β}{\ensuremath{\beta}}
\newunicodechar{γ}{\ensuremath{\gamma}}
\newunicodechar{δ}{\ensuremath{\delta}}
\newunicodechar{ε}{\ensuremath{\varepsilon}}
\newunicodechar{θ}{\ensuremath{\theta}}
\newunicodechar{λ}{\ensuremath{\lambda}}
\newunicodechar{μ}{\ensuremath{\mu}}
\newunicodechar{σ}{\ensuremath{\sigma}}
\newunicodechar{τ}{\ensuremath{\tau}}
\newunicodechar{φ}{\ensuremath{\varphi}}
\newunicodechar{ψ}{\ensuremath{\psi}}
\newunicodechar{ω}{\ensuremath{\omega}}
\newunicodechar{Σ}{\ensuremath{\Sigma}}
% majuscules produites par \MakeUppercase dans les titres (α-aminés -> Α)
\newunicodechar{Α}{\ensuremath{\alpha}}
\newunicodechar{Β}{\ensuremath{\beta}}
\newunicodechar{Γ}{\ensuremath{\Gamma}}
\newunicodechar{Δ}{\ensuremath{\Delta}}
\newunicodechar{Ε}{\ensuremath{\varepsilon}}
\newunicodechar{Θ}{\ensuremath{\Theta}}
\newunicodechar{Λ}{\ensuremath{\Lambda}}
\newunicodechar{Μ}{\ensuremath{\mu}}
\newunicodechar{Τ}{\ensuremath{\tau}}
\newunicodechar{Φ}{\ensuremath{\Phi}}
\newunicodechar{Ψ}{\ensuremath{\Psi}}
\newunicodechar{Ω}{\ensuremath{\Omega}}
\newunicodechar{↦}{\ensuremath{\mapsto}}
\newunicodechar{∈}{\ensuremath{\in}}
\newunicodechar{∉}{\ensuremath{\notin}}
\newunicodechar{∧}{\ensuremath{\wedge}}
\newunicodechar{⇔}{\ensuremath{\Leftrightarrow}}
\newunicodechar{⟺}{\ensuremath{\Longleftrightarrow}}
\newunicodechar{⇒}{\ensuremath{\Rightarrow}}
\newunicodechar{⟹}{\ensuremath{\Longrightarrow}}
\newunicodechar{∘}{\ensuremath{\circ}}
\newunicodechar{∪}{\ensuremath{\cup}}
\newunicodechar{∩}{\ensuremath{\cap}}
\newunicodechar{≡}{\ensuremath{\equiv}}
\newunicodechar{⊂}{\ensuremath{\subset}}
\newunicodechar{⊥}{\ensuremath{\perp}}
\newunicodechar{∃}{\ensuremath{\exists}}
\newunicodechar{∀}{\ensuremath{\forall}}
\newunicodechar{∛}{\ensuremath{\sqrt[3]{\,}}}
\newunicodechar{∜}{\ensuremath{\sqrt[4]{\,}}}
\newunicodechar{⃗}{\ensuremath{{}^{\to}}}
\newunicodechar{ⁿ}{\ifmmode^{n}\else\textsuperscript{\ensuremath{n}}\fi}
\newunicodechar{ˣ}{\ifmmode^{x}\else\textsuperscript{\ensuremath{x}}\fi}
\newunicodechar{ᵇ}{\ifmmode^{b}\else\textsuperscript{\ensuremath{b}}\fi}
\newunicodechar{ʳ}{\ifmmode^{r}\else\textsuperscript{\ensuremath{r}}\fi}
\newunicodechar{ᵘ}{\ifmmode^{u}\else\textsuperscript{\ensuremath{u}}\fi}
\newunicodechar{ʸ}{\ifmmode^{y}\else\textsuperscript{\ensuremath{y}}\fi}
\newunicodechar{ᵃ}{\ifmmode^{a}\else\textsuperscript{\ensuremath{a}}\fi}
\newunicodechar{ᵉ}{\ifmmode^{e}\else\textsuperscript{\ensuremath{e}}\fi}
\newunicodechar{ᵏ}{\ifmmode^{k}\else\textsuperscript{\ensuremath{k}}\fi}
\newunicodechar{ᵖ}{\ifmmode^{p}\else\textsuperscript{\ensuremath{p}}\fi}
\newunicodechar{ᴺ}{\ifmmode^{N}\else\textsuperscript{\ensuremath{N}}\fi}
\newunicodechar{ᶜ}{\ifmmode^{c}\else\textsuperscript{\ensuremath{c}}\fi}
\newunicodechar{ʰ}{\ifmmode^{h}\else\textsuperscript{\ensuremath{h}}\fi}
\newunicodechar{ᵠ}{\ifmmode^{q}\else\textsuperscript{\ensuremath{q}}\fi}
\newunicodechar{ₙ}{\ifmmode_{n}\else\textsubscript{\ensuremath{n}}\fi}
\newunicodechar{ₐ}{\ifmmode_{a}\else\textsubscript{\ensuremath{a}}\fi}
\newunicodechar{ᵢ}{\ifmmode_{i}\else\textsubscript{\ensuremath{i}}\fi}
\newunicodechar{ᵣ}{\ifmmode_{r}\else\textsubscript{\ensuremath{r}}\fi}
\newunicodechar{ₖ}{\ifmmode_{k}\else\textsubscript{\ensuremath{k}}\fi}
\newunicodechar{ᵦ}{\ifmmode_{\beta}\else\textsubscript{\ensuremath{\beta}}\fi}
\newunicodechar{ₓ}{\ifmmode_{x}\else\textsubscript{\ensuremath{x}}\fi}
\newfontfamily\popdisplay{ArchivoBlack-Regular.ttf}[Path=../../../fonts/]
\newfontfamily\poplogo{SpaceGrotesk-Bold.ttf}[Path=../../../fonts/]
\newfontfamily\popsemi{PlusJakartaSans-SemiBold.ttf}[Path=../../../fonts/]
\newfontfamily\popxbold{PlusJakartaSans-ExtraBold.ttf}[Path=../../../fonts/]
\newfontfamily\popxboldls{PlusJakartaSans-ExtraBold.ttf}[Path=../../../fonts/,LetterSpace=3.0]

\setlist[enumerate]{leftmargin=1.7em,itemsep=5pt,topsep=4pt}
\setlist[itemize]{leftmargin=1.7em,itemsep=4pt,topsep=4pt,label={\color{poporange}\textbullet}}
\setlength{\parindent}{0pt}
\setlength{\parskip}{5pt}
\renewcommand{\arraystretch}{1.25}

% pgfplots : style Pop par défaut
\pgfplotsset{
  every axis/.append style={axis line style={popink,line width=1pt},tick style={popink},
    grid style={popline,line width=0.5pt},label style={font=\small},tick label style={font=\footnotesize}},
  popcurve/.style={poporange,line width=1.8pt,no markers,smooth},
}
% Même style disponible dans un \draw TikZ simple (hors pgfplots)
\tikzset{popcurve/.style={poporange,line width=1.8pt}}
\tikzset{popgrid/.style={popline,very thin}}
\colorlet{popcurve}{poporange} % le nom du style est parfois utilisé comme couleur

% ============================================================
% Pill / squiggle
% ============================================================
\newcommand{\poppill}[3]{%
  \tikz[baseline=-0.55ex]\node[draw=popink,line width=1.2pt,rounded corners=2.6pt,%
    fill=#2,inner xsep=6.5pt,inner ysep=3pt]{%
    \popxboldls\fontsize{7}{7}\selectfont\color{#3}\MakeUppercase{#1}};%
}
\newcommand{\popsquiggle}[1]{%
  \tikz[baseline=-0.4ex]{
    \draw[#1,line width=2.6pt,line cap=round]
      plot [smooth] coordinates {(0,0.09) (0.34,0.01) (0.68,0.21) (1.02,0.01) (1.36,0.10)};
  }%
}
% Mot-clé surligné (marqueur orange sous le texte)
\newcommand{\cle}[1]{{\popxbold\color{popdark}#1}}

% ============================================================
% En-tête / pied
% ============================================================
\pagestyle{fancy}
\fancyhf{}
\renewcommand{\headrulewidth}{0pt}
\renewcommand{\footrulewidth}{0pt}
\lhead{\raisebox{-0.15ex}{\poplogo\fontsize{13}{13}\selectfont\color{popink}OnBuch\color{poporange}+}}
\rhead{\poppill{\DOCMATIERE\ \textbullet\ \DOCNIVEAU\ \textbullet\ Leçon \DOCLECON}{white}{popink}}
\lfoot{\popxbold\fontsize{7.6}{7.6}\selectfont\color{popmuted}\MakeUppercase{Cours signé OnBuch+}\ \textbullet\ {\popsemi\DOCTITRE}}
\rfoot{\tikz[baseline=-0.6ex]\node[rounded corners=8.5pt,fill=popink,minimum height=17pt,minimum width=17pt,inner xsep=5pt,inner sep=0]{\popxbold\fontsize{8}{8}\selectfont\color{popcream}\thepage\,/\,\pageref*{LastPage}};}

% ============================================================
% Titres
% ============================================================
\newcommand{\chaptitre}[2]{%
  \begin{tcolorbox}[enhanced,colback=poporange,colframe=popink,boxrule=2.6pt,arc=11pt,
      drop shadow=popink,left=15pt,right=15pt,top=12pt,bottom=11pt,before skip=3pt,after skip=18pt]
    \poppill{#2}{popink}{popcream}\par\vspace{9pt}
    {\raggedright\hyphenpenalty=10000\exhyphenpenalty=10000\popdisplay\fontsize{22}{24}\selectfont\color{popcream}\MakeUppercase{#1}\par}
  \end{tcolorbox}
}
% Section numérotée : pastille orange avec le numéro + titre Archivo
\newcounter{popsec}
\newcommand{\coursec}[1]{%
  \stepcounter{popsec}\setcounter{popsub}{0}%
  \par\vspace{16pt}\noindent
  \begin{minipage}{\linewidth}
  \tikz[baseline=-0.7ex]\node[draw=popink,line width=1.6pt,fill=poporange,rounded corners=4pt,minimum size=22pt,inner sep=0]{\popdisplay\fontsize{12}{12}\selectfont\color{popcream}\thepopsec};%
  \hspace{8pt}{\popdisplay\fontsize{15}{17}\selectfont\color{popink}\MakeUppercase{#1}}\par
  \vspace{4pt}\noindent{\color{poporange}\rule{36pt}{3.6pt}}
  \end{minipage}\par\nopagebreak\vspace{8pt}
}
\newcounter{popsub}
\newcommand{\courssub}[1]{%
  \stepcounter{popsub}%
  \par\vspace{9pt}\noindent{\popxbold\fontsize{11.5}{13}\selectfont\color{popink}{\color{poporange}\thepopsec.\thepopsub}\hspace{6pt}#1}\par\nopagebreak\vspace{4pt}
}

% ============================================================
% Boîtes pédagogiques — même recette : fond blanc, bordure épaisse colorée,
% ombre dure encre, badge plein. Titre optionnel [#1].
% ============================================================
\tcbset{popbase/.style={breakable,enhanced,colback=white,boxrule=2.3pt,arc=9pt,
  shadow={3pt}{-3pt}{0pt}{popink},left=14pt,right=14pt,top=11pt,bottom=11pt,before skip=11pt,after skip=15pt,
  fontupper=\popsemi\fontsize{10.6}{14.5}\selectfont\color{popink},
  fontlower=\popsemi\fontsize{10.6}{14.5}\selectfont\color{popink}}}
% Coche dessinée (le glyphe ✓ n'existe pas dans les polices du document)
\newcommand{\popcheck}[1]{\tikz[baseline=-0.5ex]\draw[#1,line width=1.6pt,line cap=round,line join=round](-0.13,0.0)--(-0.04,-0.09)--(0.13,0.1);}
\providecommand{\coche}{\popcheck{popgreen}}
\providecommand{\resultat}[1]{\par\smallskip\noindent\fcolorbox{popgreen}{popgreenL}{\parbox{\dimexpr\linewidth-2\fboxsep-2\fboxrule\relax}{\textbf{Résultat.} #1}}\par\smallskip}
\newcommand{\popboxhead}[3]{\poppill{#1}{#2}{white}\ifx\relax#3\relax\else\hspace{6pt}{\popxbold\fontsize{10.6}{12}\selectfont\color{popink}#3}\fi\par\vspace{7pt}}

\newtcolorbox{aretenir}[1][]{popbase,colframe=popgreen,before upper={\popboxhead{À retenir}{popgreen}{#1}}}
\newtcolorbox{exemplebox}[1][]{popbase,colframe=poporange,before upper={\popboxhead{Exemple}{poporange}{#1}}}
\newtcolorbox{definition}[1][]{popbase,colframe=popblue,before upper={\popboxhead{Définition}{popblue}{#1}}}
\newtcolorbox{propriete}[1][]{popbase,colframe=poppurple,before upper={\popboxhead{Propriété}{poppurple}{#1}}}
\newtcolorbox{methode}[1][]{popbase,colframe=popgold,before upper={\popboxhead{Méthode}{popgold}{#1}}}
\newtcolorbox{attention}[1][]{popbase,colframe=poppink,before upper={\popboxhead{Attention}{poppink}{#1}}}
\newtcolorbox{experience}[1][]{popbase,colframe=popblue,colback=popblueL,before upper={\popboxhead{Expérience}{popblue}{#1}}}
% « Le savais-tu ? » : carte violette pleine (contexte, culture, Cameroun)
\newtcolorbox{savaistu}[1][]{popbase,colback=poppurple,colframe=popink,
  fontupper=\popsemi\fontsize{10.4}{14.2}\selectfont\color{white},
  before upper={\poppill{Le savais-tu\,?}{white}{poppurple}\ifx\relax#1\relax\else\hspace{6pt}{\popxbold\color{white}#1}\fi\par\vspace{7pt}}}
% Exercice résolu : énoncé en haut, \tcblower puis la solution sur fond crème
\newcounter{popexo}\newcounter{popexr}
\newtcolorbox{exoresolu}[1][]{popbase,colframe=poporange,colbacklower=popcream,
  segmentation style={popink,line width=1.2pt,dashed},
  before upper={\stepcounter{popexr}\popboxhead{Exercice résolu \thepopexr}{poporange}{#1}},
  before lower={\poppill{Solution}{popink}{popcream}\par\vspace{6pt}}}
% Exercice (non résolu en place) — la correction est dans la partie Corrigés
\newtcolorbox{exercice}[1][]{popbase,colframe=popink,
  before upper={\stepcounter{popexo}\popboxhead{Exercice \thepopexo}{popink}{#1}}}
% Corrigé
\newtcolorbox{corrige}[1][]{popbase,colframe=popgreen,colback=popgreenL,
  before upper={\popboxhead{Corrigé}{popgreen}{#1}}}
% Objectifs / prérequis / bilan
\newtcolorbox{objectifs}{popbase,colframe=popink,colback=poporangeL,
  before upper={\popboxhead{Objectifs de la leçon}{popink}{}}}
\newtcolorbox{prerequis}{popbase,colframe=popink,colback=popblueL,
  before upper={\popboxhead{Prérequis}{popblue}{}}}
\newtcolorbox{bilan}{popbase,colframe=popink,colback=white,boxrule=2.8pt,
  before upper={\popboxhead{Fiche bilan}{poporange}{L'essentiel à connaître pour l'examen}}}
% Figure encadrée avec légende
\newtcolorbox{popfigure}[1][]{enhanced,breakable=false,colback=white,colframe=popline,boxrule=1.4pt,arc=8pt,
  left=10pt,right=10pt,top=10pt,bottom=8pt,before skip=10pt,after skip=12pt,halign=center,
  lower separated=false,halign lower=center,
  fontlower=\popsemi\fontsize{9}{11.5}\selectfont\color{popmuted},
  #1}
\newcommand{\legende}[1]{\par\vspace{6pt}{\popsemi\fontsize{9}{11.5}\selectfont\color{popmuted}#1}}

% ============================================================
% Signature OnBuch+
% ============================================================
\newcommand{\poplogoinline}[1]{{\poplogo\fontsize{#1}{#1}\selectfont\color{popink}OnBuch\color{poporange}+}}
\newcommand{\signatureonbuch}{%
  \begin{tcolorbox}[enhanced,colback=popink,colframe=popink,boxrule=2.2pt,arc=10pt,
      drop shadow=poporange,left=14pt,right=14pt,top=10pt,bottom=10pt,before skip=0pt,after skip=0pt]
    \tikz[baseline=-0.7ex]\node[circle,fill=popgreen,draw=popcream,line width=1.3pt,minimum size=22pt,inner sep=0]{\popcheck{white}};%
    \hspace{9pt}\begin{minipage}[c]{0.62\linewidth}
      {\popxbold\fontsize{12}{13}\selectfont\color{popcream}Signé\ \textbullet\ {\poplogo OnBuch\color{poporange}+}}\par\vspace{2pt}
      {\raggedright\popsemi\fontsize{8}{10.5}\selectfont\color{popline}\MakeUppercase{Équipe pédagogique OnBuch+}\\\MakeUppercase{Conforme au programme officiel MINESEC}\par}
    \end{minipage}\hfill
    {\poplogo\fontsize{22}{22}\selectfont\color{popcream}OnBuch\color{poporange}+}
  \end{tcolorbox}%
}

% ============================================================
% Page de garde
% #1 titre · #2 matière · #3 niveau · #4 module · #5 n° leçon · #6 durée ·
% #7 description / objectifs (texte ou itemize)
% ============================================================
\newcommand{\pagegarde}[7]{%
  \thispagestyle{empty}
  \newgeometry{a4paper,margin=1.7cm,top=1.6cm,bottom=1.4cm}%
  \begin{tikzpicture}[remember picture,overlay]
    \fill[poporange!18] ([xshift=-3.2cm,yshift=-3.6cm]current page.north east) circle (4.6cm);
    \fill[poppurple!14] ([xshift=2.4cm,yshift=7.6cm]current page.south west) circle (3.8cm);
  \end{tikzpicture}%
  {\poplogo\fontsize{30}{30}\selectfont\color{popink}OnBuch\color{poporange}+}\hfill
  \raisebox{0.6ex}{\poppill{Cours signé \textbullet\ Premium}{popink}{popcream}}\par
  \vspace{1pt}\popsquiggle{poppurple}\par
  \vspace{8pt}
  \begin{tcolorbox}[enhanced,colback=poporange,colframe=popink,boxrule=2.8pt,arc=13pt,
      drop shadow=popink,left=18pt,right=18pt,top=12pt,bottom=13pt,after skip=10pt]
    \poppill{#2 \textbullet\ #3}{popink}{popcream}\hspace{5pt}\poppill{#4}{white}{popink}\par\vspace{8pt}
    {\popdisplay\fontsize{14}{14}\selectfont\color{popink}LEÇON #5}\par\vspace{4pt}
    {\raggedright\hyphenpenalty=10000\exhyphenpenalty=10000\popdisplay\fontsize{25}{27}\selectfont\color{popcream}\MakeUppercase{#1}\par}
    \vspace{8pt}
    {\popxbold\fontsize{9.5}{9.5}\selectfont\color{popink}\MakeUppercase{Durée conseillée : #6}}
  \end{tcolorbox}
  \begin{tcolorbox}[enhanced,colback=white,colframe=popink,boxrule=2.2pt,arc=10pt,
      drop shadow=popink,left=16pt,right=16pt,top=10pt,bottom=10pt,after skip=8pt]
    \poppill{Dans cette leçon}{poppurple}{white}\par\vspace{5pt}
    \popsemi\fontsize{9.3}{12.6}\selectfont\color{popink}#7
  \end{tcolorbox}
  \noindent\begin{minipage}[t]{0.487\linewidth}
    \begin{tcolorbox}[enhanced,colback=popgreen,colframe=popink,boxrule=2.2pt,arc=10pt,
        drop shadow=popink,left=11pt,right=11pt,top=10pt,bottom=10pt,height=3.1cm,valign=center]
      \tcbox[colback=white,colframe=popink,boxrule=1.3pt,arc=5pt,boxsep=0pt,nobeforeafter,
        left=3pt,right=3pt,top=3pt,bottom=3pt,valign=center]{\qrcode[height=1.9cm]{https://play.google.com/store/apps/details?id=com.luvvix.onbuch&hl=fr}}%
      \hspace{8pt}\begin{minipage}[c]{0.5\linewidth}
        {\popdisplay\fontsize{11}{12.5}\selectfont\color{white}\MakeUppercase{Télécharge OnBuch}}\par\vspace{4pt}
        {\raggedright\popsemi\fontsize{8.4}{11}\selectfont\color{white}Gratuit \textbullet\ Google Play\\Tuteur IA, annales, concours, résultats\par}
      \end{minipage}
    \end{tcolorbox}
  \end{minipage}\hfill
  \begin{minipage}[t]{0.487\linewidth}
    \begin{tcolorbox}[enhanced,colback=poppurple,colframe=popink,boxrule=2.2pt,arc=10pt,
        drop shadow=popink,left=11pt,right=11pt,top=10pt,bottom=10pt,height=3.1cm,valign=center]
      \tikz[baseline=-0.7ex]\node[circle,fill=white,draw=popink,line width=1.3pt,minimum size=1.6cm,inner sep=0]{\popdisplay\fontsize{24}{24}\selectfont\color{poppurple}+};%
      \hspace{8pt}\begin{minipage}[c]{0.6\linewidth}
        {\popdisplay\fontsize{11}{12.5}\selectfont\color{white}\MakeUppercase{Passe à OnBuch+}}\par\vspace{4pt}
        {\raggedright\popsemi\fontsize{8.4}{11}\selectfont\color{white}Tous les cours signés, Tuteur IA illimité, fiches PDF — onglet OnBuch+ de l'app\par}
      \end{minipage}
    \end{tcolorbox}
  \end{minipage}
  \vspace{8pt}
  \popcontenu
  \vfill
  \signatureonbuch
  \restoregeometry
  \newpage
}

% Rangée « Ce cours contient » (page de garde)
\newcommand{\poptile}[3]{% #1 couleur #2 gros texte #3 légende
  \begin{tcolorbox}[enhanced,colback=#1,colframe=popink,boxrule=2pt,arc=9pt,shadow={2.5pt}{-2.5pt}{0pt}{popink},
      left=6pt,right=6pt,top=9pt,bottom=9pt,halign=center,valign=center,height=1.75cm,width=\linewidth,nobeforeafter]
    {\popdisplay\fontsize{12.5}{13}\selectfont\color{white}\MakeUppercase{#2}}\par\vspace{3pt}
    {\centering\popsemi\fontsize{7.8}{9.5}\selectfont\color{white}#3\par}
  \end{tcolorbox}}
\newcommand{\popcontenu}{%
  \par\noindent{\popxboldls\fontsize{7.5}{7.5}\selectfont\color{popmuted}CE COURS SIGNÉ CONTIENT}\par\vspace{7pt}
  \noindent\begin{minipage}[t]{0.235\linewidth}\poptile{popblue}{Cours}{complet, illustré, pas à pas}\end{minipage}\hfill
  \begin{minipage}[t]{0.235\linewidth}\poptile{poporange}{Méthodes}{et exercices résolus}\end{minipage}\hfill
  \begin{minipage}[t]{0.235\linewidth}\poptile{poppink}{Exercices}{type examen + corrigés}\end{minipage}\hfill
  \begin{minipage}[t]{0.235\linewidth}\poptile{popgreen}{Fiche bilan}{l'essentiel à retenir}\end{minipage}\par}

% Sommaire
\newtcolorbox{sommaire}{enhanced,colback=white,colframe=popink,boxrule=2.2pt,arc=10pt,
  drop shadow=popink,left=18pt,right=18pt,top=13pt,bottom=13pt,before skip=6pt,after skip=20pt,
  before upper={\poppill{Sommaire}{poppurple}{white}\par\vspace{9pt}\popsemi\fontsize{10.6}{14.5}\selectfont\color{popink}}}

% Signature de fin de cours — poussée tout en bas de la dernière page
\newcommand{\findecours}{%
  \par\vfill\nopagebreak
  \begin{center}\popsquiggle{poporange}\end{center}\vspace{4pt}
  \signatureonbuch
}
\AtBeginDocument{\def\llbracket{\mathopen{[\mkern-3.5mu[}}\def\rrbracket{\mathclose{]\mkern-3.5mu]}}} % intervalles d'entiers (glyphe ⟦ absent de la police)
% Glyphes absents de Fira Math : redéfinis à partir de glyphes disponibles
\AtBeginDocument{%
  \def\setminus{\mathbin{\backslash}}\let\smallsetminus\setminus
  \def\vdots{\mathord{\vbox{\baselineskip3.2pt\lineskiplimit0pt\kern4pt\hbox{.}\hbox{.}\hbox{.}}}}%
  \def\ddots{\mathinner{\mkern1mu\raise6.5pt\hbox{.}\mkern2mu\raise3.6pt\hbox{.}\mkern2mu\raise0.7pt\hbox{.}\mkern1mu}}%
  \def\triangle{\mathord{\scalebox{0.9}{$\symup{\Delta}$}}}%
  \def\nabla{\mathord{\raisebox{\depth}{\scalebox{1}[-1]{$\symup{\Delta}$}}}}%
  \def\checkmark{\popcheck{popdark}}%
  \def\star{\mathbin{\ast}}\def\bigstar{\mathord{\ast}}%
  \def\diamond{\mathbin{\scalebox{0.6}{$\Diamond$}}}%
  \def\bigcup{\mathop{\mathchoice{\raisebox{-0.3em}{\scalebox{1.7}{$\cup$}}}{\scalebox{1.25}{$\cup$}}{\cup}{\cup}}\displaylimits}%
  \def\bigcap{\mathop{\mathchoice{\raisebox{-0.3em}{\scalebox{1.7}{$\cap$}}}{\scalebox{1.25}{$\cap$}}{\cap}{\cap}}\displaylimits}%
  \def\lesseqqgtr{\mathrel{\lessgtr}}\def\bigoplus{\mathop{\oplus}\displaylimits}%
}
\providecommand{\ssi}{si et seulement si} % abréviation fréquente
\AtBeginDocument{\providecommand{\wideparen}{\overparen}} % arc de cercle

\begin{document}

\pagegarde{La population mondiale}{Géographie}{1re Tech}{Géographie – classe de Première technique}{N01}{3 séances (fiche de progression harmonisée)}{Cette leçon analyse la dynamique de la population mondiale à travers ses indicateurs, ses déséquilibres et les politiques mises en place pour les corriger, en s'appuyant sur des études de cas camerounaises et internationales. Elle développe les compétences cartographiques, graphiques et argumentatives exigées au Baccalauréat technique.
\vspace{4pt}
\begin{multicols}{2}\small
\begin{itemize}
\item À la fin de cette leçon, je sais définir et calculer les principaux indicateurs démographiques (taux de natalité, mortalité, fécondité, accroissement naturel, indice de migration).
\item À la fin de cette leçon, je sais expliquer le modèle de la transition démographique et situer n'importe quel pays dans l'une de ses phases à partir de ses indicateurs.
\item À la fin de cette leçon, je sais identifier et caractériser les grands problèmes démographiques mondiaux : surpopulation, sous-population, vieillissement, transition inachevée et dividende démographique.
\item À la fin de cette leçon, je sais analyser et comparer des politiques démographiques (pro-natalistes, anti-natalistes, migratoires) à travers des exemples précis (Chine, France, Suède, Cameroun).
\item À la fin de cette leçon, je sais construire et interpréter une pyramide des âges, un croquis de répartition de la population et un diagramme d'évolution démographique.
\item À la fin de cette leçon, je sais rédiger une réponse argumentée justifiant les choix démographiques d'un État en mobilisant des données chiffrées et des documents cartographiques.
\end{itemize}\end{multicols}}

\begin{sommaire}
\begin{multicols}{2}
\begin{enumerate}
\item 1. Les outils du démographe : indicateurs et transition démographique
\item 2. Les grands problèmes démographiques mondiaux : diagnostics contrastés
\item 3. Dossier 1 : Les politiques démographiques – Types, acteurs et bilans
\item 4. Focus Cameroun : Dynamiques internes, défis et prospective 2035
\item 5. Atelier méthodologique : Cartographie, graphiques et analyse de documents (TP)
\item Activité d'intégration
\item Méthodes et exercices résolus
\item Exercices
\item Corrigés des exercices
\item Fiche bilan
\end{enumerate}
\end{multicols}
\end{sommaire}

\begin{objectifs}
\begin{itemize}
\item À la fin de cette leçon, je sais définir et calculer les principaux indicateurs démographiques (taux de natalité, mortalité, fécondité, accroissement naturel, indice de migration).
\item À la fin de cette leçon, je sais expliquer le modèle de la transition démographique et situer n'importe quel pays dans l'une de ses phases à partir de ses indicateurs.
\item À la fin de cette leçon, je sais identifier et caractériser les grands problèmes démographiques mondiaux : surpopulation, sous-population, vieillissement, transition inachevée et dividende démographique.
\item À la fin de cette leçon, je sais analyser et comparer des politiques démographiques (pro-natalistes, anti-natalistes, migratoires) à travers des exemples précis (Chine, France, Suède, Cameroun).
\item À la fin de cette leçon, je sais construire et interpréter une pyramide des âges, un croquis de répartition de la population et un diagramme d'évolution démographique.
\item À la fin de cette leçon, je sais rédiger une réponse argumentée justifiant les choix démographiques d'un État en mobilisant des données chiffrées et des documents cartographiques.
\item À la fin de cette leçon, je sais évaluer les défis démographiques actuels du Cameroun (jeunesse, urbanisation, pression sur les ressources) et proposer des pistes d'action cohérentes.
\end{itemize}
\end{objectifs}

\begin{prerequis}
\begin{itemize}
\item Notions de base sur la population mondiale (effectifs, répartition inégale) vues en 4ème.
\item Lecture et interprétation de graphiques simples (courbes, histogrammes) et de cartes choroplèthes.
\item Calcul de pourcentages, de taux pour 1000 et de densités de population.
\item Localisation des grands ensembles géographiques et du Cameroun sur une planisphère.
\item Vocabulaire de base : natalité, mortalité, migration, densité, urbanisation.
\end{itemize}
\end{prerequis}

\newpage

\coursec{Les outils du démographe : indicateurs et transition démographique}

\textbf{Situation d'accroche.} Dans la cour d'un lycée de Ngaoundéré, 300 élèves s'entassent dans des salles prévues pour 50 ; au même moment, un village du Mayo-Rey voit son école fermer faute d'enfants. Pendant que le Cameroun compte environ 28 millions d'habitants avec un âge médian d'environ 18 ans, le Japon vieillit et la Chine vient de perdre sa place de pays le plus peuplé au profit de l'Inde. Comment expliquer ces contrastes saisissants et quelles politiques les États mettent-ils en œuvre pour gérer leur population ?

\textbf{Problématique.} Pourquoi les populations du monde ne croissent-elles pas au même rythme, et comment les démographes mesurent-ils et expliquent-ils ces différences ?

\courssub{Les indicateurs bruts : mesurer la dynamique d'une population}

Pour comparer les populations, il ne suffit pas de compter les habitants : il faut des \cle{indicateurs} qui rapportent les naissances et les décès à la taille de la population. L'intuition est simple : dire qu'il y a eu environ 900 000 naissances au Cameroun ne veut rien dire si l'on ignore que le pays compte environ 28,6 millions d'habitants. On calcule donc des \cle{taux}, généralement exprimés pour mille habitants (symbole ‰, à ne pas confondre avec le pourcentage \%).

\begin{definition}[Taux de natalité (TN)]
Le \cle{taux de natalité} est le nombre de naissances vivantes pour 1000 habitants par an.
\[ TN = \frac{\text{nombre de naissances vivantes dans l'année}}{\text{population moyenne de l'année}} \times 1000 \quad (\text{en ‰}) \]
\end{definition}

\begin{definition}[Taux de mortalité (TM) et accroissement naturel]
Le \cle{taux de mortalité} est le nombre de décès pour 1000 habitants par an. L'\cle{accroissement naturel} (ou solde naturel) est la différence entre les naissances et les décès :
\[ AN = TN - TM \]
Exprimé en pourcentage, il donne le \cle{taux d'accroissement naturel}. Il ne tient pas compte des migrations : la croissance totale d'une population ajoute le \cle{solde migratoire} (différence entre les immigrations et les émigrations).
\end{definition}

\begin{definition}[Indice conjoncturel de fécondité (ICF)]
L'\cle{indice conjoncturel de fécondité} est le nombre moyen d'enfants qu'une femme aurait au cours de sa vie féconde si elle subissait les taux de fécondité par âge de l'année considérée. Le \cle{seuil de renouvellement des générations} est d'environ 2,1 enfants par femme : en dessous, une population finit par décliner (sans immigration).
\end{definition}

\begin{definition}[Mortalité infantile et espérance de vie]
Le \cle{taux de mortalité infantile} (TMI) est le nombre de décès d'enfants de moins d'un an pour 1000 naissances vivantes : c'est un excellent indicateur du niveau sanitaire d'un pays. L'\cle{espérance de vie à la naissance} (EV) est le nombre moyen d'années qu'un nouveau-né peut espérer vivre si les conditions de mortalité de l'année se maintiennent.
\end{definition}

\begin{exemplebox}[Calculer un taux d'accroissement naturel]
Cameroun, 2023 : $TN \approx 35,4$‰ et $TM \approx 8,7$‰. Donc $AN = 35,4 - 8,7 = 26,7$‰, soit un taux d'accroissement naturel de $2,67\,\%$ par an. À ce rythme, la population \emph{double en moins de 30 ans} (règle pratique : temps de doublement $\approx 70 / \text{taux en \%}$, soit $70/2,67 \approx 26$ ans).
\end{exemplebox}

\begin{attention}[Ne pas confondre population totale et densité]
Erreur fréquente : confondre la \cle{densité de population} (habitants par km\textsuperscript{2}) et la \cle{population totale} (effectif). Un pays peut être très peuplé mais peu dense, ou inversement : la Chine est le deuxième pays le plus peuplé du monde mais n'est pas le plus dense ; Monaco est minuscule en effectif mais est l'État le plus dense du monde ; la Mongolie est très peu dense malgré sa superficie immense. De même, un taux de natalité élevé ne signifie pas forcément un grand nombre de naissances en valeur absolue.
\end{attention}

\courssub{Taux bruts et taux standardisés : les limites des comparaisons}

Les taux présentés ci-dessus sont des \cle{taux bruts} : ils portent sur l'ensemble de la population sans tenir compte de sa structure par âge. Or cette structure fausse les comparaisons. Intuition : un pays jeune comme le Cameroun (âge médian d'environ 18 ans) compte peu de personnes âgées, donc mécaniquement peu de décès : son taux brut de mortalité est faible même si les conditions sanitaires restent difficiles. À l'inverse, le Japon, très vieilli, affiche un taux brut de mortalité élevé malgré une excellente médecine.

\begin{definition}[Taux standardisé]
Un \cle{taux standardisé} est un taux recalculé en appliquant les taux par âge d'un pays à une structure d'âge de référence commune. Il permet de comparer réellement les niveaux de mortalité ou de fécondité entre pays aux pyramides des âges différentes.
\end{definition}

\begin{exemplebox}[Le paradoxe de la mortalité]
En 2023, le taux brut de mortalité de la France (environ $9,4$‰) est \emph{supérieur} à celui du Cameroun (environ $8,7$‰). Faut-il en conclure que l'on meurt davantage en France ? Non : la France compte plus de 20\,\% de personnes de 65 ans et plus, contre environ 3\,\% au Cameroun. Standardisée sur la même structure d'âge, la mortalité camerounaise apparaît bien plus élevée — ce que confirment l'espérance de vie (environ 62 ans contre 82,5 ans) et la mortalité infantile (environ 48‰ au Cameroun contre moins de 4‰ en France).
\end{exemplebox}

\begin{aretenir}[Les limites des taux bruts]
Les taux bruts dépendent de la structure par âge : ils ne permettent pas de comparer directement deux pays très différents. Pour des comparaisons rigoureuses, on utilise l'ICF, la mortalité infantile, l'espérance de vie ou des taux standardisés.
\end{aretenir}

\courssub{Le modèle de la transition démographique}

Comment passer d'une natalité de 35‰ à une natalité de 10‰ ? Les démographes Adolphe Landry et Frank Notestein ont proposé un modèle général, construit à partir de l'histoire européenne puis vérifié à l'échelle mondiale.

\begin{propriete}[Modèle de transition démographique (Landry, Notestein)]
Toute société passe d'un \cle{équilibre ancien} (natalité forte, mortalité forte, croissance quasi nulle) à un \cle{équilibre moderne} (natalité faible, mortalité faible, croissance faible ou nulle), en traversant une phase intermédiaire de \cle{forte croissance} où la mortalité baisse avant la natalité. Ce modèle, dont la connaissance est exigible, se décompose en phases :
\begin{enumerate}
\item \textbf{Phase 1 — pré-transition} : TN et TM élevés (30 à 40‰ et plus), croissance très faible. Régime démographique traditionnel.
\item \textbf{Phase 2 — transition débutante} : la mortalité chute (révolution sanitaire) mais la natalité reste élevée : explosion démographique.
\item \textbf{Phase 3 — transition avancée} : la natalité baisse à son tour (éducation, planning familial), la croissance ralentit.
\item \textbf{Phase 4 — post-transition} : TN et TM faibles et proches, croissance quasi nulle, ICF proche ou inférieur à 2,1.
\item \textbf{Phase 5 — vieillissement} (ajoutée par certains auteurs) : la natalité passe sous la mortalité, la population décline et vieillit (Japon, Allemagne, Italie).
\end{enumerate}
\end{propriete}

\begin{popfigure}
\begin{tikzpicture}
\begin{axis}[width=12cm,height=7cm,legend pos=north west,xlabel=Temps,ylabel=Taux (‰),ytick={0,10,20,30,40,50},xtick={1,2,3,4,5},xticklabels={Phase 1,Phase 2,Phase 3,Phase 4,Phase 5}]
\addplot[thick,red] coordinates {(1,40) (2,40) (3,30) (4,12) (5,10)};
\addlegendentry{Natalité}
\addplot[thick,blue] coordinates {(1,40) (2,20) (3,15) (4,10) (5,10)};
\addlegendentry{Mortalité}
\addplot[thick,green!60!black] coordinates {(1,0) (2,20) (3,15) (4,2) (5,0)};
\addlegendentry{Accroissement naturel}
\end{axis}
\end{tikzpicture}
\legende{Le modèle de la transition démographique : la mortalité baisse d'abord (phase 2), la natalité ensuite (phase 3) ; l'écart entre les deux courbes mesure la croissance naturelle, maximale en phase 2.}
\end{popfigure}

\courssub{Pourquoi la mortalité baisse-t-elle d'abord, puis la natalité ?}

La \cle{baisse de la mortalité} s'explique par la \cle{révolution sanitaire} : vaccinations, antibiotiques, hygiène, accès à l'eau potable, progrès de l'alimentation et de la médecine. Ces progrès se diffusent vite car ils ne demandent pas de changer les mentalités : une campagne de vaccination protège un enfant en quelques minutes.

La \cle{baisse de la natalité} est plus lente car elle touche aux comportements et aux valeurs. Ses facteurs principaux sont :
\begin{itemize}
\item l'\cle{éducation des filles} : une femme scolarisée se marie plus tard, connaît la contraception et peut travailler ;
\item le \cle{coût des enfants} : en ville, un enfant doit être nourri, logé, scolarisé longtemps avant de rapporter un revenu, alors qu'à la campagne il aide tôt aux champs ;
\item le \cle{planning familial} et la diffusion de la contraception ;
\item l'\cle{urbanisation} : la vie urbaine réduit la place des familles nombreuses ;
\item la baisse de la mortalité infantile elle-même : quand les enfants survivent, les parents en font moins.
\end{itemize}

\begin{savaistu}[La transition la plus rapide de l'histoire]
L'Europe a mis près de deux siècles (1750--1950) pour achever sa transition démographique. La Corée du Sud n'a mis que quelques décennies : son ICF est passé d'environ 6 enfants par femme en 1960 à moins de 1 aujourd'hui, le plus bas du monde. Les transitions contemporaines sont donc beaucoup plus rapides, mais aussi plus brutales à gérer pour les États.
\end{savaistu}

\courssub{Situer les pays dans la transition : méthode et applications}

\begin{methode}[Situer un pays dans la transition démographique]
\begin{enumerate}
\item Regarder le TN et le TM (en ‰).
\item Calculer l'accroissement naturel : $AN = TN - TM$.
\item Observer l'ICF et l'espérance de vie.
\item Comparer aux seuils des phases : $TN > 30$‰ et $TM > 20$‰ = phase 1 ; $TN > 30$‰ et $TM < 15$‰ = phase 2 ; TN en baisse nette (15 à 30‰) = phase 3 ; TN et TM faibles et proches ($\approx 10$‰) = phase 4 ; $TM > TN$ = phase 5.
\end{enumerate}
\end{methode}

\begin{exemplebox}[Quatre pays, quatre phases (données 2023, INS et ONU)]
\begin{center}
\begin{tabularx}{\linewidth}{|l|X|X|X|X|}
\hline
\rowcolor{popblueL} \textbf{Indicateur} & \textbf{Cameroun} & \textbf{Niger} & \textbf{France} & \textbf{Japon} \\
\hline
Population (M) & 28,6 & 27 & 68 & 124 \\
\hline
TN (‰) & 35,4 & 45 & 10,2 & 6,3 \\
\hline
TM (‰) & 8,7 & 9 & 9,4 & 12,9 \\
\hline
AN (\%) & 2,67 & 3,6 & 0,08 & $-0,66$ \\
\hline
ICF (enfants/femme) & 4,6 & 6,7 & 1,68 & 1,26 \\
\hline
EV (années) & 62,4 & 62 & 82,5 & 84,8 \\
\hline
Phase de transition & 2 & 1--2 & 4 & 5 \\
\hline
\end{tabularx}
\end{center}
\end{exemplebox}

Le \textbf{Cameroun} se situe clairement en \cle{phase 2} (transition débutante) : la mortalité a fortement baissé depuis les années 1960 (grâce aux vaccinations, à l'hygiène, aux moustiquaires imprégnées), mais la natalité reste élevée (environ $35$‰) et l'ICF de 4,6 enfants par femme témoigne d'une fécondité encore forte, surtout en milieu rural et dans les régions du Nord. Sa pyramide des âges l'illustre parfaitement.

\begin{popfigure}
\begin{tikzpicture}[x=0.32cm,y=0.75cm]
% Pyramide schématique Cameroun (barres horizontales)
\foreach \y/\w in {0/8.2,1/7.0,2/5.6,3/4.2,4/2.8,5/1.6,6/0.7}{
  \fill[popblue!70] (-\w,\y) rectangle (0,\y+0.8);
  \fill[poppink!70] (0,\y) rectangle (\w,\y+0.8);
}
\draw[->,thick] (-9,0) -- (9,0);
\draw[->,thick] (0,0) -- (0,7.6) node[above]{\footnotesize Âge};
\node[left] at (-9,0.4) {\footnotesize Hommes};
\node[right] at (9,0.4) {\footnotesize Femmes};
\node[below] at (0,-0.2) {\footnotesize 0\,\% \hspace{2.2cm} part de la population \hspace{2.2cm} 0\,\%};
\node[left,font=\footnotesize] at (-8.6,0.4) {0--14 ans};
\node[left,font=\footnotesize] at (-8.6,6.4) {65 ans et +};
\draw[<-,poporange,very thick] (4,0.6) -- (10,2.2) node[right,popdark,align=left,font=\footnotesize] {Base très large :\\ plus de 42\,\% de moins de 15 ans};
\draw[<-,poporange,very thick] (0.5,6.4) -- (10,6.2) node[right,popdark,align=left,font=\footnotesize] {Sommet étroit :\\ faible espérance de vie (62 ans)};
\end{tikzpicture}
\legende{Pyramide des âges schématique du Cameroun (2023, d'après INS/ONU) : base très large (jeunesse) et sommet étroit, profil typique de la phase 2 de la transition.}
\end{popfigure}

Le \textbf{Niger} est en fin de phase 1 / début de phase 2 : mortalité encore élevée, ICF record de près de 7 enfants par femme. La \textbf{France} est en phase 4 : natalité et mortalité faibles et proches, croissance naturelle quasi nulle ($0,08\,\%$). Le \textbf{Japon} est entré en phase 5 : sa mortalité dépasse sa natalité, sa population diminue et vieillit fortement.

\begin{definition}[Transition inachevée]
On parle de \cle{transition inachevée} lorsque la mortalité a fortement baissé mais que la natalité reste élevée : l'écart entre les deux entretient une croissance rapide et persistante. C'est la situation de nombreux pays africains, dont le Cameroun : la population double en une génération, ce qui exige chaque année davantage d'écoles, de soins, d'emplois et de logements.
\end{definition}

\begin{exoresolu}[Situer un pays dans la transition]
Énoncé. Un pays A affiche : $TN = 38$‰, $TM = 12$‰, $ICF = 5,2$, $EV = 58$ ans. Un pays B affiche : $TN = 11$‰, $TM = 11,5$‰, $ICF = 1,4$, $EV = 83$ ans. (1) Calculer l'accroissement naturel de chaque pays. (2) Situer chaque pays dans la transition démographique en justifiant.
\tcblower
Solution détaillée. (1) Pays A : $AN = 38 - 12 = 26$‰, soit $2,6\,\%$ par an (croissance très forte). Pays B : $AN = 11 - 11,5 = -0,5$‰, soit $-0,05\,\%$ (décroissance naturelle). (2) Pays A : natalité élevée ($> 30$‰) et mortalité déjà abaissée ($< 15$‰), ICF de 5,2 et EV faible : il est en \textbf{phase 2}, transition débutante, avec une transition inachevée (mortalité baissée, natalité restée forte). Pays B : natalité et mortalité faibles, mortalité supérieure à la natalité, ICF très inférieur au seuil de renouvellement (2,1) : il est en \textbf{phase 5}, phase de vieillissement et de déclin démographique.
\end{exoresolu}

\begin{exercice}[À vous de jouer]
À partir des données du tableau comparatif (Cameroun, Niger, France, Japon) : (1) calculez le temps de doublement approximatif de la population du Niger ; (2) expliquez pourquoi le taux brut de mortalité du Japon est plus élevé que celui du Cameroun ; (3) rédigez un paragraphe de 8 à 10 lignes situant le Cameroun dans la transition démographique et justifiant votre réponse par au moins trois indicateurs.
\end{exercice}

\begin{corrige}[Exercice 1]
(1) $AN_{\text{Niger}} = 45 - 9 = 36$‰ $= 3,6\,\%$ ; temps de doublement $\approx 70 / 3,6 \approx 19$ ans. (2) Le Japon a une population très âgée (plus de 29\,\% de 65 ans et plus) : les décès sont nombreux en proportion, d'où un TM brut élevé malgré une excellente médecine ; le Cameroun a une population très jeune, donc peu de décès en proportion — c'est la limite des taux bruts, qui dépendent de la structure d'âge. (3) Le Cameroun est en phase 2 (transition débutante) : sa mortalité a fortement reculé ($TM \approx 8,7$‰, EV de 62,4 ans en hausse) grâce à la révolution sanitaire, mais sa natalité demeure élevée ($TN \approx 35,4$‰, ICF de 4,6). L'écart produit un accroissement naturel de $2,67\,\%$ par an et une pyramide à base très large : la transition est engagée mais inachevée.
\end{corrige}

\begin{aretenir}[L'essentiel de la section]
Les démographes mesurent la dynamique des populations avec des taux (natalité, mortalité, accroissement naturel, ICF, mortalité infantile, espérance de vie). Les taux bruts dépendent de la structure d'âge : prudence dans les comparaisons. Le modèle de la transition démographique décrit le passage d'un régime ancien à un régime moderne en 4 ou 5 phases ; la mortalité baisse d'abord (révolution sanitaire), la natalité ensuite (éducation des filles, urbanisation, planning familial). Le Cameroun est en phase 2 : transition inachevée et croissance rapide.
\end{aretenir}

\coursec{Les grands problèmes démographiques mondiaux : diagnostics contrastés}

Dans la section précédente, nous avons appris à manier les \cle{outils du démographe} : taux de natalité, de mortalité, indice conjoncturel de fécondité (ICF), taux d'accroissement naturel, et nous avons décrit le modèle de la \cle{transition démographique}. Ces instruments ne sont pas de simples statistiques : ils servent à poser des \textbf{diagnostics}. Or, à l'échelle mondiale, les diagnostics sont profondément \emph{contrastés} : certains pays s'inquiètent d'avoir « trop d'enfants », d'autres de ne plus en avoir assez ; certains territoires étouffent sous la densité, d'autres restent presque vides. Cette section dresse la typologie des grands problèmes démographiques mondiaux et fournit les clés pour les analyser sans caricature.

\courssub{Une répartition mondiale profondément inégale}

L'humanité a franchi le cap des \textbf{8 milliards d'habitants en novembre 2022}. Mais cette population est répartie de manière très inégale entre les continents, et sa croissance — environ \num{0,9} \% par an, en \emph{ralentissement} — ne se répartit pas non plus équitablement.

\begin{popfigure}
\begin{center}
\begin{tikzpicture}
\begin{axis}[width=\linewidth,height=7cm,ybar,bar width=12pt,
ylabel={Part dans la population mondiale (\%)},
symbolic x coords={Asie,Afrique,Europe,Amérique N,Amérique S,Océanie},
xtick=data,ymin=0,ymax=70,enlargelimits=0.15,
nodes near coords,nodes near coords align={vertical},
every node near coord/.append style={font=\small},
axis line style={draw=popline},
tick label style={font=\small},
major grid style={draw=popline!30},ymajorgrids]
\addplot[fill=popblue!40,draw=popblue] coordinates {(Asie,60) (Afrique,18) (Europe,9) (Amérique N,7.6) (Amérique S,5.4) (Océanie,0.5)};
\end{axis}
\end{tikzpicture}
\end{center}
\legende{Répartition de la population mondiale par continent (environ 8 milliards d'habitants, 2022). L'Asie concentre 60 \% de l'humanité ; l'Afrique, avec 18 \%, assurera pourtant près de la moitié de la croissance future. Source : ONU, Division de la population, 2022.}
\end{popfigure}

Trois enseignements majeurs se dégagent de cette répartition :
\begin{itemize}
\item L'\textbf{Asie} concentre à elle seule 60 \% de l'humanité (Chine et Inde dépassent chacune 1,4 milliard d'habitants).
\item L'\textbf{Afrique} ne représente que 18 \% de la population mondiale, mais elle assurera environ \textbf{50 \% de la croissance démographique future} d'ici 2050 : c'est là que se joue l'avenir démographique de la planète.
\item L'\textbf{Europe} (9 \%) voit sa part décliner régulièrement ; l'Océanie reste marginale (0,5 \%).
\end{itemize}

\begin{popfigure}
\begin{center}
\begin{tikzpicture}[scale=0.95]
% Planisphère très simplifié en blocs (coordonnées avec points décimaux)
\fill[popgreen!25,draw=popdark,rounded corners=2pt] (0.2,2.6) -- (2.2,3.4) -- (2.6,2.2) -- (1.6,1.4) -- (0.4,1.6) -- cycle; % Amérique N
\fill[popgreen!25,draw=popdark,rounded corners=2pt] (2.0,1.2) -- (2.8,0.9) -- (2.4,-0.6) -- (1.8,-0.2) -- cycle; % Amérique S
\fill[popblue!25,draw=popdark,rounded corners=2pt] (4.2,2.8) -- (5.4,3.2) -- (5.6,2.4) -- (4.4,2.2) -- cycle; % Europe
\fill[poporange!30,draw=popdark,rounded corners=2pt] (4.0,2.0) -- (5.8,2.2) -- (6.0,0.4) -- (5.0,-0.6) -- (4.2,0.2) -- cycle; % Afrique
\fill[popblue!40,draw=popdark,rounded corners=2pt] (5.8,3.4) -- (9.4,3.4) -- (9.8,1.8) -- (8.2,1.0) -- (6.4,1.6) -- (5.8,2.4) -- cycle; % Asie
\fill[popgold!35,draw=popdark,rounded corners=2pt] (8.6,0.2) -- (9.8,0.3) -- (9.6,-0.5) -- (8.8,-0.4) -- cycle; % Océanie
\node[font=\scriptsize] at (1.4,2.4) {Am. N};
\node[font=\scriptsize] at (2.3,0.3) {Am. S};
\node[font=\scriptsize] at (4.9,2.7) {Europe};
\node[font=\scriptsize] at (5.0,0.9) {Afrique};
\node[font=\scriptsize] at (7.8,2.6) {Asie};
\node[font=\scriptsize] at (9.2,-0.1) {Océanie};
% Flèches de croissance
\draw[->,very thick,popink] (5.0,0.0) -- (5.0,-1.2) node[below,black,font=\small] {Forte croissance $> 2\,\%$/an (Afrique subsaharienne)};
\draw[->,very thick,popblue] (4.9,3.3) -- (4.9,4.2) node[above,black,font=\small] {Croissance nulle ou négative (Europe, Japon)};
\draw[->,very thick,popgold] (8.6,1.6) -- (9.9,1.2) node[right,black,font=\small] {Ralentissement (Chine)};
% Nord
\draw[->,thick,popdark] (10.6,3.4) -- (10.6,4.2) node[above,font=\small] {N};
\end{tikzpicture}
\end{center}
\legende{Carte schématique des contrastes mondiaux d'accroissement naturel (croquis simplifié, sans précision cartographique). Les flèches indiquent les grandes situations : forte croissance en Afrique subsaharienne, stagnation ou déclin en Europe et au Japon, ralentissement marqué en Chine.}
\end{popfigure}

\courssub{Typologie des grands problèmes démographiques}

\textbf{1) L'explosion démographique et la transition inachevée.} En Afrique subsaharienne et dans une partie de l'Asie du Sud, la mortalité a fortement baissé depuis 1950 (grâce à la médecine, à l'hygiène, aux vaccins) mais la natalité reste élevée : la transition démographique est \emph{inachevée}. Il en résulte des taux d'accroissement naturel supérieurs à 2,5 \% par an, c'est-à-dire des populations qui \textbf{doublent en moins de 30 ans}. Les États peinent à construire assez vite écoles, hôpitaux, routes et emplois.

\begin{propriete}[Loi de la transition démographique]
La baisse de la mortalité précède toujours la baisse de la natalité. Entre les deux, l'écart entre naissances et décès crée un \textbf{pic de croissance transitoire} : c'est la phase d'« explosion démographique ». Tant que la fécondité ne rejoint pas la mortalité à un niveau bas, la population continue de croître fortement.
\end{propriete}

\textbf{2) Le vieillissement démographique et le dépeuplement.} À l'opposé, l'Europe, le Japon, la Corée du Sud et désormais la Chine connaissent une fécondité durablement inférieure au seuil de renouvellement (2,1 enfants par femme) et une espérance de vie élevée. La part des personnes âgées s'accroît, la population active diminue, et certains pays (Japon, Allemagne, Italie) voient leur population totale \emph{décliner} : c'est le dépeuplement relatif, aggravé dans les campagnes abandonnées par les jeunes.

\textbf{3) La sous-population relative.} De vastes espaces restent presque vides parce que les conditions naturelles y sont hostiles : l'Amazonie (forêt équatoriale dense), la Sibérie (froid extrême), le Sahara (aridité), et au Cameroun l'Extrême-Nord septentrional et certaines zones de l'Est. Ces « déserts humains » posent des problèmes de mise en valeur et de contrôle du territoire.

\textbf{4) Les déséquilibres hommes/femmes.} En Chine et en Inde, la préférence culturelle pour les garçons, combinée aux politiques de limitation des naissances et à l'échographie, a permis la \textbf{sélection prénatale} : il naît parfois 115 à 120 garçons pour 100 filles (contre 105 naturellement). Il en résulte des dizaines de millions d'« hommes surnuméraires » qui ne pourront se marier, source de tensions sociales.

\textbf{5) Les migrations forcées et les réfugiés.} Conflits (Syrie, Soudan, Ukraine, bassin du lac Tchad), persécutions et de plus en plus dérèglements climatiques (sécheresses, montée des eaux) jettent sur les routes plus de 100 millions de personnes déplacées ou réfugiées. Ces mouvements subits désorganisent les régions de départ et pèsent sur les régions d'accueil.

\courssub{Surpopulation et surpeuplement : une distinction essentielle}

\begin{definition}[Surpopulation et surpeuplement]
La \textbf{surpopulation} désigne une situation où les \emph{ressources} disponibles (terres, eau, emplois) sont insuffisantes pour assurer correctement la subsistance de la population : c'est un rapport \cle{population/ressources}. Le \textbf{surpeuplement} désigne une \emph{densité} excessive par rapport à l'espace disponible : c'est un rapport \cle{population/espace}.
\end{definition}

Les deux notions ne se recouvrent pas toujours. Le \textbf{Rwanda} connaît un \emph{surpeuplement rural} : plus de 400 habitants au km², des parcelles morcelées devenues trop petites pour nourrir une famille. Le \textbf{Bangladesh} (plus de 1 100 hab/km²) cumule surpeuplement et \emph{surpopulation} : les terres agricoles et les emplois ne suffisent pas à la population, malgré une remarquable baisse de la fécondité. À l'inverse, un pays peut être peu dense mais en surpopulation si ses ressources sont quasi nulles (pays sahéliens), ou très dense sans surpopulation s'il est riche et importe ses ressources (Singapour, Pays-Bas).

\begin{attention}[Erreur fréquente : « l'Afrique est surpeuplée »]
Non. La densité moyenne de l'Afrique (environ 45 hab/km²) est \emph{inférieure} à celle de l'Europe (environ 75 hab/km²) et très inférieure à celle de l'Asie (environ 150 hab/km²). Le vrai problème africain n'est pas le nombre absolu d'habitants, mais la \textbf{rapidité de la croissance} : les infrastructures (écoles, hôpitaux, logements, emplois) ne suivent pas le rythme. Dans une copie, ne jamais écrire « l'Afrique est surpeuplée » sans nuance.
\end{attention}

\courssub{Le vieillissement : indicateurs et conséquences}

\begin{definition}[Taux de dépendance]
Le \textbf{taux de dépendance} est le rapport entre la population inactive (les 0--14 ans et les 65 ans et plus) et la population en âge de travailler (15--64 ans), exprimé en pourcentage :
\[ \text{Taux de dépendance} = \frac{\text{population 0--14 ans} + \text{population 65 ans et plus}}{\text{population 15--64 ans}} \times 100 \]
On distingue le taux de dépendance des \emph{jeunes} (0--14 ans) et celui des \emph{personnes âgées} (65 ans et plus).
\end{definition}

On mesure aussi le vieillissement par l'\textbf{âge médian} (âge qui partage la population en deux moitiés égales : environ 15 ans au Niger, plus de 48 ans au Japon) et par l'\textbf{index de vieillissement} (rapport des 65 ans et plus aux moins de 15 ans). Les conséquences du vieillissement sont lourdes :
\begin{itemize}
\item \textbf{Financement des retraites} : de moins en moins d'actifs cotisent pour de plus en plus de retraités ; les systèmes par répartition sont sous tension.
\item \textbf{Santé} : les dépenses liées aux maladies chroniques et à la dépendance (perte d'autonomie) explosent.
\item \textbf{Pénurie de main-d'œuvre} : certains secteurs (bâtiment, soins, agriculture) ne trouvent plus de travailleurs, ce qui pousse à l'immigration ou à l'automatisation.
\item \textbf{« Silver economy »} : en contrepartie, un nouveau marché s'ouvre autour des besoins des seniors (maisons médicalisées, technologies d'assistance, loisirs adaptés), source d'emplois.
\end{itemize}

\courssub{Le dividende démographique : une fenêtre d'opportunité}

\begin{definition}[Dividende démographique]
Le \textbf{dividende démographique} est le gain de croissance économique potentiel qui résulte de l'augmentation de la part de la population en âge de travailler (15--64 ans) par rapport à la population inactive (jeunes et personnes âgées). Il s'ouvre lorsque la fécondité baisse : les générations d'enfants deviennent moins nombreuses alors que les générations d'actifs sont encore nombreuses.
\end{definition}

L'intuition est simple : si, pour 100 actifs, il n'y a plus que 50 personnes à charge au lieu de 90, la société dégage de l'épargne et des forces de travail pour investir. Mais cette « fenêtre » ne se transforme en prospérité que sous \textbf{conditions} : baisse effective de la fécondité, \emph{éducation} de qualité, \emph{emplois} créés en masse, \emph{investissements} publics et privés. Les « \textbf{Tigres asiatiques} » (Corée du Sud, Taïwan, Singapour, Thaïlande) ont su capter ce dividende entre 1970 et 2000 grâce à la scolarisation massive et à l'industrialisation. À l'inverse, si les jeunes actifs ne trouvent pas d'emplois, la fenêtre devient une « \textbf{bombe démographique} » : chômage de masse, émigration, instabilité — risque réel au Sahel.

\begin{methode}[Analyser une pyramide des âges]
\begin{enumerate}
\item Observer la \textbf{forme générale} : pyramide (base large, sommet pointu = pays jeune), colonne (pays en transition avancée), bulbe ou pyramide inversée (pays vieillissant).
\item Examiner la \textbf{base} : large = natalité forte ; étroite = natalité faible.
\item Examiner le \textbf{sommet} : pointu = espérance de vie faible ; large et évasé = espérance de vie forte, vieillissement.
\item Repérer les \textbf{creux et les bosses} : ils traduisent des événements passés (guerres, épidémies, baby-boom, politiques de limitation des naissances).
\item Comparer les \textbf{sex-ratios} à la base : un excès anormal de garçons peut révéler une sélection sexuée prénatale.
\end{enumerate}
\end{methode}

\courssub{Le cas du Cameroun : une jeunesse à la fois atout et défi}

Le Cameroun illustre parfaitement la transition inachevée. Sa pyramide des âges est une \textbf{pyramide classique à base très large} : environ \textbf{43 \% de la population a moins de 15 ans}. Le taux de dépendance des jeunes atteint environ \textbf{85 \%} : pour 100 actifs, il faut scolariser, soigner et nourrir 85 jeunes. L'\textbf{urbanisation est rapide} : près de \textbf{58 \% des Camerounais vivent en ville}, souvent dans des quartiers précaires de Douala ou Yaoundé, ce qui crée une forte pression sur le logement, l'eau et les transports. La pression sur l'éducation, la santé et l'emploi est considérable : il faudrait créer des centaines de milliers d'emplois par an. Enfin, le pays connaît de fortes \textbf{disparités internes} : le Nord (Extrême-Nord, Nord, Adamaoua) garde une fécondité élevée (souvent plus de 6 enfants par femme) et une pauvreté forte, tandis que le Sud (Centre, Littoral, Ouest) est engagé dans une transition plus avancée. C'est tout l'enjeu de la prospective « Cameroun 2035 » : transformer cette jeunesse en dividende démographique plutôt qu'en facteur de crise.

\begin{exoresolu}[Calculer et interpréter un taux de dépendance]
\textbf{Énoncé.} Dans un pays A, on compte 9 millions de 0--14 ans, 14 millions de 15--64 ans et 1 million de 65 ans et plus. Dans un pays B : 2 millions de 0--14 ans, 13 millions de 15--64 ans et 5 millions de 65 ans et plus. 1) Calculer le taux de dépendance total de chaque pays. 2) Comparer la nature des problèmes démographiques des deux pays.
\tcblower
\textbf{Solution.} 1) Pays A :
\[ \frac{9 + 1}{14} \times 100 \approx 71{,}4\ \% \]
Pays B :
\[ \frac{2 + 5}{13} \times 100 \approx 53{,}8\ \% \]
2) Dans le pays A, la charge vient surtout des \emph{jeunes} (9 millions contre 1 million de personnes âgées) : c'est un pays jeune, en transition inachevée, dont les priorités sont les écoles, la santé maternelle et infantile et la création massive d'emplois. Dans le pays B, la charge vient surtout des \emph{personnes âgées} (5 millions) : c'est un pays vieillissant, confronté au financement des retraites, aux dépenses de santé des seniors et à la pénurie de main-d'œuvre. Même ordre de grandeur du taux global, mais diagnostics et politiques à mener totalement différents.
\end{exoresolu}

\begin{savaistu}[Records démographiques]
Le \textbf{Niger} détient l'ICF le plus élevé du monde (environ 6,7 enfants par femme) et l'âge médian le plus bas (environ 15 ans) : la moitié des Nigériens a moins de 15 ans ! À l'opposé, \textbf{Monaco} affiche la densité la plus forte de la planète (environ 26 000 hab/km²) tout en ayant une population très vieillissante et très riche : preuve qu'une forte densité n'implique pas automatiquement la pauvreté.
\end{savaistu}

\begin{aretenir}[L'essentiel de la section]
\begin{itemize}
\item Les problèmes démographiques mondiaux sont \textbf{contrastés} : explosion démographique (Afrique subsaharienne), vieillissement et dépeuplement (Europe, Japon, Chine), sous-population (Amazonie, Sibérie, Sahara), déséquilibres des sexes (Chine, Inde), migrations forcées.
\item \textbf{Surpopulation} $\neq$ \textbf{surpeuplement} : le premier oppose population et ressources, le second population et espace.
\item Le \textbf{vieillissement} se mesure par le taux de dépendance des personnes âgées, l'âge médian et l'index de vieillissement ; il menace retraites, santé et main-d'œuvre.
\item Le \textbf{dividende démographique} est une fenêtre d'opportunité conditionnée par l'éducation, l'emploi et l'investissement ; sans emplois, il devient « bombe démographique ».
\item Le \textbf{Cameroun} (43 \% de moins de 15 ans, 58 \% d'urbains) doit transformer sa jeunesse en atout d'ici 2035.
\end{itemize}
\end{aretenir}

\coursec{Dossier 1 : Les politiques démographiques – Types, acteurs et bilans}

Dans la section précédente, nous avons dressé le diagnostic des grands problèmes démographiques mondiaux : explosion démographique au sud du Sahara, vieillissement en Europe et en Asie de l'Est, déséquilibres spatiaux entre villes surpeuplées et campagnes désertées. Face à ces déséquilibres, les États ne restent pas passifs : ils conçoivent des \cle{politiques démographiques}, c'est-à-dire des ensembles cohérents de mesures destinées à infléchir volontairement la dynamique des populations. Ce dossier étudie leurs types, leurs acteurs et leurs bilans, à travers trois études de cas : la Chine, la France et le Cameroun.

\courssub{Qu'est-ce qu'une politique démographique ?}

\begin{definition}[Politique démographique]
Une \cle{politique démographique} est l'ensemble des mesures prises par un État (seul ou avec des partenaires) pour modifier volontairement la taille, le rythme de croissance, la structure par âge ou la répartition spatiale de sa population. Elle agit sur trois variables : la \cle{natalité}, la \cle{mortalité} et les \cle{migrations}.
\end{definition}

L'intuition est simple : si la démographie est une « machine » réglée par la natalité, la mortalité et les migrations, alors l'État peut chercher à tourner ces trois boutons. Mais contrairement à une machine, la population est faite d'êtres humains dotés de droits : toute politique démographique se situe donc sur une ligne de crête entre l'intérêt collectif et la liberté individuelle de procréer. C'est ce qui rend ce dossier à la fois géographique et éthique.

On distingue quatre grandes familles de politiques démographiques :

\begin{enumerate}
\item les politiques \cle{anti-natalistes}, qui cherchent à freiner les naissances ;
\item les politiques \cle{pro-natalistes}, qui cherchent à stimuler les naissances ;
\item les politiques \cle{migratoires}, qui règlent les entrées et les sorties de population ;
\item les politiques de \cle{répartition spatiale}, qui déplacent les hommes sur le territoire.
\end{enumerate}

\courssub{Les politiques anti-natalistes : freiner la croissance}

\begin{definition}[Politique anti-nataliste]
Ensemble de mesures incitatives ou coercitives visant à réduire le taux de natalité et la fécondité (exemples : planning familial, âge minimum au mariage, stérilisation, politique de l'enfant unique).
\end{definition}

Ces politiques apparaissent surtout après 1950, dans les pays du Sud confrontés à une croissance démographique rapide que l'économie ne parvient pas à suivre. Leurs instruments vont du plus doux au plus brutal :

\begin{itemize}
\item la diffusion de la \cle{contraception} et le planning familial (Iran, Cameroun) ;
\item l'éducation des filles et le recul de l'âge au mariage, qui raccourcissent mécaniquement la période de fécondité ;
\item les incitations financières (primes, avantages fiscaux pour les petites familles) ;
\item les mesures coercitives : amendes, stérilisations forcées, limitation légale du nombre d'enfants.
\end{itemize}

\begin{exemplebox}[Trois expériences anti-natalistes contrastées]
\begin{itemize}
\item \textbf{Chine (1979--2015)} : la politique de l'enfant unique interdit à la plupart des couples urbains d'avoir plus d'un enfant, sous peine de lourdes amendes. Pékin affirme avoir ainsi « évité » environ 400 millions de naissances.
\item \textbf{Inde (années 1970)} : sous l'état d'urgence (1975--1977), des millions de stérilisations, souvent forcées, sont pratiquées ; le scandale provoque un rejet durable qui freine encore le planning familial indien.
\item \textbf{Iran (années 1990)} : succès remarquable d'un planning familial volontaire, appuyé par les autorités religieuses : l'ICF passe d'environ 6,5 enfants par femme en 1985 à moins de 2 en 2005, sans coercition majeure.
\end{itemize}
\end{exemplebox}

\begin{attention}[Erreur fréquente]
Croire que la politique de l'enfant unique est la seule cause de la baisse de la fécondité en Chine. La baisse a commencé \emph{avant} 1979 : dans les années 1970, la campagne « \emph{wan xi shao} » (mariage tardif, naissances espacées, peu d'enfants), le développement économique et l'éducation des filles avaient déjà fait chuter l'ICF de 5,8 (1970) à 2,7 (1980). La politique coercitive a accéléré une tendance déjà engagée.
\end{attention}

\courssub{Les politiques pro-natalistes : relancer la natalité}

\begin{definition}[Politique pro-nataliste]
Ensemble de mesures incitatives (allocations, avantages fiscaux, services de garde d'enfants, congé parental) visant à encourager les naissances pour contrer le vieillissement ou le dépeuplement.
\end{definition}

Ces politiques concernent les pays ayant achevé leur transition démographique, où l'ICF est durablement inférieur au seuil de renouvellement des générations (2,1 enfants par femme). Leur logique : réduire le « coût » économique et social de l'enfant pour les parents.

\begin{exemplebox}[Quatre modèles pro-natalistes]
\begin{itemize}
\item \textbf{France} : politique familiale universaliste (allocations pour tous), crèches et congé parental (voir étude de cas 3.6).
\item \textbf{Suède} : modèle social et égalitaire — congé parental long (480 jours) partagé entre les deux parents, garde d'enfants quasi gratuite : l'emploi des mères est compatible avec la fécondité.
\item \textbf{Singapour} : campagnes directes (« \emph{Have 3 or more if you can afford it} »), primes de naissance et avantages au logement, ciblées sur les couples diplômés.
\item \textbf{Hongrie} : depuis 2019, les femmes ayant quatre enfants sont exonérées d'impôt sur le revenu \emph{à vie} ; prêts immobiliers effacés à la troisième naissance.
\end{itemize}
\end{exemplebox}

\courssub{Politiques migratoires et politiques de répartition spatiale}

\textbf{Les politiques migratoires} règlent les flux internationaux :

\begin{itemize}
\item \textbf{Sélectives} : le Canada et l'Australie attribuent des visas selon un système de points (âge, diplômes, langue) pour attirer une main-d'œuvre jeune et qualifiée qui compense le vieillissement.
\item \textbf{Restrictives} : l'Europe (politique de Schengen, accords de réadmission) et les États-Unis durcissent l'entrée des migrants non désirés, tout en tolérant une immigration de main-d'œuvre peu qualifiée.
\item \textbf{Accords bilatéraux} : conventions de main-d'œuvre entre pays de départ et d'accueil (ex. Philippines et Golfe ; Maroc et Espagne).
\end{itemize}

\textbf{Les politiques de répartition spatiale} agissent à l'intérieur du territoire :

\begin{itemize}
\item \textbf{Transmigration en Indonésie} : transfert organisé de millions de paysans de Java (surpeuplée) vers Sumatra, Kalimantan ou la Papouasie — bilan mitigé (conflits fonciers, déforestation).
\item \textbf{Nouvelles capitales} : Brasília (Brésil, 1960) pour peupler l'intérieur ; Abuja (Nigeria, 1991) pour décongestionner Lagos ; la nouvelle capitale administrative égyptienne à l'est du Caire pour désengorger une mégapole de plus de 20 millions d'habitants.
\end{itemize}

\begin{popfigure}
\begin{center}
\begin{tikzpicture}[scale=0.9]
% Planisphère très simplifié
\draw[fill=popblueL,draw=popblue] (-5.2,0.6) .. controls (-5.4,2.2) and (-3.2,2.6) .. (-2.6,1.8) .. controls (-2.2,1.2) and (-3.4,0.4) .. (-4.2,-0.2) .. controls (-4.8,-0.6) and (-5.0,-0.2) .. (-5.2,0.6);
\draw[fill=popblueL,draw=popblue] (-3.4,-1.2) .. controls (-3.0,-0.6) and (-2.4,-0.8) .. (-2.4,-1.8) .. controls (-2.4,-2.8) and (-3.0,-3.2) .. (-3.3,-2.6) .. controls (-3.6,-2.0) and (-3.7,-1.7) .. (-3.4,-1.2);
\draw[fill=popgreenL,draw=popgreen] (-0.6,1.4) .. controls (0.2,2.4) and (1.2,2.2) .. (1.4,1.4) .. controls (1.6,0.8) and (0.4,0.6) .. (-0.2,0.8) .. controls (-0.8,1.0) and (-1.0,1.0) .. (-0.6,1.4);
\draw[fill=popgoldL,draw=popgold] (-0.2,0.2) .. controls (0.6,0.8) and (1.4,0.4) .. (1.4,-0.6) .. controls (1.4,-1.8) and (0.8,-2.8) .. (0.4,-2.6) .. controls (-0.2,-2.2) and (-0.6,-0.8) .. (-0.2,0.2);
\draw[fill=poporangeL,draw=poporange] (2.0,1.2) .. controls (2.6,2.4) and (4.6,2.4) .. (5.2,1.4) .. controls (5.6,0.6) and (4.6,0.0) .. (3.8,0.2) .. controls (3.0,0.4) and (1.6,0.4) .. (2.0,1.2);
\draw[fill=poporangeL,draw=poporange] (4.6,-0.6) .. controls (5.0,-0.2) and (5.4,-0.6) .. (5.2,-1.2) .. controls (5.0,-1.7) and (4.4,-1.6) .. (4.3,-1.2) .. controls (4.2,-0.9) and (4.3,-0.9) .. (4.6,-0.6);
% Flèches et étiquettes
\draw[<-,very thick,popgreen] (0.6,1.7) -- (-1.6,3.0) node[left,popdark,align=right] {\textbf{Pro-nataliste}\\ (Europe, France)};
\draw[<-,very thick,poporange] (3.8,1.6) -- (5.8,3.0) node[right,popdark,align=left] {\textbf{Anti-nataliste historique}\\ (Chine, Inde)};
\draw[<-,very thick,popblue] (-4.0,1.6) -- (-5.8,3.0) node[left,popdark,align=right] {\textbf{Migratoire sélective}\\ (Canada, Australie)};
\draw[<-,very thick,popgold] (0.6,-1.4) -- (2.6,-3.2) node[right,popdark,align=left] {\textbf{Transition en cours}\\ (Afrique : planning familial)};
\end{tikzpicture}
\end{center}
\legende{Croquis schématique des grandes politiques démographiques dans le monde (contours très simplifiés). 1 : Europe pro-nataliste ; 2 : Asie anti-nataliste historique ; 3 : Amérique du Nord, immigration sélective ; 4 : Afrique, politiques de planning familial en cours.}
\end{popfigure}

\courssub{Étude de cas approfondie : la Chine, du tout-coercitif au revirement}

\textbf{Contexte.} En 1979, la Chine compte près d'un milliard d'habitants et craint de ne pouvoir ni nourrir ni employer sa population. Deng Xiaoping impose alors la \cle{politique de l'enfant unique} : un enfant par couple urbain (deux tolérés dans certaines campagnes et pour les minorités). Les contrevenants paient des « frais de super-naissance » pouvant atteindre plusieurs années de revenu d'un paysan ; les fonctionnaires récalcitrants risquent leur emploi.

\textbf{Résultats affichés.} Selon Pékin, environ \textbf{400 millions de naissances} ont été évitées entre 1979 et 2015. L'ICF tombe sous le seuil de renouvellement dès les années 1990.

\textbf{Effets secondaires lourds :}
\begin{itemize}
\item un \cle{vieillissement} très rapide : la Chine « vieillira avant d'être riche » ;
\item un déséquilibre des sexes à la naissance : environ \textbf{115 garçons pour 100 filles} (contre 105 naturellement), lié aux avortements sélectifs dans une culture valorisant les garçons — des dizaines de millions d'hommes ne pourront se marier ;
\item une \cle{main-d'œuvre} qui diminue depuis 2012, menaçant le modèle économique.
\end{itemize}

\textbf{Revirement.} Devant ces déséquilibres, Pékin autorise deux enfants en 2016, puis trois en 2021, et multiplie les incitations financières (primes, congés, avantages au logement). Échec relatif : la natalité chute à son plus bas niveau historique en 2023 (ICF proche de 1,0) et la population chinoise \emph{diminue} depuis 2022. Les jeunes urbains, confrontés au coût du logement et de l'éducation, ne veulent plus de familles nombreuses : la coercition a durablement transformé les mentalités.

\courssub{Étude de cas : la France, une politique familiale continue}

La France mène une politique familiale \emph{continue} depuis le \textbf{Code de la famille de 1939}, ce qui en fait un cas unique en Europe.

\textbf{Principes et instruments :}
\begin{itemize}
\item caractère \cle{universaliste} : les allocations familiales concernent toutes les familles, sans condition de revenus (pour l'essentiel) ;
\item soutien à la \cle{conciliation travail-famille} : réseau dense de crèches, école maternelle dès 3 ans, congé parental pouvant aller jusqu'à 3 ans ;
\item avantages fiscaux (quotient familial).
\end{itemize}

\textbf{Résultats et limites.} Avec un ICF d'environ 1,68 (2023), la France reste parmi les plus fécondes de l'Union européenne, mais \emph{sous} le seuil de renouvellement (2,1). Le coût est élevé : environ \textbf{80 milliards d'euros par an}, soit près de 4 \% du PIB.

\begin{exemplebox}[Chiffres de la politique familiale française]
Allocations familiales : environ \SI{140}{\euro}/mois pour le 2\textsuperscript{e} enfant, \SI{315}{\euro}/mois à partir du 3\textsuperscript{e} (montants 2024). Congé parental : jusqu'à 3 ans par enfant. À comparer avec la Chine d'avant 2015 : l'amende pour un enfant « hors quota » pouvait représenter plusieurs années de revenu paysan — deux philosophies opposées de l'action publique.
\end{exemplebox}

\courssub{Étude de cas : le Cameroun, une politique de population à bilan mitigé}

Le Cameroun adopte sa \cle{Déclaration de politique de population} en \textbf{1991}, révisée en \textbf{2007}, dans le cadre de la vision « Cameroun émergent 2035 ».

\textbf{Objectifs :}
\begin{itemize}
\item réduire la fécondité (cible : ICF de 3,5 en 2035) ;
\item réduire la mortalité maternelle et infantile ;
\item intégrer la variable population dans la planification du développement ;
\item mieux gérer les migrations internes et internationales.
\end{itemize}

\textbf{Moyens mobilisés :} promotion du planning familial et de la santé reproductive, scolarisation des filles (levier majeur de baisse de la fécondité), campagnes de sensibilisation, appui des partenaires internationaux (UNFPA, USAID).

\textbf{Bilan mitigé :}
\begin{itemize}
\item l'ICF baisse, mais \emph{lentement} : de 6,4 enfants par femme en 1991 à environ 4,6 en 2023 — loin de la cible ;
\item la mortalité maternelle reste élevée : environ \textbf{406 décès pour 100~000 naissances vivantes} (EDS 2018) ;
\item la pression foncière s'aggrave dans l'Extrême-Nord, où la densité dépasse 100 hab/km\textsuperscript{2} sur des terres fragiles, source de tensions entre agriculteurs et éleveurs.
\end{itemize}

\begin{popfigure}
\begin{center}
\begin{tikzpicture}
\begin{axis}[width=12cm,height=7cm,ylabel={ICF (enfants par femme)},xlabel={Ann\'ees},legend pos=north east,ymin=0,ymax=7,grid=major,grid style={popline!40}]
\addplot[thick,red,mark=*] coordinates {(1970,5.8) (1980,2.7) (1990,2.3) (2000,1.6) (2010,1.6) (2020,1.3) (2023,1.0)};
\addlegendentry{Chine}
\addplot[thick,blue,mark=square*] coordinates {(1970,2.5) (1980,1.9) (1990,1.8) (2000,1.9) (2010,2.0) (2020,1.8) (2023,1.68)};
\addlegendentry{France}
\addplot[thick,green!60!black,mark=triangle*] coordinates {(1970,6.4) (1980,6.2) (1990,5.8) (2000,5.2) (2010,5.0) (2020,4.6) (2023,4.6)};
\addlegendentry{Cameroun}
\addplot[dashed,popmuted] coordinates {(1970,2.1) (2023,2.1)};
\addlegendentry{Seuil de renouvellement (2,1)}
\end{axis}
\end{tikzpicture}
\end{center}
\legende{Évolution de l'indice conjoncturel de fécondité (ICF) en Chine, en France et au Cameroun, 1970--2023. La Chine est passée sous le seuil de renouvellement dès les années 1990 ; le Cameroun reste au-dessus de 4,5 ; la France se stabilise autour de 1,8.}
\end{popfigure}

\courssub{Les acteurs des politiques démographiques}

Une politique démographique n'est jamais l'œuvre du seul État :

\begin{itemize}
\item \textbf{L'État} : législation (âge au mariage, avortement, allocations), financement, services publics de santé et d'éducation ;
\item \textbf{Les ONG et agences internationales} : l'\cle{UNFPA} (Fonds des Nations unies pour la population), l'USAID, la Fondation Gates financent le planning familial, forment les personnels de santé et fournissent les contraceptifs ;
\item \textbf{Les leaders religieux et traditionnels} : leur influence sur les normes sociales est décisive — l'appui du clergé iranien a fait réussir le planning familial, tandis qu'au Cameroun certaines autorités coutumières freinent la contraception ;
\item \textbf{Le secteur privé} : production et distribution des contraceptifs, marketing social ;
\item \textbf{La société civile} : associations de femmes, jeunesse, médias, qui font évoluer les mentalités.
\end{itemize}

\courssub{Évaluer l'efficacité d'une politique démographique}

\begin{propriete}[Efficacité comparée des politiques démographiques]
Les politiques anti-natalistes \emph{coercitives} donnent des résultats rapides mais au prix de violations des droits et d'effets secondaires durables (Chine : sex-ratio déséquilibré, vieillissement brutal). Les politiques pro-natalistes sont coûteuses, lentes, et ramènent rarement l'ICF au seuil de renouvellement (2,1) sans apport migratoire. Les politiques \emph{volontaires} fondées sur l'éducation des filles et la santé reproductive sont les plus durables.
\end{propriete}

\begin{methode}[Analyser une politique démographique : grille d'évaluation]
\begin{enumerate}
\item \textbf{Contexte initial} : quels indicateurs (ICF, mortalité, croissance) justifient l'action ?
\item \textbf{Objectifs affichés} : freiner, stimuler, redistribuer ?
\item \textbf{Instruments} : légaux, financiers, sanitaires, éducatifs ?
\item \textbf{Acteurs et acceptabilité} : qui décide, qui finance, qui résiste ?
\item \textbf{Résultats chiffrés à 10--20 ans} : évolution de l'ICF, de la mortalité, du taux de contraception, de l'espérance de vie ;
\item \textbf{Effets secondaires} : sex-ratio, vieillissement, respect des droits humains ;
\item \textbf{Leçons transférables} : que peut en tirer un autre pays, comme le Cameroun ?
\end{enumerate}
\end{methode}

\courssub{Le débat éthique : malthusianisme contre droit à la procréation}

Depuis Malthus (1798), qui craignait que la population ne croisse plus vite que les ressources, le \cle{malthusianisme} justifie le contrôle des naissances au nom de l'intérêt collectif. À l'opposé, la \textbf{Conférence internationale sur la population et le développement du Caire (1994)} consacre la \cle{santé reproductive} comme un \emph{droit humain} : chaque couple doit décider \emph{librement} du nombre et de l'espacement de ses enfants. Toute coercition (stérilisations forcées, amendes) est désormais condamnée par le droit international. Se pose aussi la question de la \cle{conditionnalité de l'aide} : certains pays du Nord ont subordonné leur aide au développement à l'adoption de programmes de limitation des naissances, pratique jugée néo-coloniale par de nombreux États africains.

\begin{savaistu}[Le consensus du Caire (1994)]
La CIPD du Caire marque un tournant : on passe du « contrôle de la population » à la « santé reproductive et droits reproductifs ». Le Programme d'action adopté par 179 pays affirme que le développement durable ne peut se faire sans l'égalité femmes-hommes, l'éducation des filles et l'accès universel aux services de santé reproductive. Ce cadre guide encore aujourd'hui les politiques de l'UNFPA et les Objectifs de développement durable (ODD 3 et 5).
\end{savaistu}

\begin{exoresolu}[Comparer deux politiques démographiques]
Énoncé : à l'aide de la grille d'évaluation, comparez la politique chinoise de l'enfant unique (1979--2015) et la politique familiale française. Concluez sur leur efficacité respective.
\tcblower
\textbf{Solution détaillée.}
\begin{itemize}
\item \emph{Contexte} : Chine, croissance rapide menaçant le développement (ICF 2,7 en 1980, déjà en baisse) ; France, fécondité insuffisante et vieillissement naissant dès les années 1930.
\item \emph{Objectifs et instruments} : la Chine choisit la coercition (interdiction légale, amendes de plusieurs années de revenu) ; la France choisit l'incitation universaliste (allocations d'environ \SI{140}{\euro}/mois dès le 2\textsuperscript{e} enfant, crèches, congé parental de 3 ans, coût total voisin de 4 \% du PIB).
\item \emph{Résultats} : la Chine fait chuter son ICF à 1,6 en 2000 puis proche de 1,0 en 2023 — objectif dépassé, mais avec 400 millions de naissances « évitées » au prix de violations des droits ; la France maintient un ICF de 1,68, le plus élevé d'Europe, mais sous le seuil de 2,1.
\item \emph{Effets secondaires} : Chine — 115 garçons pour 100 filles à la naissance, vieillissement brutal, main-d'œuvre en baisse depuis 2012 ; France — coût budgétaire élevé, mais aucun effet pervers majeur.
\item \emph{Conclusion} : la politique chinoise a été « efficace » à court terme mais non durable et contraire aux droits humains, au point que Pékin échoue aujourd'hui à remonter sa natalité ; la politique française, coûteuse mais respectueuse des libertés, stabilise une fécondité modérée. La leçon transférable au Cameroun : privilégier les leviers volontaires (scolarisation des filles, santé reproductive) plutôt que la contrainte.
\end{itemize}
\end{exoresolu}

\begin{aretenir}[L'essentiel du dossier]
\begin{itemize}
\item Quatre types de politiques démographiques : anti-natalistes, pro-natalistes, migratoires, de répartition spatiale.
\item La coercition (Chine) agit vite mais viole les droits et crée des déséquilibres durables (sex-ratio, vieillissement) ; l'incitation (France) est coûteuse et plafonne sous 2,1.
\item Le Cameroun (politique de population de 1991, révisée en 2007) progresse lentement : ICF de 6,4 à 4,6 entre 1991 et 2023, mortalité maternelle encore à 406 pour 100~000.
\item Depuis le Caire (1994), la santé reproductive est un droit humain : le consentement libre prime sur les objectifs chiffrés.
\end{itemize}
\end{aretenir}

\begin{exercice}[Pour s'entraîner]
\begin{enumerate}
\item Définis les termes : politique pro-nataliste, seuil de renouvellement, transmigration.
\item Explique pourquoi la baisse de la fécondité chinoise ne peut pas être attribuée uniquement à la politique de l'enfant unique.
\item Rédige un paragraphe argumenté (10 lignes) : « Le Cameroun doit-il s'inspirer du modèle iranien de planning familial ? » en mobilisant au moins trois chiffres du cours.
\end{enumerate}
\end{exercice}

\begin{corrige}[Exercice 1]
\begin{enumerate}
\item Politique pro-nataliste : ensemble de mesures incitatives (allocations, avantages fiscaux, garde d'enfants, congé parental) visant à encourager les naissances. Seuil de renouvellement : ICF de 2,1 enfants par femme assurant le remplacement des générations. Transmigration : déplacement organisé de populations (Indonésie : de Java vers les îles extérieures).
\item L'ICF chinois était déjà passé de 5,8 (1970) à 2,7 (1980) \emph{avant} 1979, grâce à la campagne « wan xi shao », au développement et à l'éducation des filles ; la politique de l'enfant unique a accéléré, non créé, la tendance.
\item Attendu : rappel du succès iranien (ICF de 6,5 à moins de 2 entre 1985 et 2005) fondé sur le volontariat et l'appui religieux ; comparaison avec le Cameroun (ICF 4,6 en 2023, mortalité maternelle 406/100~000, cible 3,5 en 2035) ; argument sur la nécessité d'associer leaders religieux et traditionnels et de scolariser les filles plutôt que de contraindre.
\end{enumerate}
\end{corrige}

\coursec{Focus Cameroun : Dynamiques internes, défis et prospective 2035}

Après avoir analysé les politiques démographiques à l'échelle mondiale (section 3), nous recentrons notre regard sur le terrain camerounais. Le Cameroun, souvent qualifié d'\og Afrique en miniature \fg{} par sa diversité écologique et culturelle, concentre en son sein l'essentiel des dynamiques et des tensions démographiques observées sur le continent. Comprendre son cas particulier permet d'illustrer concrètement les concepts de transition démographique, de dividende et de pression sur les ressources.

\courssub{Une croissance rapide : le défi du doublement démographique}

Le Cameroun traverse une \cle{transition démographique inachevée} : la mortalité a chuté grâce aux progrès sanitaires (vaccination, lutte contre le paludisme, accès aux soins), mais la fécondité reste résistante.

\begin{definition}[Transition démographique inachevée]
Situation où la mortalité a fortement baissé mais où la natalité reste élevée, entraînant une croissance démographique rapide qui dépasse les capacités d'investissement social (écoles, santé, emplois, logements, infrastructures).
\end{definition}

\begin{popfigure}
\begin{tikzpicture}
\begin{axis}[
    width=\linewidth,
    height=7cm,
    ylabel={Population (millions)},
    xlabel={Année},
    legend pos=north west,
    grid=major,
    tick label style={font=\small},
    label style={font=\small},
    legend style={font=\small},
    xmin=1985, xmax=2055,
    ymin=0, ymax=55,
    xtick={1987,1997,2007,2017,2023,2030,2050},
    ytick={0,10,20,30,40,50},
]
\addplot[thick,popblue,mark=*] coordinates {(1987,10) (1997,14) (2007,19) (2017,24) (2023,28.6) (2030,35) (2050,50)};
\addlegendentry{Tendance ONU moyenne (variante médiane)}
\addplot[thick,popred,dashed,mark=square*] coordinates {(2023,28.6) (2030,33) (2050,42)};
\addlegendentry{Variante basse (baisse fécondité rapide)}
\end{axis}
\end{tikzpicture}
\legende{Évolution de la population camerounaise : observations (1987--2023) et projections ONU 2024 (variante médiane vs basse). Le doublement en 25--30 ans est le scénario central si l'indice conjoncturel de fécondité (ICF) demeure autour de 4,0--4,2. Source : ONU, Perspectives de la population mondiale 2024 ; INS, RGPH 2005 et estimations.}
\end{popfigure}

\begin{exemplebox}[Lecture du graphique]
En 2023, la population est estimée à \textbf{28,6 millions} d'habitants. La variante médiane (bleu) projette \textbf{35 millions en 2030} et \textbf{50 millions en 2050}. La variante basse (rouge), qui suppose une chute rapide de l'ICF vers 2,5 dès 2030, aboutirait à 42 millions en 2050. L'écart de 8 millions entre les deux scénarios en 2050 représente le \og coût de l'inaction \fg{} en matière de planification familiale et d'éducation des filles.
\end{exemplebox}

\textbf{Facteurs explicatifs.}\par
\begin{itemize}
    \item \textbf{Mortalité en baisse :} Espérance de vie à la naissance $\approx$ 60 ans (2023) contre 50 ans en 1990. Mortalité infantile : 48\,\textperthousand\ (2023) vs 85\,\textperthousand\ (1990).
    \item \textbf{Fécondité résistante :} ICF $\approx$ 4,2 enfants/femme (2023), quasi stable depuis 20 ans. Préférence culturelle pour les grandes familles, faible prévalence contraceptive moderne (\textbf{$\sim$ 20\%} seulement, EDS-MICS 2018), mariage précoce (âge médian 1\textsuperscript{re} union : 18,5 ans femmes).
    \item \textbf{Migrations :} Solde migratoire légèrement positif (retour de réfugiés, immigration nigériane/centrafricaine), mais faible impact sur la croissance globale comparé au solde naturel ($>$ 2,5\%/an).
\end{itemize}

\begin{attention}[Piège fréquent]
Ne pas confondre \emph{taux de croissance} et \emph{effectifs annuels}. Même si le taux baisse (ex: 2,6\% $\to$ 2,3\%), les effectifs annuels d'accroissement augmentent car la population de base est plus large (effet d'inertie démographique). En 2023, l'accroissement naturel est d'environ \textbf{750 000 habitants/an} (équivalent à une ville comme Bamenda chaque année).
\end{attention}

\courssub{Structure par âge : \og Bombe à retardement \fg{} ou \og Dividende démographique \fg{} ?}

La pyramide des âges camerounaise est de type \og expansif \fg{} (base large, sommet étroit), caractéristique des pays en phase de transition précoce.

\begin{definition}[Dividende démographique (cas Cameroun)]
Fenêtre d'opportunité (environ \textbf{2020--2050}) où la part des 15--64 ans (âge actif) augmente plus vite que la population totale, \textbf{à condition d'investir massivement dans le capital humain (éducation, santé) et de créer des emplois décents}. Sans ces investissements, le surplus de main-d'œuvre devient une charge (chômage, instabilité) : c'est la \og bombe à retardement \fg{}.
\end{definition}

\begin{aretenir}[Chiffres clés structure 2023 (INS/ONU)]
\begin{itemize}
    \item \textbf{43\%} de la population a \textbf{moins de 15 ans} (charge de jeunesse très forte).
    \item \textbf{54\%} ont 15--64 ans (potentiel actif).
    \item \textbf{3\%} ont 65 ans et plus.
    \item \textbf{Ratio de dépendance total :} $\approx$ 85\% (85 inactifs pour 100 actifs potentiels).
    \item \textbf{Arrivées annuelles sur le marché du travail :} $\sim$ \textbf{300 000 jeunes/an}.
    \item \textbf{Emplois formels créés/an :} $<$ 30 000 (secteur public + formel privé). \textbf{Déficit : 270 000/an.}
\end{itemize}
\end{aretenir}

\begin{exoresolu}[Analyse du défi emploi-éducation]
\textbf{Énoncé :} En 2023, le taux brut de scolarisation au primaire dépasse 100\% (redoublements, entrées tardives), mais le taux d'achèvement du primaire est faible ($\sim$ 60\%) et les résultats PASEC placent le Cameroun en queue de peloton en lecture/maths. En parallèle, 300 000 jeunes arrivent sur le marché du travail pour $<$ 30 000 emplois formels.

\textbf{Analyse :}
\begin{enumerate}
    \item \textbf{Quantité vs Qualité :} L'accès à l'école (quantité) a progressé, mais la qualité (apprentissages réels) stagne. Un diplôme sans compétences ne crée pas d'employabilité.
    \item \textbf{Adéquation formation-emploi :} L'enseignement technique/professionnel (le vôtre !) est le levier principal pour réduire le décalage. Les filières industrielles, agricoles, tertiaires (BTS, HND, Probatoire technique) doivent répondre aux besoins réels des bassins d'emploi (agro-industrie, BTP, numérique, énergies renouvelables).
    \item \textbf{Économie informelle :} 90\% des emplois sont informels. La politique publique (PAMPEJ, PIAASI, Guichet Unique) doit faciliter la formalisation et l'accès au financement pour les jeunes entrepreneurs.
\end{enumerate}
\textbf{Conclusion :} Capturer le dividende impose une \textbf{révolution de la qualité éducative} et une \textbf{politique industrielle créatrice d'emplois}, pas seulement une expansion quantitative du système scolaire.
\tcblower
\textbf{Solution détaillée :} Voir analyse ci-dessus (structure en 3 points + conclusion). Les chiffres clés (300k entrants vs <30k emplois formels) justifient l'urgence d'une réforme structurelle de l'appareil de formation et de l'économie.
\end{exoresolu}

\courssub{Répartition spatiale inégale et urbanisation accélérée}

Le territoire camerounais présente une \cle{dichotomie Nord/Sud} et une \cle{métropolisation} côtière extrêmes.

\begin{popfigure}
\begin{tikzpicture}
% Fond de carte simplifié Cameroun (approximatif)
% Contour très grossier
\draw[thick, popdark] (0.5,0.5) -- (1,1) -- (2,0.8) -- (3,1.5) -- (4,1) -- (5,2) -- (6,1.5) -- (7,2.5) -- (8,2) -- (9,3) -- (9.5,4) -- (9,5) -- (8,6) -- (7,5.5) -- (6,6.5) -- (5,6) -- (4,7) -- (3,6.5) -- (2,7) -- (1,6) -- (0.5,5) -- (0,4) -- (0.5,3) -- (0,2) -- cycle;

% Zones de densité (choroplèthe simplifié)
% Nord/Est/Adamaoua : Faible densité (<10-20)
\fill[popblue!15] (0.5,0.5) -- (1,1) -- (2,0.8) -- (3,1.5) -- (4,1) -- (5,2) -- (4.5,3) -- (3,3.5) -- (2,3) -- (1,2) -- cycle; % Nord
\fill[popblue!15] (4,4) -- (5,4.5) -- (6,4) -- (5.5,5) -- (4.5,5.5) -- (3.5,5) -- cycle; % Est
\fill[popblue!15] (5,2.5) -- (6,3) -- (5.5,4) -- (4.5,3.5) -- cycle; % Adamaoua

% Ouest/Littoral/NO/SO : Forte densité (>100-200)
\fill[popred!30] (3,3.5) -- (4,4) -- (5,4.5) -- (5.5,4) -- (5,3) -- (4.5,3.5) -- cycle; % Ouest (Bamiléké)
\fill[popred!30] (5,2) -- (6,1.5) -- (7,2.5) -- (6.5,3) -- (5.5,2.5) -- cycle; % Littoral (Douala)
\fill[popred!30] (2.5,4.5) -- (3.5,5) -- (3,6) -- (2,5.5) -- cycle; % Nord-Ouest (Bamenda)
\fill[popred!30] (3.5,5.5) -- (4.5,6) -- (4,6.5) -- (3,6) -- cycle; % Sud-Ouest

% Villes principales (symboles proportionnels)
\node[circle,fill=popdark,minimum size=6pt,inner sep=0] (Douala) at (5.8,1.8) {};
\node[left,font=\tiny] at (Douala.west) {Douala $\sim$ 4 M};
\node[circle,fill=popdark,minimum size=6pt,inner sep=0] (Yaounde) at (4.2,4.2) {};
\node[left,font=\tiny] at (Yaounde.west) {Yaoundé $\sim$ 4 M};
\node[circle,fill=popdark,minimum size=4pt,inner sep=0] (Bamenda) at (2.8,5.2) {};
\node[left,font=\tiny] at (Bamenda.west) {Bamenda};
\node[circle,fill=popdark,minimum size=4pt,inner sep=0] (Bafoussam) at (3.8,3.8) {};
\node[above,font=\tiny] at (Bafoussam.north) {Bafoussam};
\node[circle,fill=popdark,minimum size=4pt,inner sep=0] (Garoua) at (2.5,1.5) {};
\node[below,font=\tiny] at (Garoua.south) {Garoua};
\node[circle,fill=popdark,minimum size=4pt,inner sep=0] (Maroua) at (1.5,0.8) {};
\node[below,font=\tiny] at (Maroua.south) {Maroua};
\node[circle,fill=popdark,minimum size=4pt,inner sep=0] (Ngaoundere) at (5.2,3.2) {};
\node[right,font=\tiny] at (Ngaoundere.east) {Ngaoundéré};

% Flèches migrations internes
\draw[->, thick, popgreen] (1.5,1.2) -- (3.5,2.5) node[midway,above,font=\tiny,sloped] {Nord $\to$ Sud (plantations)};
\draw[->, thick, popgreen] (4.5,5) -- (5.5,3.5) node[midway,above,font=\tiny,sloped] {Ouest $\to$ Littoral/SO};
\draw[->, thick, poporange] (3,5.5) -- (4.2,4.3) node[midway,above,font=\tiny,sloped] {Crise NO/SO $\to$ Centre/Littoral};

% Légende manuelle
\node[draw,fill=white,font=\tiny,anchor=south east, align=center] at (9.5,0.5){
    \begin{tabular}{ll}
    \cellcolor{popred!30} & $>$ 200 hab/km² (Ouest, Littoral, NO, SO) \\
    \cellcolor{popblue!15} & $<$ 20 hab/km² (Est, Adamaoua, Nord, EN) \\
    $\bullet$ & Villes $>$ 300 000 hab. \\
    \textcolor{popgreen}{$\rightarrow$} & Migrations économiques \\
    \textcolor{poporange}{$\rightarrow$} & Migrations forcées (crise)
    \end{tabular}
};

\end{tikzpicture}
\legende{Croquis schématique : Densité de population, principales villes et flux migratoires internes au Cameroun. Les densités $>$ 200 hab/km² (Ouest, Littoral, NO, SO) contrastent avec les vastes espaces $<$ 20 hab/km² (Est, Adamaoua, Nord, Extrême-Nord). Douala et Yaoundé polarisent l'urbanisation (métropolisation). Flèches vertes = migrations économiques traditionnelles ; flèches orange = déplacements liés à la crise anglophone (2016--).}
\end{popfigure}

\begin{propriete}[Relation population-développement (Conférence du Caire 1994)]
Le développement est le meilleur contraceptif (amélioration niveau de vie $\to$ baisse fécondité), \textbf{mais} la maîtrise de la démographie est une condition \emph{nécessaire} au développement durable (investissement par tête, qualité du capital humain). Au Cameroun, l'investissement par habitant (santé, éducation, infrastructure) stagne car le dénominateur (population) croît plus vite que le numérateur (budget public).
\end{propriete}

\textbf{Urbanisation : 58\% (2023), taux 3,6\%/an.}\par
\begin{itemize}
    \item \textbf{Métropolisation :} Douala (capitale économique, port, $\sim$ 4 M) et Yaoundé (capitale politique, $\sim$ 4 M) concentrent 50\% de la population urbaine et 80\% du PIB formel.
    \item \textbf{Problèmes urbains :} Habitat précaire (\og bidonvilles \fg{} : Nkolbisson, Mimboman, New-Bell, etc. $\to$ 60\% des citadins), congestion routière, gestion déchets/eau défaillante, insécurité, inondations récurrentes (Douala).
    \item \textbf{Exode rural :} Vieillissement des campagnes, perte de main-d'œuvre agricole, pression foncière dans les zones d'origine (Ouest) et d'accueil (périurbain).
\end{itemize}

\begin{attention}[Hétérogénéité intra-nationale : erreur fréquente]
Le Cameroun n'est \textbf{pas un bloc uniforme}.
\begin{itemize}
    \item \textbf{Nord (Extrême-Nord, Nord, Adamaoua) :} Indicateurs proches du Sahel. ICF $>$ 6, scolarisation faible (filles $<$ 30\% au secondaire), mortalité maternelle/infantile très élevée, pauvreté monétaire $>$ 50\%.
    \item \textbf{Sud (Centre, Sud, Est, Littoral, SO, NO, Ouest) :} Transition plus avancée. ICF 3,5--4,5, scolarisation meilleure, urbanisation forte.
\end{itemize}
Toute politique nationale \textbf{doit} être territorialisée (décentralisation effective, budgets régionaux différenciés).
\end{attention}

\courssub{Migrations : carrefour complexe}

Le Cameroun est à la fois pays d'émigration, d'immigration, de transit et de refuge.

\begin{center}
\begin{tabularx}{\linewidth}{|l|X|X|}\hline
\textbf{Type} & \textbf{Flux principaux} & \textbf{Enjeux / Chiffres (2023)} \\ \hline
\textbf{Internes} & Nord $\to$ Sud (plantations CDC, HEVECAM, agro-industrie) ; Ouest $\to$ Littoral/SO (commerce, plantations) ; Est $\to$ Ouest/Centre (mines, forêt) ; \textbf{NO/SO $\to$ Centre/Littoral/Ouest} (crise anglophone, 700k+ déplacés internes). & Redistribution main-d'œuvre ; pression foncière/urbaine ; traumatismes psychosociaux (PDI). \\ \hline
\textbf{Immigration} & Nigéria (commerce, main-d'œuvre, réfugiés Borno), RCA (réfugiés, commerce), Tchad, Diaspora France/USA (retour investissement, compétences). & $\sim$ 470 000 réfugiés/demandeurs d'asile (HCR 2023, mostly RCA/Nigéria). Intégration locale difficile. \\ \hline
\textbf{Émigration} & Élites, étudiants, main-d'œuvre qualifiée (médecins, ingénieurs, IT) $\to$ France, USA, Canada, Allemagne, Pays du Golfe. & \textbf{Fuite des cerveaux (Brain drain)} : 1 médecin sur 3 formé au Cameroun exerce à l'étranger. Transferts financiers (remittances) $\sim$ 300 millions \$US/an (soutien ménages). \\ \hline
\end{tabularx}
\end{center}

\courssub{Vision 2035, PND30 et le pari du dividende démographique}

La \textbf{Vision 2035} (\og Cameroun pays émergent, démocratique et uni dans sa diversité \fg{}) et le \textbf{PND30 (Plan National de Développement 2020--2030)} identifient la démographie comme \textbf{levier stratégique} ou \textbf{frein majeur}.

\begin{methode}[Capturer le dividende démographique : les 4 piliers du PND30]
\begin{enumerate}
    \item \textbf{Capital humain -- Éducation/Formation :} Généralisation enseignement technique/professionnel (visée 50\% flux secondaire), qualité (PASEC), alphabétisation adultes, éducation non formelle. \textit{Objectif : employabilité.}
    \item \textbf{Capital humain -- Santé reproductive :} Planning familial (objectif prévalence contraceptive moderne 35\% en 2030), santé maternelle/infantile, lutte VIH/Sida/paludisme, nutrition. \textit{Objectif : baisse ICF vers 3,0--3,5.}
    \item \textbf{Emploi jeunes \& Entrepreneuriat :} PAMPEJ, PIAASI, Fonds Jeunes, Guichet Unique, zones économiques spéciales, appui PME/PMI, économie numérique. \textit{Objectif : 300 000 emplois décents/an.}
    \item \textbf{Gouvernance \& Décentralisation :} Régionalisation effective (transferts compétences/ressources), gestion transparente finances publiques, lutte corruption, statistiques fiables (RGPH prévu). \textit{Objectif : efficacité dépense publique, équité territoriale.}
\end{enumerate}
\end{methode}

\begin{attention}[Défi financement]
Recettes fiscales $\approx$ 13\% PIB (faible), dette publique $\approx$ 45\% PIB (risque surendettement modéré). L'espace budgétaire pour les investissements massifs (écoles, hôpitaux, routes, usines) est contraint. \textbf{Efficience de la dépense} et \textbf{mobilisation ressources internes} (impôt foncier, fiscalité numérique, secteur extractif) sont critiques.
\end{attention}

\courssub{Risques climatiques et migrations environnementales}

Le changement climatique agit comme \textbf{multiplicateur de vulnérabilité} démographique.

\begin{itemize}
    \item \textbf{Nord (Soudano-sahélien) :} Sécheresses récurrentes, désertification, raréfaction eau/pâturages $\to$ insécurité alimentaire chronique, conflits éleveurs/agriculteurs (Logone-et-Chari, Mayo-Danay), migrations de survie vers villes (Maroua, Garoua) ou Sud.
    \item \textbf{Côte (Littoral, Sud-Ouest) :} Montée des eaux, érosion côtière (Douala-Bonabéri, Limbé, Kribi), inondations urbaines, salinisation nappes $\to$ menace sur 40\% du PIB (zone portuaire/industrielle), déplacements populations riveraines.
    \item \textbf{Forêt (Est, Sud, Sud-Ouest) :} Déforestation (agriculture itinérante, exploitation forestière illégale, mines), perte biodiversité, perturbation régimes pluviométriques locaux.
\end{itemize}
\emph{Migration climatique interne :} Déjà une réalité (ex: déplacements Lac Tchad, vallée du Logone). Projections : plusieurs millions de \og migrants climatiques internes \fg{} d'ici 2050 si non adaptation.

\courssub{Scénario prospectif 2035 : Tendanciel vs Volontariste}

\begin{center}
\begin{tabularx}{\linewidth}{|l|X|X|}\hline
\textbf{Dimension} & \textbf{Scénario Tendanciel (Inertie)} & \textbf{Scénario Volontariste (Vision 2035 atteinte)} \\ \hline
\textbf{ICF (2035)} & $\sim$ 3,8 (baisse lente) & $\sim$ 3,0 (accès PF généralisé, scolarisation filles) \\ \hline
\textbf{Population (2035)} & $\sim$ 38 millions & $\sim$ 35 millions \\ \hline
\textbf{Structure âge} & 40\% $<$ 15 ans (dépendance haute) & 35\% $<$ 15 ans (fenêtre dividende ouverte) \\ \hline
\textbf{Capital humain} & Qualité éducation stagnante, santé sous-financée & Éducation technique de qualité, couverture santé universelle (CSU) \\ \hline
\textbf{Emploi} & $<$ 100k formels/an, informel précaire dominant & 250k--300k décents/an, transformation structurelle économie \\ \hline
\textbf{Urbanisation} & Métropolisation chaotique, bidonvilles étendus & Villes intermédiaires développées, habitat planifié, mobilité durable \\ \hline
\textbf{Gouvernance} & Centralisation, corruption, données faibles & Décentralisation effective, reddition comptes, data pour décision \\ \hline
\textbf{Résultat 2035} & \textbf{Piège à faible revenu}, instabilité sociale, pression ressources & \textbf{Pays émergent}, dividende capté, cohésion territoriale \\ \hline
\end{tabularx}
\end{center}

\begin{savaistu}[Le saviez-vous ?]
Le Cameroun est le \textbf{2e pays le plus peuplé d'Afrique centrale} après la RDC (100M+). La \textbf{région de l'Ouest} (pays Bamiléké, sols volcaniques fertiles) voit certains de ses départements (Mifi, Menoua, Haut-Nkam) atteindre des densités \textbf{$>$ 300 hab/km²} (record national), avec une forte tradition d'\textbf{émigration interne} vers les plantations du Sud-Ouest/Littoral (café, cacao, banane, palme, hévéa) et les villes (Douala, Yaoundé, Bafoussam). Cette mobilité historique a façonné l'unité nationale mais crée aujourd'hui des tensions foncières dans les zones d'accueil.
\end{savaistu}

\begin{exercice}[Application : Croquis de synthèse]
À partir des données de cette section, réalisez un \textbf{croquis légendé} de la population du Cameroun sur un fond de carte A4 (régions, fleuves principaux, villes $>$ 100 000 hab).
\begin{enumerate}
    \item Choroplèthe des densités (5 classes : $<$10 ; 10--50 ; 50--100 ; 100--200 ; $>$200 hab/km²).
    \item Symboles proportionnels pour les 7 métropoles régionales (Douala, Yaoundé, Bamenda, Bafoussam, Garoua, Maroua, Ngaoundéré).
    \item Flèches des 3 grands flux migratoires internes (Nord$\to$Sud plantations ; Ouest$\to$Littoral/SO ; Crise NO/SO$\to$Centre/Littoral).
    \item Légende organisée, titre précis, échelle graphique, flèche Nord.
\end{enumerate}
\end{exercice}

\begin{corrige}[Exercice : Croquis de synthèse]
\textbf{Méthode (rappel) :} Voir boîte \og Méthode : Construire un croquis de la population du Cameroun \fg{} ci-dessous.
\textbf{Résultat attendu :}
\begin{itemize}
    \item \textbf{Titre :} \og Cameroun : Répartition de la population, grands pôles urbains et flux migratoires internes (2023) \fg.
    \item \textbf{Choroplèthe :} Ouest/Littoral/NO/SO en rouge foncé ($>$200) ; Centre/Sud/Adamaoua en orange (50--100) ; Nord/Est/Extrême-Nord en jaune/beige ($<$20, sauf poches urbaines Garoua/Maroua/Ngaoundéré en point rouge).
    \item \textbf{Symboles :} Cercles proportionnels (Douala/Yaoundé max), carrés ou triangles pour autres.
    \item \textbf{Flèches :} Épaisseur proportionnelle à l'importance du flux. Flèches vertes (économiques), orange (forcées).
    \item \textbf{Légende :} Hiérarchisée (fond, symboles, flèches).
\end{itemize}
\end{corrige}

\begin{methode}[Construire un croquis de la population du Cameroun]
\begin{enumerate}
    \item \textbf{Fond de carte :} Tracer les 10 régions, fleuves principaux (Sanaga, Bénoué, Logone, Nyong), villes $>$ 100 000 hab. (points de repère).
    \item \textbf{Choroplèthe densités :} 5 classes (ex: $<$10 ; 10--50 ; 50--100 ; 100--200 ; $>$200 hab/km²). Colorier chaque région/unité écologique selon sa densité moyenne (données RGPH/INS). Hachurer les zones de peuplement dense très localisé (ex: contreforts Mandara, vallée Mbéré).
    \item \textbf{Symboles proportionnels villes :} Douala/Yaoundé (très grands cercles), Bamenda/Bafoussam/Garoua/Maroua/Ngaoundéré (moyens), autres chefs-lieux (petits). Surface $\propto$ population.
    \item \textbf{Flèches migrations internes :} 3 flèches principales : (1) Nord $\to$ Sud-Ouest/Littoral (plantations) ; (2) Ouest $\to$ Littoral/Sud-Ouest/Centre (commerce, terre) ; (3) Nord-Ouest/Sud-Ouest $\to$ Centre/Littoral/Ouest (crise 2016--). Épaisseur = importance relative.
    \item \textbf{Légende + Titre :} Titre informatif (quoi, où, quand). Légende en 3 blocs : Densités (carrés couleurs), Villes (cercles gradués), Migrations (flèches typées). Échelle graphique (0--100--200 km) et Nord.
\end{enumerate}
\end{methode}

\begin{exoresolu}[Analyse documentaire : Extrait discours YouthConnekt 2023]
\textbf{Document :} \og Le dividende démographique n'est pas un cadeau du ciel. C'est un choix politique. Si nous n'investissons pas aujourd'hui dans l'éducation technique, la santé reproductive et l'entrepreneuriat jeune, les 300 000 qui arrivent chaque année sur le marché du travail deviendront la matière première de l'instabilité. \fg{} (Paraphrase allocution Minjec/Minfopra).

\textbf{Questions :}
\begin{enumerate}
    \item Identifiez les trois piliers d'action cités.
    \item Expliquez la métaphore \og matière première de l'instabilité \fg{} à la lumière du ratio dépendance et du déficit d'emplois formels.
    \item Proposez deux indicateurs quantitatifs pour suivre la \og capture du dividende \fg{} d'ici 2030.
\end{enumerate}

\textbf{Solution :}
\begin{enumerate}
    \item \textbf{Éducation technique/formation} (capital humain productif), \textbf{Santé reproductive} (baisse fécondité, santé maternelle), \textbf{Entrepreneuriat/Emploi} (création richesses, insertion).
    \item \og Matière première de l'instabilité \fg{} : 300 000 jeunes/an sans perspective d'emploi décent + 43\% $<$ 15 ans (armée de réserve future) + inégalités spatiales (Nord/Sud) + frustrations diplômés chômeurs $\to$ risque radicalisation, criminalité urbaine, migrations irrégulières, adhésion groupes armés (Boko Haram, séparatistes). Le \textbf{dividende non capté se transforme en fardeau}.
    \item Indicateurs de suivi 2030 : (1) \textbf{Taux d'emploi formel des 15--35 ans} (cible : $>$ 20\% vs $<$ 10\% actuel) ; (2) \textbf{ICF} (cible : $<$ 3,5) ; (3) \textbf{Taux d'achèvement secondaire technique/professionnel} (cible : $>$ 60\%) ; (4) \textbf{Prévalence contraceptive moderne} (cible : 35\%).
\end{enumerate}
\tcblower
\textbf{Solution détaillée :} Les réponses ci-dessus intègrent les données du cours (ratio dépendance 85\%, déficit 270k emplois/an, ICF 4,2, PC moderne 20\%). La métaphore souligne le basculement possible d'un atout (jeunesse nombreuse) en menace (jeunesse désœuvrée) en l'absence de politiques publiques volontaristes.
\end{exoresolu}

\coursec{Atelier méthodologique : Cartographie, graphiques et analyse de documents (TP)}

Dans la section précédente, nous avons étudié les dynamiques démographiques internes du Cameroun et ses défis à l'horizon 2035. Mais un géographe ne se contente pas de connaître les faits : il doit savoir les \cle{représenter} (cartes, graphiques) et les \cle{argumenter} (analyse de documents, démonstration écrite). Cet atelier pratique vous donne les outils techniques et méthodologiques exigés au Probatoire et au Baccalauréat technique : construire une pyramide des âges, réaliser un croquis choroplèthe, tracer un diagramme d'évolution et rédiger une synthèse argumentée à partir d'un corpus documentaire.

\courssub{TP 1 : Construire une pyramide des âges}

\textbf{Objectif.} Représenter graphiquement la structure par âge et par sexe d'une population, puis comparer trois profils contrastés : Cameroun 2023 (population jeune), France 2023 (population mature), Japon 2023 (population vieillissante).

\begin{definition}[Pyramide des âges]
Une \cle{pyramide des âges} est un graphique en barres horizontales superposées qui représente la répartition d'une population par \textbf{classe d'âge} (généralement quinquennale, c'est-à-dire par tranches de 5 ans) et par \textbf{sexe}. Par convention, les hommes sont placés à \textbf{gauche} (valeurs négatives) et les femmes à \textbf{droite} (valeurs positives), les classes d'âge étant ordonnées de bas en haut (les 0--4 ans en bas).
\end{definition}

\begin{methode}[Construire une pyramide des âges (Excel ou papier millimétré)]
\begin{enumerate}
\item \textbf{Données brutes} : recueillir les effectifs Hommes/Femmes par classe de 5 ans (source : INS, INSEE, Banque mondiale).
\item \textbf{Calcul des pourcentages} : $\% = \dfrac{\text{effectif de la classe}}{\text{population totale}} \times 100$, calculé séparément pour chaque sexe.
\item \textbf{Convention de signe} : Hommes = valeurs \textbf{négatives}, Femmes = valeurs \textbf{positives}.
\item \textbf{Insertion} : dans Excel, Insertion $>$ Graphique à barres groupées (barres horizontales).
\item \textbf{Options} : séries superposées (chevauchement 100 \%), écart entre barres 0 \%, \textbf{inverser l'ordre des catégories} pour placer les 0--4 ans en bas.
\item \textbf{Mise en forme} : code couleur constant (ex. bleu = hommes, rouge = femmes), étiquettes d'axes, titre précis (lieu + date), source obligatoire.
\end{enumerate}
\end{methode}

\begin{attention}[Erreurs fréquentes au TP]
\begin{itemize}
\item Oublier de mettre les \textbf{hommes à gauche} (valeurs négatives) et les \textbf{femmes à droite} : c'est la convention universelle, son non-respect rend le graphique illisible.
\item Ne pas écrire le symbole \% sur l'axe horizontal : on écrit « Population (\%) » en titre d'axe, pas 5,2 \% sur chaque graduation.
\item Utiliser des \textbf{échelles différentes} pour les hommes et les femmes : cela déforme visuellement la comparaison. Les deux axes doivent être symétriques.
\item Oublier le titre, la date ou la source : toute production non sourcée est pénalisée à l'examen.
\end{itemize}
\end{attention}

\begin{exemplebox}[Données pour le TP : structure par âge et sexe en 2023 (\% de la population totale)]
\begin{center}
\begin{tabular}{|c|cc|cc|cc|cc|}\hline
\rowcolor{popblueL}
\textbf{Classe d'âge} & \multicolumn{2}{c|}{\textbf{Cameroun}} & \multicolumn{2}{c|}{\textbf{France}} & \multicolumn{2}{c|}{\textbf{Japon}} & \multicolumn{2}{c|}{\textbf{Total Cameroun}} \\ \cline{2-9}
\rowcolor{popblueL}
 & \textbf{H} & \textbf{F} & \textbf{H} & \textbf{F} & \textbf{H} & \textbf{F} & \textbf{H+F} \\ \hline
0--4 ans     & 5,2 & 5,0 & 2,5 & 2,4 & 1,8 & 1,7 & 10,2 \\
5--9 ans     & 4,8 & 4,7 & 2,6 & 2,5 & 1,9 & 1,8 & 9,5 \\
10--14 ans   & 4,5 & 4,4 & 2,7 & 2,6 & 2,1 & 2,0 & 8,9 \\
15--19 ans   & 4,0 & 3,9 & 2,8 & 2,7 & 2,3 & 2,2 & 7,9 \\
20--24 ans   & 3,5 & 3,4 & 3,0 & 2,9 & 2,5 & 2,4 & 6,9 \\
25--29 ans   & 3,0 & 2,9 & 3,2 & 3,1 & 2,8 & 2,7 & 5,9 \\
30--34 ans   & 2,5 & 2,4 & 3,3 & 3,2 & 3,0 & 2,9 & 4,9 \\
35--39 ans   & 2,0 & 1,9 & 3,4 & 3,3 & 3,2 & 3,1 & 3,9 \\
40--44 ans   & 1,5 & 1,4 & 3,5 & 3,4 & 3,5 & 3,4 & 2,9 \\
45--49 ans   & 1,2 & 1,1 & 3,3 & 3,2 & 3,3 & 3,2 & 2,3 \\
50--54 ans   & 1,0 & 0,9 & 3,0 & 2,9 & 3,0 & 2,9 & 1,9 \\
55--59 ans   & 0,8 & 0,7 & 2,8 & 2,7 & 2,5 & 2,4 & 1,5 \\
60--64 ans   & 0,6 & 0,5 & 2,5 & 2,4 & 2,0 & 1,9 & 1,1 \\
65--69 ans   & 0,4 & 0,4 & 2,0 & 1,9 & 1,8 & 1,7 & 0,8 \\
70--74 ans   & 0,3 & 0,3 & 1,5 & 1,4 & 1,5 & 1,4 & 0,6 \\
75--79 ans   & 0,2 & 0,2 & 1,0 & 0,9 & 1,0 & 0,9 & 0,4 \\
80 ans et +  & 0,1 & 0,15 & 1,2 & 1,5 & 1,2 & 1,8 & 0,25 \\ \hline
\end{tabular}
\end{center}
\legende{Données approximatives issues des estimations INS (Cameroun), INSEE (France) et Bureau des statistiques (Japon) pour 2023. Les totaux H+F pour le Cameroun permettent de vérifier la part des moins de 20 ans (somme des 4 premières lignes = 36,5 \%).}
\end{exemplebox}

\begin{popfigure}
\begin{tikzpicture}
\begin{axis}[
    width=\linewidth,
    height=9cm,
    xbar,
    bar width=3.5pt,
    ytick=data,
    symbolic y coords={0-4,5-9,10-14,15-19,20-24,25-29,30-34,35-39,40-44,45-49,50-54,55-59,60-64,65-69,70-74,75-79,80+},
    yticklabels={0-4,5-9,10-14,15-19,20-24,25-29,30-34,35-39,40-44,45-49,50-54,55-59,60-64,65-69,70-74,75-79,80+},
    xlabel={Population (\% de la population totale)},
    xmin=-6, xmax=6,
    xtick={-5,-4,-3,-2,-1,0,1,2,3,4,5},
    xticklabels={5\%,4\%,3\%,2\%,1\%,0\%,1\%,2\%,3\%,4\%,5\%},
    ylabel={Classe d'âge},
    ylabel style={align=center},
    legend style={at={(0.5,-0.12)},anchor=north,legend columns=-1},
    axis lines*=left,
    y axis line style={opacity=0},
    tickwidth=0pt,
    enlarge y limits=0.02
]
\addplot[fill=popblue!60, draw=popblue!80!black] coordinates {
    (-5.2,0-4) (-4.8,5-9) (-4.5,10-14) (-4.0,15-19) (-3.5,20-24) (-3.0,25-29) (-2.5,30-34) (-2.0,35-39) (-1.5,40-44) (-1.2,45-49) (-1.0,50-54) (-0.8,55-59) (-0.6,60-64) (-0.4,65-69) (-0.3,70-74) (-0.2,75-79) (-0.1,80+)
};
\addlegendentry{Hommes}
\addplot[fill=poppink!60, draw=poppink!80!black] coordinates {
    (5.0,0-4) (4.7,5-9) (4.4,10-14) (3.9,15-19) (3.4,20-24) (2.9,25-29) (2.4,30-34) (1.9,35-39) (1.4,40-44) (1.1,45-49) (0.9,50-54) (0.7,55-59) (0.5,60-64) (0.4,65-69) (0.3,70-74) (0.2,75-79) (0.15,80+)
};
\addlegendentry{Femmes}
\end{axis}
\end{tikzpicture}
\legende{Pyramide des âges du Cameroun (2023, en \% de la population totale). Base très large (plus de 36 \% de moins de 20 ans) et sommet effilé : profil typique d'une population jeune en début de transition démographique. Source : INS, estimations 2023.}
\end{popfigure}

\textbf{Interprétation comparative.} Trois lectures s'imposent :
\begin{itemize}
\item \textbf{Jeunesse} : au Cameroun, les moins de 15 ans représentent environ 37 \% de la population (somme 0-14 ans) ; la base large traduit une natalité encore élevée (ICF proche de 4,5).
\item \textbf{Vieillissement} : au Japon, la base est rétrécie (moins de 12 \% de moins de 15 ans) et le sommet élargi (près de 10 \% de 80 ans et plus) : la pyramide a une forme de « tonneau » inversé.
\item \textbf{Sex-ratio} : à la naissance, il naît environ 105 garçons pour 100 filles, mais la surmortalité masculine fait que les femmes dominent dans les classes âgées (d'où un sommet de pyramide penché à droite, visible au Cameroun comme en France).
\end{itemize}

\courssub{TP 2 : Croquis de la répartition de la population au Cameroun}

\textbf{Objectif.} Réaliser un croquis régional combinant trois modes de représentation : la \cle{choroplèthe} (densités par région), les \cle{symboles proportionnels} (chefs-lieux) et les \cle{flèches} (migrations internes).

\begin{methode}[Choisir les seuils d'une choroplèthe]
\begin{enumerate}
\item Constituer la \textbf{série statistique} : densité = population / superficie, calculée pour chacune des 10 régions (données INS, RGPH 2005 projetées 2023).
\item Observer la \textbf{distribution} : diagramme en bâtons ou boîte à moustaches pour repérer les ruptures.
\item Choisir la méthode de discrétisation :
\begin{itemize}
\item \textbf{Classes équidistantes} : amplitude constante, simple mais masque les ruptures ;
\item \textbf{Quantiles} : même nombre de régions par classe, utile pour comparer ;
\item \textbf{Seuils naturels (Jenks)} : optimise l'homogénéité à l'intérieur des classes et l'hétérogénéité entre classes — méthode recommandée.
\end{itemize}
\item Limiter à \textbf{5 ou 6 classes maximum} pour rester lisible.
\item Utiliser un \textbf{dégradé ordonné} de couleurs (clair $\rightarrow$ foncé) adapté à un ratio comme la densité.
\end{enumerate}
\end{methode}

\textbf{Application.} Les densités régionales camerounaises (estimations 2023) vont de moins de 10 hab/km$^2$ dans l'Est (5,8) et l'Adamaoua (9,2) à plus de 200 hab/km$^2$ dans l'Ouest (210) et le Littoral (180, hors Douala). Une discrétisation en 5 classes (moins de 10 ; 10--25 ; 25--50 ; 50--100 ; plus de 100 hab/km$^2$) fait apparaître le contraste fondamental : un \textbf{Sud-Ouest dense et urbanisé} (Ouest, Littoral, Nord-Ouest, Sud-Ouest, Centre) contre un \textbf{Nord-Est et un Est faiblement peuplés} (Adamaoua, Nord, Extrême-Nord, Est). Pour les villes, on dessine des cercles dont l'\textbf{aire} est proportionnelle à la population : le rayon est donc proportionnel à la \textbf{racine carrée} de l'effectif ($r = k\sqrt{P}$), sinon Douala (environ 3,9 millions d'habitants) et Yaoundé (environ 4,1 millions) écraseraient toute la carte. Enfin, des flèches d'épaisseur proportionnelle figurent les migrations internes (données ENEC 2022 / RGPH) : flux majeurs du Nord, de l'Adamaoua, de l'Est et de l'Ouest vers Yaoundé et Douala.

\begin{popfigure}
\begin{tikzpicture}
% Fond grille
\draw[step=1cm,gray!20] (0,0) grid (12,9);
% Contour Cameroun simplifié (polygone approximatif)
\draw[thick, popdark] (1,1) -- (10,1) -- (11,3) -- (11,7) -- (9,8) -- (6,8) -- (4,7) -- (1,5) -- cycle;
% Choroplèthe : hachures pour simuler les classes
\foreach \x/\y/\pattern in {
    2,2/north east lines, 3,2/north east lines, 4,2/north east lines, % Extrême-Nord, Nord, Adamaoua (Classe 1 <10)
    5,4/north west lines, 6,4/north west lines, % Est, Centre (Classe 2 10-25)
    7,6/horizontal lines, 8,6/horizontal lines, % Sud, Sud-Ouest (Classe 3 25-50)
    7,4/vertical lines, 8,4/vertical lines, % Littoral, Nord-Ouest (Classe 4 50-100)
    6,3/crosshatch % Ouest (Classe 5 >100)} {
    \fill[pattern=\pattern, pattern color=popblue!50!black, opacity=0.6] (\x,\y) rectangle +(1,1);
}
% Villes : symboles proportionnels (rayon ~ sqrt(pop))
% Douala (3.9M) -> r=0.55 ; Yaoundé (4.1M) -> r=0.57 ; Bafoussam (0.5M) -> r=0.2 ; Garoua (0.4M) -> r=0.18
\fill[poporange!80!black, opacity=0.8] (9.5,2.5) circle (0.55cm) node[below=0.1cm, font=\tiny, text=popdark] {Douala};
\fill[poporange!80!black, opacity=0.8] (5.5,5.5) circle (0.57cm) node[above=0.1cm, font=\tiny, text=popdark] {Yaoundé};
\fill[poporange!80!black, opacity=0.8] (4.5,3.5) circle (0.2cm) node[right, font=\tiny, text=popdark] {Bafoussam};
\fill[poporange!80!black, opacity=0.8] (3.5,2.5) circle (0.18cm) node[right, font=\tiny, text=popdark] {Garoua};
\fill[poporange!80!black, opacity=0.8] (8.5,6.5) circle (0.15cm) node[left, font=\tiny, text=popdark] {Buea};
% Flèches migrations
\draw[->, thick, popgreen!60!black] (3,2) .. controls (4,3) and (5,4) .. (5.5,5.5); % Nord -> Yaoundé
\draw[->, thick, popgreen!60!black] (4,3) .. controls (5,3.5) .. (5.5,5.5); % Adamaoua -> Yaoundé
\draw[->, thick, popgreen!60!black] (6,3) -- (5.5,5.5); % Est -> Yaoundé
\draw[->, thick, popgreen!60!black] (4.5,3.5) -- (5.5,5.5); % Ouest -> Yaoundé
\draw[->, thick, popgreen!60!black] (4.5,3.5) -- (9.5,2.5); % Ouest -> Douala
% Légende manuelle (schématique)
\node[anchor=south west, font=\small, text=popdark] at (0.5,8.5) {\textbf{Légende} :};
\node[anchor=west, font=\tiny] at (0.5,8.0) {\textcolor{popblue!50!black}{\rule{0.3cm}{0.15cm}} $<$ 10 hab/km$^2$};
\node[anchor=west, font=\tiny] at (0.5,7.7) {\textcolor{popblue!50!black}{\rule{0.3cm}{0.15cm}} 10 -- 25};
\node[anchor=west, font=\tiny] at (0.5,7.4) {\textcolor{popblue!50!black}{\rule{0.3cm}{0.15cm}} 25 -- 50};
\node[anchor=west, font=\tiny] at (0.5,7.1) {\textcolor{popblue!50!black}{\rule{0.3cm}{0.15cm}} 50 -- 100};
\node[anchor=west, font=\tiny] at (0.5,6.8) {\textcolor{popblue!50!black}{\rule{0.3cm}{0.15cm}} $>$ 100};
\node[anchor=west, font=\tiny] at (0.5,6.4) {\textcolor{poporange!80!black}{\rule{0.3cm}{0.15cm}} Chef-lieu (aire $\propto$ pop.)};
\node[anchor=west, font=\tiny] at (0.5,6.1) {\textcolor{popgreen!60!black}{\rule{0.3cm}{0.15cm}} Flux migratoires internes};
% Nord
\draw[->, thick] (11.5, 8.5) -- (11.5, 8.0) node[right] {\textbf{N}};
\node[anchor=south, font=\small\bfseries] at (6, 0.3) {Répartition de la population au Cameroun (2023)};
\node[anchor=north, font=\tiny] at (6, -0.2) {Source : INS, RGPH projections 2023 ; ENEC 2022.};
\end{tikzpicture}
\legende{Croquis schématique « Répartition de la population au Cameroun » : choroplèthe des densités (5 classes), symboles proportionnels pour les principales villes, flèches des migrations internes. Rappel : titre complet, légende structurée et numérotée, flèche du nord, échelle indicative (ex. 1 cm $\approx$ 100 km), source.}
\end{popfigure}

\begin{aretenir}[Les règles d'or du croquis géographique]
Titre complet (thème + espace + date), légende \textbf{structurée et numérotée} (dans l'ordre de la démonstration), flèche du nord, échelle indicative, source. Le croquis n'est pas une carte muette : chaque figuré sert une idée.
\end{aretenir}

\courssub{TP 3 : Diagramme d'évolution de l'ICF (1960--2023)}

\textbf{Objectif.} Tracer sur un même graphique l'évolution de l'\cle{Indice Conjoncturel de Fécondité} (nombre moyen d'enfants par femme) de quatre pays : Cameroun, France, Chine, Corée du Sud (données Banque mondiale).

\textbf{Protocole.} Sur Excel ou Google Sheets : saisir les années en colonne A (1960, 1970, 1980, 1990, 2000, 2010, 2023) et les ICF des quatre pays en colonnes B à E ; insérer un graphique en courbes avec axes légendés (années en abscisse, ICF en ordonnée), grille légère, légende et titre. Annoter ensuite les \textbf{événements clés} : loi Neuwirth (1967, légalisation de la contraception en France), politique de l'enfant unique en Chine (1979--2015), politique de planning familial en Corée du Sud (années 1960--1990).

\begin{center}
\begin{tabular}{|c|cccc|}\hline
\rowcolor{popblueL}
\textbf{Année} & \textbf{Cameroun} & \textbf{France} & \textbf{Chine} & \textbf{Corée du Sud} \\ \hline
1960 & 6,5 & 2,8 & 6,0 & 6,0 \\
1970 & 6,5 & 2,5 & 5,5 & 4,5 \\
1980 & 6,2 & 1,9 & 2,6 & 2,8 \\
1990 & 5,8 & 1,8 & 2,3 & 1,6 \\
2000 & 5,2 & 1,9 & 1,6 & 1,5 \\
2010 & 4,8 & 2,0 & 1,6 & 1,2 \\
2023 & 4,5 & 1,8 & 1,2 & 0,7 \\ \hline
\end{tabular}
\end{center}

\textbf{Analyse.} Le graphique révèle des \textbf{vitesses de transition très contrastées} :
\begin{itemize}
\item La \textbf{France} a achevé sa transition tôt et doucement : ICF déjà proche de 2,8 en 1960, autour de 1,8 en 2023, sans politique coercitive.
\item La \textbf{Corée du Sud} illustre une transition \textbf{éclair} : de 6 enfants par femme en 1960 à moins de 1 en 2023 — la plus rapide de l'histoire, sous l'effet de l'industrialisation, de l'éducation des filles et des politiques de planning familial.
\item La \textbf{Chine} montre le rôle décisif d'une \textbf{politique volontariste} : chute de 6 à 1,6 entre 1960 et 2000, accélérée par l'enfant unique (1979--2015).
\item Le \textbf{Cameroun} amorce une transition \textbf{lente et inachevée} : d'environ 6,5 en 1960 à 4,5 en 2023 ; le seuil critique de renouvellement des générations (2,1) n'est pas encore atteint.
\end{itemize}

\begin{attention}[Seuil critique à connaître]
Le \textbf{seuil de renouvellement des générations} est fixé à \textbf{2,1 enfants par femme} (et non 2,0, pour compenser la mortalité infantile et la surmortalité masculine). En dessous, sans immigration, la population décroît à terme. Ne confondez pas baisse de l'ICF et baisse immédiate de la population : l'inertie démographique prolonge la croissance plusieurs décennies.
\end{attention}

\courssub{TP 4 : Analyse de documents type Bac (Probatoire)}

\textbf{Corpus proposé (3 à 4 documents)} : une carte choroplèthe de la croissance démographique mondiale ; deux pyramides des âges contrastées (Niger / Japon) ; un tableau d'indicateurs comparés Cameroun / Suède (natalité, mortalité, ICF, espérance de vie, PIB/habitant) ; un extrait de discours (ONU ou Ministre camerounais de la Santé Publique).

\begin{methode}[Répondre à un sujet type Bac (analyse de documents)]
\begin{enumerate}
\item \textbf{Lecture active} : pour CHAQUE document, identifier nature, source, date, échelle et légende.
\item \textbf{Repérage des informations clés} : chiffres marquants, contrastes spatiaux, anomalies.
\item \textbf{Liens inter-documents} : confirmation, contradiction, changement d'échelle.
\item \textbf{Plan thématique} : organiser par idées (2 ou 3 parties), JAMAIS document par document.
\item \textbf{Rédaction} : phrase d'introduction + parties (Idée force + Preuve tirée du document + Connaissance personnelle) + conclusion synthétique.
\end{enumerate}
\end{methode}

\textbf{Questions guidées types} : Décrire (que voit-on ? où ?) ; Expliquer (quels facteurs : économiques, culturels, sanitaires, politiques ?) ; Comparer (quels contrastes entre Niger et Japon, Cameroun et Suède ?) ; Évaluer (quelle politique démographique est adaptée à chaque situation ?) ; Prospective (quelles évolutions probables à l'horizon 2050 ?).

\begin{exoresolu}[Synthèse argumentée (modèle, 200--300 mots)]
\textbf{Sujet} : À partir des documents et de vos connaissances, montrez que le monde fait face à des problèmes démographiques contrastés et que les politiques mises en œuvre diffèrent selon les contextes.
\tcblower
\textbf{Corrigé rédigé (extrait).} Alors que le Niger affiche un ICF de près de 7 enfants par femme, le Japon est descendu sous 1,3 : le monde contemporain connaît des situations démographiques radicalement opposées. Comment expliquer ces contrastes et quelles réponses politiques y sont apportées ?

D'une part, les pays du Sud, comme le Niger ou le Cameroun, connaissent une \textbf{croissance rapide} : natalité élevée (plus de 35 \textperthousand), mortalité en baisse, population très jeune (près de 50 \% de moins de 15 ans au Niger, doc. 2). Cette jeunesse exerce une pression forte sur l'école, l'emploi et la santé. Les politiques y sont \textbf{malthusiennes} : planning familial, éducation des filles, campagnes de contraception, comme le préconise l'ONU (doc. 4).

D'autre part, les pays développés affrontent le \textbf{vieillissement} : au Japon, près de 30 \% de la population a plus de 65 ans, la pyramide est inversée (doc. 2), la population active diminue et les dépenses de santé explosent. Les politiques y sont \textbf{natalistes} (allocations familiales, congés parentaux en Suède) ou \textbf{migratoires} (ouverture contrôlée des frontières).

Ainsi, il n'existe pas un mais des problèmes démographiques, appelant des réponses opposées : freiner les naissances ici, les encourager là. Demain, la transition démographique en cours en Afrique subsaharienne déterminera l'équilibre démographique mondial du XXI\up{e} siècle.
\end{exoresolu}

\courssub{Méthode : la démonstration géographique}

\begin{methode}[Rédiger une démonstration géographique]
\begin{enumerate}
\item \textbf{Accroche} : un chiffre clé ou un paradoxe (ex. « 8 milliards d'humains en 2022, mais deux mondes démographiques opposés »).
\item \textbf{Problématique} : une question qui engage toute la réflexion.
\item \textbf{Annonce du plan} : 2 ou 3 parties, formulées clairement.
\item \textbf{Développement} : pour chaque partie, Idée force + Argument (donnée chiffrée, exemple localisé, mécanisme) + Document (carte, graphique, citation) $\rightarrow$ phrase de lien vers la partie suivante.
\item \textbf{Conclusion} : réponse explicite à la problématique + ouverture (autre échelle, perspective d'avenir).
\end{enumerate}
\end{methode}

\courssub{Les outils numériques du géographe}

\begin{itemize}
\item \textbf{QGIS} (logiciel libre) : cartographie professionnelle — choroplèthes, symboles proportionnels, export PDF vectoriel.
\item \textbf{Excel / Google Sheets} : calculs de pourcentages et de densités, pyramides des âges, courbes d'évolution.
\item \textbf{Gapminder Tools et World Bank Data} : exploration interactive et animée des indicateurs démographiques de tous les pays depuis 1800.
\item \textbf{Google Earth} : visualisation des densités, de l'urbanisation et des contrastes de peuplement à toutes les échelles.
\end{itemize}

\begin{savaistu}[QGIS et Gapminder]
Le logiciel \textbf{QGIS} (gratuit, open source) permet de réaliser des choroplèthes professionnelles avec les fonds de carte GADM (le niveau 1 correspond aux régions du Cameroun) et d'exporter en PDF vectoriel de qualité examen. Le site \textbf{Gapminder Tools} (gapminder.org/tools) anime la transition démographique de tous les pays depuis 1800 : on y voit littéralement les pays « voyager » de la fécondité élevée vers la fécondité faible, à des vitesses très différentes — un excellent support pour réviser le TP 3.
\end{savaistu}

\begin{exercice}[Entraînement : Pyramide simplifiée et comparaison]
À partir du tableau de données de l'\textbf{Exemplebox} (TP 1) :
\begin{enumerate}
\item Calculez la part des moins de 20 ans dans la population totale pour le Cameroun et pour la France.
\item Construisez sur papier millimétré une pyramide des âges \textbf{simplifiée du Cameroun} en 4 grandes classes d'âge : 0--19 ans, 20--39 ans, 40--59 ans, 60 ans et plus. Respectez la convention (hommes à gauche/négatif, femmes à droite/positif), l'échelle symétrique, le titre et la source.
\item Rédigez en 10 lignes maximum une comparaison des deux structures démographiques (Cameroun vs France), en mobilisant les notions de \cle{jeunesse}, de \cle{vieillissement} et de \cle{transition démographique}.
\end{enumerate}
\end{exercice}

\begin{corrige}[Exercice 1]
\begin{enumerate}
\item \textbf{Calculs :}
\begin{itemize}
\item \textbf{Cameroun} : Somme des pourcentages H+F pour les classes 0-4, 5-9, 10-14, 15-19.
$ (5,2+5,0) + (4,8+4,7) + (4,5+4,4) + (4,0+3,9) = 10,2 + 9,5 + 8,9 + 7,9 = \textbf{36,5 \%} $.
\item \textbf{France} : Somme des pourcentages H+F pour les mêmes classes.
$ (2,5+2,4) + (2,6+2,5) + (2,7+2,6) + (2,8+2,7) = 4,9 + 5,1 + 5,3 + 5,5 = \textbf{20,8 \%} $.
\end{itemize}
\item \textbf{Pyramide simplifiée (Cameroun) :} Regrouper les données du tableau par grandes classes (somme H / somme F), reporter en valeurs négatives (H) et positives (F) sur axe gradué symétrique (ex. 1 cm pour 5 \%), barres horizontales superposées, titre « Structure par âge et sexe du Cameroun, 2023 », source « INS, 2023 ».
\begin{center}
\begin{tabular}{|c|cc|c|}\hline
\textbf{Grande classe} & \textbf{Hommes (\%)} & \textbf{Femmes (\%)} & \textbf{Total (\%)} \\ \hline
0--19 ans & 18,5 & 18,0 & 36,5 \\
20--39 ans & 12,0 & 11,6 & 23,6 \\
40--59 ans & 5,5 & 5,1 & 10,6 \\
60 ans et + & 1,6 & 1,55 & 3,15 \\ \hline
\end{tabular}
\end{center}
\item \textbf{Comparaison rédigée :}
La pyramide du Cameroun présente une base très large (36,5 \% de moins de 20 ans) traduisant une natalité encore forte (ICF $\approx$ 4,5) et une mortalité infantile en recul : c'est le profil d'une population \textbf{jeune} en plein \textbf{début de transition démographique} (phase 2--3). Celle de la France a une base étroite (20,8 \% de moins de 20 ans) et un corps élargi aux âges adultes, signe d'une \textbf{transition achevée} (phase 4) et d'un \textbf{vieillissement} structurel (ICF $\approx$ 1,8 < 2,1). Le défi camerounais est l'investissement massif dans l'éducation et l'emploi des jeunes (dividende démographique à capter) ; le défi français est le financement des retraites et de la dépendance par une population active en proportion décroissante.
\end{enumerate}
\end{corrige}

\newpage

\coursec{Activité d'intégration}

\courssub{Objectifs de l'activité}

À l'issue de cette activité d'intégration, l'élève doit être capable de :

\begin{itemize}
\item mobiliser les \cle{indicateurs démographiques} (taux de natalité, de mortalité, d'accroissement naturel, indice synthétique de fécondité, taux de dépendance) pour établir un diagnostic chiffré ;
\item situer le Cameroun dans la \cle{transition démographique} et identifier les problèmes démographiques qui en découlent ;
\item distinguer les grandes catégories de \cle{politiques démographiques} (nataliste, malthusienne, migratoire, sociale) et en évaluer les bilans ;
\item analyser des documents variés (pyramides des âges, cartes de densité, tableaux statistiques, extraits de rapports) ;
\item rédiger et justifier une démarche, une réponse argumentée et une conclusion, dans le cadre d'une prise de parole institutionnelle simulée ;
\item travailler en groupe, débattre de manière contradictoire et construire une synthèse collective.
\end{itemize}

\courssub{Situation}

\textbf{Yaoundé, Conseil des Ministres extraordinaire.} Le Chef de l'État a inscrit à l'ordre du jour la question suivante : \emph{« Quelle politique de population le Cameroun doit-il adopter pour transformer sa dynamique démographique en atout à l'horizon 2035 ? »} Vous êtes les experts convoqués.

\textbf{Les données dont dispose le Conseil (dossier documentaire commun) :}

\begin{center}
\begin{tabularx}{\linewidth}{|l|X|X|}\hline
\rowcolor{popblueL} \textbf{Indicateur} & \textbf{Cameroun (vers 2023)} & \textbf{Repère mondial} \\ \hline
Population totale & environ 28,6 millions d'habitants & 8 milliards \\ \hline
Taux d'accroissement naturel & environ 2,6 \% par an & 0,9 \% par an \\ \hline
Indice synthétique de fécondité (ISF) & environ 4,4 enfants par femme & 2,3 \\ \hline
Taux brut de natalité & environ 35 \textperthousand & 17 \textperthousand \\ \hline
Taux brut de mortalité & environ 9 \textperthousand & 8 \textperthousand \\ \hline
Part des moins de 15 ans & environ 43 \% & 25 \% \\ \hline
Espérance de vie à la naissance & environ 60 ans & 73 ans \\ \hline
Taux d'urbanisation & environ 58 \% & 57 \% \\ \hline
Prévalence contraceptive & environ 20 \% des femmes en union & plus de 60 \% dans les pays développés \\ \hline
\end{tabularx}
\end{center}

\textbf{Documents complémentaires remis à chaque commission :}
\begin{itemize}
\item Document 1 : pyramide des âges du Cameroun 2023 (base très élargie, sommet étroit) et projection 2035 (base encore large, gonflement des classes 15-35 ans) ;
\item Document 2 : carte schématique des densités (fortes densités à l'Ouest, au Littoral et dans l'Extrême-Nord ; faibles densités à l'Est et au Sud) et des métropoles (Douala plus de 3,5 millions, Yaoundé plus de 4 millions) ;
\item Document 3 : tableau des indicateurs régionaux (ISF proche de 6 dans l'Extrême-Nord contre environ 3,5 au Centre ; mortalité maternelle élevée dans les régions septentrionales) ;
\item Document 4 : extraits du PND30 (Plan National de Développement 2020-2030) et de la vision « Cameroun émergent 2035 » ;
\item Document 5 : extraits de rapports UNFPA et Banque mondiale sur le « dividende démographique ».
\end{itemize}

\textbf{Le problème posé :} avec un doublement de la population tous les 27 ans environ, le Cameroun devra scolariser, soigner, nourrir et employer près de 40 millions d'habitants en 2035. Faut-il freiner la natalité ? Investir massivement dans la jeunesse ? Rééquilibrer le peuplement du territoire ? Et avec quel budget ?

\courssub{Consignes}

La classe est répartie en \textbf{cinq commissions ministérielles} : Planification-Économie ; Santé-Affaires sociales ; Éducation-Emploi des jeunes ; Décentralisation-Aménagement du territoire ; Finances-Budget. Chaque commission produit une \textbf{fiche-action d'une page}.

\begin{enumerate}
\item \textbf{Diagnostiquer.} À partir du tableau d'indicateurs et des documents 1 à 3, identifiez et chiffrez les \textbf{trois urgences démographiques} de votre secteur. Calculez au moins un indicateur vous-même (taux de dépendance, taux d'accroissement, évolution en pourcentage).
\item \textbf{Situer.} Précisez à quelle étape de la transition démographique se trouve le Cameroun et ce que cela implique pour votre secteur.
\item \textbf{Proposer.} Formulez \textbf{deux mesures concrètes} : coût estimé, cibles précises (populations, régions prioritaires), indicateur de suivi chiffré (exemple : « porter la prévalence contraceptive de 20 \% à 35 \% en 2030 »).
\item \textbf{Arbitrer.} Identifiez au moins un \textbf{conflit d'arbitrage} avec une autre commission (budget santé contre budget éducation, incitations à la baisse de la fécondité contre autonomie des femmes, investissements urbains contre monde rural) et proposez une solution de compromis.
\item \textbf{Débattre et synthétiser.} En plénière, chaque commission présente sa fiche en 3 minutes. Le « Premier Ministre » (le professeur) modère le débat contradictoire. La classe rédige ensuite collectivement une \textbf{« Déclaration de politique de population 2025-2035 »} de 10 lignes maximum, hiérarchisant les priorités.
\end{enumerate}

\courssub{Ressources à mobiliser}

\begin{aretenir}[Boîte à outils du démographe]
\begin{itemize}
\item Taux d'accroissement naturel (\textperthousand) $=$ taux brut de natalité $-$ taux brut de mortalité ;
\item Taux d'accroissement total $=$ accroissement naturel $+$ solde migratoire ;
\item Taux de dépendance $= \dfrac{\text{jeunes (0-14 ans)} + \text{vieux (65 ans et +)}}{\text{actifs potentiels (15-64 ans)}} \times 100$ ;
\item Temps de doublement : à 2,6 \% par an, une population double en environ $70/2{,}6 \approx 27$ ans ;
\item Étapes de la transition démographique : régime traditionnel (natalité et mortalité élevées) $\rightarrow$ transition (mortalité en baisse, forte croissance) $\rightarrow$ régime moderne (natalité et mortalité faibles) ;
\item Types de politiques démographiques : nataliste, malthusienne (limitation des naissances), migratoire, sociale (santé, éducation, statut des femmes).
\end{itemize}
\end{aretenir}

\courssub{Démarche et corrigé commenté}

\begin{exoresolu}[Exemple de fiche-action : commission Santé-Affaires sociales]
Énoncé : produire la fiche-action de la commission Santé selon les consignes 1 à 4.
\tcblower
\textbf{Étape 1 — Diagnostic chiffré.} Trois urgences sont identifiées à partir des documents :
\begin{itemize}
\item \emph{Urgence 1 : une fécondité élevée et mal maîtrisée.} ISF de 4,4 enfants par femme (jusqu'à 6 dans l'Extrême-Nord) et prévalence contraceptive d'environ 20 \% seulement : la demande de planning familial est insatisfaite.
\item \emph{Urgence 2 : une mortalité encore forte.} Taux brut de mortalité de 9 \textperthousand, espérance de vie de 60 ans (13 ans de moins que la moyenne mondiale), mortalité maternelle et infantile élevées, surtout dans les régions septentrionales (document 3).
\item \emph{Urgence 3 : une pression sanitaire croissante.} Avec 2,6 \% d'accroissement naturel, il naît chaque année environ $28{,}6 \text{ millions} \times 35/1000 \approx 1$ million de bébés : maternités, vaccins et personnel de santé doivent suivre ce rythme.
\end{itemize}

\textbf{Étape 2 — Situation dans la transition.} Le Cameroun est au début de la phase de transition : la mortalité a nettement baissé depuis l'indépendance mais la natalité reste élevée (35 \textperthousand). Conséquence pour la santé : une croissance démographique rapide qui exige d'élargir l'offre de soins plus vite que la population ne croît.

\textbf{Étape 3 — Deux mesures concrètes.}
\begin{itemize}
\item \emph{Mesure 1 : programme national de planning familial.} Objectif : porter la prévalence contraceptive de 20 \% à 35 \% en 2030. Cibles : femmes en union des régions Extrême-Nord, Nord et Adamaoua (ISF supérieur à 5). Coût estimé : 15 milliards de FCFA par an (subvention des contraceptifs, formation de 2 000 sages-femmes). Indicateur de suivi : prévalence contraceptive mesurée par enquête tous les deux ans.
\item \emph{Mesure 2 : couverture santé maternelle et infantile.} Objectif : ramener la mortalité maternelle de moitié d'ici 2035 et faire passer l'espérance de vie de 60 à 65 ans. Moyens : gratuité effective de l'accouchement assisté dans les districts ruraux, 500 centres de santé intégrés réhabilités. Coût estimé : 40 milliards de FCFA par an. Indicateur : taux d'accouchements assistés par personnel qualifié.
\end{itemize}

\textbf{Étape 4 — Arbitrage.} Conflit identifié avec la commission Finances : les 55 milliards annuels demandés réduisent la marge du budget éducation. Compromis proposé : mutualiser les infrastructures (les centres de santé communautaires accueillent aussi les séances d'éducation à la santé reproductive des adolescentes, en lien avec la commission Éducation), ce qui réduit le coût de 10 \% et renforce l'articulation intersectorielle.

\textbf{Étape 5 — Argumentaire de plénière (résumé oral).} « Sans maîtrise de la fécondité ni baisse de la mortalité maternelle, aucun dividende démographique n'est possible : chaque franc investi dans la santé des mères et le planning familial économise demain des dépenses de scolarisation et d'emploi. La santé est la condition préalable de toute la stratégie 2035. »
\end{exoresolu}

\begin{methode}[Méthode commune à toutes les commissions]
\begin{enumerate}
\item Lire le tableau d'indicateurs et repérer les valeurs extrêmes (fortes ou faibles) qui concernent votre secteur ;
\item croiser avec les documents cartographiques : localiser les régions prioritaires ;
\item calculer au moins un indicateur simple pour étayer le diagnostic ;
\item formuler des mesures SMART : chiffrées, datées, localisées, avec un indicateur de suivi ;
\item anticiper les objections des autres commissions et préparer un compromis.
\end{enumerate}
\end{methode}

\begin{attention}[Pièges à éviter pendant le débat]
\begin{itemize}
\item Ne pas confondre \cle{politique malthusienne} coercitive (stérilisations forcées, enfant unique) et politique volontariste de planning familial fondée sur le libre choix : la seconde est la seule acceptable et réaliste au Cameroun ;
\item ne pas présenter la forte natalité comme un « problème » en soi : c'est le décalage entre croissance démographique et croissance économique qui crée la difficulté ;
\item éviter les mesures non chiffrées (« améliorer les hôpitaux ») : toute proposition doit comporter un coût, une cible et un indicateur ;
\item ne pas oublier la dimension territoriale : les contrastes régionaux (Extrême-Nord contre Centre) imposent des mesures différenciées.
\end{itemize}
\end{attention}

\courssub{Bilan de l'activité}

Cette simulation a permis de mettre en évidence plusieurs acquis essentiels de la leçon :

\begin{itemize}
\item la \textbf{complémentarité des indicateurs} : aucun diagnostic sérieux ne repose sur un seul chiffre ; il faut croiser natalité, mortalité, fécondité, structure par âge et répartition spatiale ;
\item la \textbf{position du Cameroun dans la transition démographique} : mortalité en baisse, natalité encore élevée, d'où une croissance rapide (2,6 \% par an, doublement en 27 ans) et une population très jeune (43 \% de moins de 15 ans) ;
\item le \textbf{caractère intersectoriel des politiques démographiques} : santé, éducation, emploi, aménagement et budget sont indissociables — une baisse de la fécondité passe autant par la scolarisation des filles que par le planning familial ;
\item la notion de \cle{dividende démographique} : une population jeune n'est un atout que si l'État investit massivement et tôt dans la santé, l'éducation et l'emploi ; sinon elle devient une charge (taux de dépendance élevé, chômage des jeunes, exode rural et urbanisation désordonnée) ;
\item la pratique des \textbf{savoir-faire} : analyser des documents variés, calculer des indicateurs, argumenter à l'oral et à l'écrit, rédiger une synthèse collective justifiée.
\end{itemize}

La « Déclaration de politique de population 2025-2035 » produite par la classe constitue la trace écrite finale : elle montre qu'une politique démographique cohérente n'est ni purement malthusienne ni purement nataliste, mais une stratégie globale adaptée aux réalités d'un pays en pleine transition démographique.

\newpage

\coursec{Méthodes et exercices résolus}

\coursec{Les outils du démographe : indicateurs et transition démographique}

\courssub{Les indicateurs démographiques de base}

\begin{definition}[Indicateurs de flux et de stock]
La démographie utilise des indicateurs standardisés pour comparer les populations dans l'espace et dans le temps.
\begin{itemize}
\item \textbf{Taux de natalité (TN)} : nombre de naissances vivantes pour 1\,000 habitants (pour mille, \textperthousand) sur une année. $TN = \dfrac{\text{Naissances}}{\text{Population totale}} \times 1000$.
\item \textbf{Taux de mortalité (TM)} : nombre de décès pour 1\,000 habitants sur une année. $TM = \dfrac{\text{Décès}}{\text{Population totale}} \times 1000$.
\item \textbf{Taux d'accroissement naturel (TAN)} : différence entre natalité et mortalité. $TAN = TN - TM$ (en \textperthousand). En pourcentage : $TAN(\%) = \dfrac{TN - TM}{10}$.
\item \textbf{Indice conjoncturel de fécondité (ICF ou IF)} : nombre moyen d'enfants qu'aurait une femme si elle traversait sa vie féconde soumise aux taux de fécondité par âge de l'année considérée. Seuil de renouvellement des générations : \textbf{2,1 enfants par femme}.
\item \textbf{Densité de population} : rapport population / superficie (hab/km\textsuperscript{2}). $D = \dfrac{P}{S}$.
\item \textbf{Temps de doublement (approché)} : $t = \dfrac{70}{TAN(\%)}$ (règle du 70).
\end{itemize}
\end{definition}

\begin{aretenir}[Ordres de grandeur de référence (monde, vers 2023)]
\begin{center}
\begin{tabularx}{\linewidth}{|l|X|X|}\hline
\textbf{Indicateur} & \textbf{Monde} & \textbf{Afrique subsaharienne} \\ \hline
Taux de natalité & $\approx 17\,\textperthousand$ & $\approx 34\,\textperthousand$ \\
Taux de mortalité & $\approx 7\,\textperthousand$ & $\approx 8\,\textperthousand$ \\
TAN & $\approx 1\,\%$ & $\approx 2,6\,\%$ \\
Indice de fécondité & $\approx 2,3$ & $\approx 4,6$ \\
\% moins de 15 ans & 25\,\% & 43\,\% \\
\% plus de 65 ans & 10\,\% & 3\,\% \\ \hline
\end{tabularx}
\end{center}
\legende{Source : ONU, Perspectives de la population mondiale 2022. Valeurs arrondies.}
\end{aretenir}

\courssub{Le modèle de la transition démographique}

\begin{propriete}[Phases de la transition démographique (modèle classique)]
La transition démographique décrit le passage d'un régime ancien (natalité forte, mortalité forte, croissance faible) à un régime moderne (natalité faible, mortalité faible, croissance faible) via une phase intermédiaire de croissance rapide.
\begin{enumerate}
\item \textbf{Phase 1 (Régime ancien)} : Natalité forte ($\approx 35-45\,\textperthousand$), mortalité forte et variable ($\approx 30-40\,\textperthousand$), croissance quasi nulle. Population jeune, espérance de vie faible ($< 40$ ans).
\item \textbf{Phase 2 (Début de transition)} : Natalité forte, \textbf{mortalité chute brutalement} (progrès sanitaires, nutrition, hygiène). $\Rightarrow$ \textbf{Explosion démographique} (TAN $> 2\,\%$). Ex. : Afrique subsaharienne actuelle.
\item \textbf{Phase 3 (Transition avancée)} : \textbf{Natalité baisse} (urbanisation, scolarisation filles, contraception, coût enfants), mortalité faible. Croissance ralentit. Ex. : Amérique latine, Asie du Sud-Est.
\item \textbf{Phase 4 (Régime moderne)} : Natalité faible ($\approx 10-14\,\textperthousand$), mortalité faible ($\approx 8-10\,\textperthousand$), croissance faible ou nulle. Vieillissement. Ex. : Europe, Amérique du Nord, Japon.
\item \textbf{Phase 5 (Hypothétique / Post-transition)} : Natalité \emph{durablement} sous le seuil de renouvellement ($< 1,5$), mortalité très faible, population décroît et vieillit fortement. Ex. : Corée du Sud (0,72 en 2023), Italie, Espagne.
\end{enumerate}
\end{propriete}

\begin{popfigure}
\begin{tikzpicture}[scale=0.8, font=\small]
% Axes
\draw[->] (0,0) -- (11,0) node[right] {Temps / Développement};
\draw[->] (0,0) -- (0,5) node[above] {Taux (\textperthousand)};
% Ticks
\foreach \y/\ylab in {1/10, 2/20, 3/30, 4/40, 4.5/45} \draw (0,\y) -- (-0.1,\y) node[left] {\ylab};
% Courbe Mortalité (bleu)
\draw[popblue, thick, domain=0:10, samples=100, smooth] plot (\x, {4.5*exp(-\x/2.5) + 0.5});
\node[popblue, align=center] at (2,4.2) {Mortalité};
% Courbe Natalité (rouge)
\draw[poporange, thick, domain=0:10, samples=100, smooth] plot (\x, {4.5*exp(-\x/5) + 0.8});
\node[poporange, align=center] at (7,2.5) {Natalité};
% Zones phases
\draw[dashed, popmuted] (2.5,0) -- (2.5,4.5);
\draw[dashed, popmuted] (5.5,0) -- (5.5,4.5);
\draw[dashed, popmuted] (8.5,0) -- (8.5,4.5);
\node[align=center, popmuted] at (1.25,4.7) {Phase 1\\(Ancien)};
\node[align=center, popgreen] at (4,4.7) {Phase 2\\(Explosion)};
\node[align=center, poppurple] at (7,4.7) {Phase 3\\(Baisse natalité)};
\node[align=center, popblue] at (9.75,4.7) {Phase 4\\(Moderne)};
% Flèche accroissement
\draw[popgreen, <->, thick] (4, {4.5*exp(-4/2.5)+0.5}) -- (4, {4.5*exp(-4/5)+0.8}) node[midway, right] {TAN fort};
\draw[poporange, <->, thick] (8.5, {4.5*exp(-8.5/2.5)+0.5}) -- (8.5, {4.5*exp(-8.5/5)+0.8}) node[midway, right] {TAN faible};
\end{tikzpicture}
\legende{Schéma simplifié de la transition démographique : écart entre natalité et mortalité = accroissement naturel.}
\end{popfigure}

\courssub{Application : Le cas du Cameroun (2023)}

\begin{exoresolu}[Calcul des indicateurs démographiques : le cas du Cameroun]
\textbf{Énoncé.} En 2023, le Cameroun compte environ \num{29000000} habitants sur une superficie de \SI{475442}{km^2}. On y enregistre environ \num{1044000} naissances vivantes et \num{261000} décès dans l'année.
\begin{enumerate}
\item Calculez le taux de natalité et le taux de mortalité du Cameroun en 2023.
\item Calculez le taux d'accroissement naturel en \textperthousand puis en \%. Interprétez.
\item Calculez la densité de population et le temps de doublement de la population camerounaise.
\end{enumerate}
\tcblower
\textbf{Solution détaillée.}

\textbf{1. Taux de natalité et taux de mortalité.}
\begin{align*}
TN &= \dfrac{1044000}{29000000} \times 1000 = 36,0\,\text{\textperthousand} \\
TM &= \dfrac{261000}{29000000} \times 1000 = 9,0\,\text{\textperthousand}
\end{align*}
Le taux de natalité camerounais (36\,\textperthousand) est très élevé : il est plus du double de la moyenne mondiale (environ 17\,\textperthousand). Le taux de mortalité (9\,\textperthousand) est proche de la moyenne mondiale (7\,\textperthousand), traduisant une baisse de la mortalité non encore suivie d'une baisse équivalente de la natalité : le Cameroun est en \textbf{phase 2 de la transition démographique}.

\textbf{2. Taux d'accroissement naturel.}
\[ TAN = TN - TM = 36 - 9 = 27\,\text{\textperthousand} = 2,7\,\% \]
\textbf{Interprétation} : un TAN de 2,7\,\% est supérieur au seuil de 2\,\% qui caractérise une croissance démographique rapide. La population camerounaise augmente donc vite, essentiellement par excédent naturel (les naissances dépassent largement les décès). Vérification : $2,7\,\%$ de \num{29000000} donne environ \num{783000} habitants supplémentaires par an, ce qui correspond bien à $1044000 - 261000 = 783000$.

\textbf{3. Densité et temps de doublement.}
\[ D = \dfrac{29000000}{475442} \approx 61\ \text{hab/km}^2 \]
Cette densité moyenne masque de forts contrastes internes (plus de 300 hab/km\textsuperscript{2} dans l'Ouest, moins de 10 hab/km\textsuperscript{2} dans l'Est et l'Adamaoua).
\[ t = \dfrac{70}{2,7} \approx 26\ \text{ans} \]
\textbf{Conclusion} : au rythme actuel, la population du Cameroun doublerait en environ 26 ans (vers 2049), ce qui pose des défis majeurs en matière d'écoles, d'emplois, de logements et de services de santé.
\end{exoresolu}

\coursec{Les grands problèmes démographiques mondiaux : diagnostics contrastés}

\courssub{Une croissance mondiale qui ralentit mais reste forte en valeur absolue}

\begin{exoresolu}[Exploitation d'un tableau statistique : l'explosion démographique mondiale]
\textbf{Énoncé.} Population mondiale : 1 milliard en 1800 ; 1,6 milliard en 1900 ; 6,1 milliards en 2000 ; 8 milliards en novembre 2022.
\begin{enumerate}
\item Calculez le taux de variation de la population mondiale entre 1900 et 2000.
\item Calculez le taux de variation entre 2000 et 2022. Comparez et interprétez.
\end{enumerate}
\tcblower
\textbf{Solution détaillée.}

\textbf{1. Entre 1900 et 2000 (le « siècle de l'explosion ») :}
\[ \text{Taux de variation} = \dfrac{V_f - V_i}{V_i} \times 100 = \dfrac{6,1 - 1,6}{1,6} \times 100 = \dfrac{4,5}{1,6} \times 100 = 281,25\,\% \]
La population mondiale a été presque multipliée par 4 en un siècle ($6,1 / 1,6 \approx 3,8$). Le taux de croissance annuel moyen a culminé vers 1965-1970 à \textbf{2,1\,\%} (pic historique).

\textbf{2. Entre 2000 et 2022 (ralentissement relatif) :}
\[ \text{Taux de variation} = \dfrac{8 - 6,1}{6,1} \times 100 = \dfrac{1,9}{6,1} \times 100 \approx 31,1\,\% \]

\textbf{Comparaison et interprétation} : la croissance reste forte en valeur absolue (+1,9 milliard en 22 ans, soit ~86 millions/an), mais le rythme relatif ralentit nettement : 31\,\% en 22 ans contre 281\,\% en 100 ans. Le taux annuel moyen est passé de $\approx 1,3\,\%$ (sur le siècle) à $\approx 1,2\,\%$ (2000-2010) puis $< 1\,\%$ aujourd'hui (0,8\,\% en 2023). Ce ralentissement traduit l'entrée de la plupart des pays du Sud (Asie, Amérique latine) dans les phases 3 et 4 de la transition (baisse de la fécondité), tandis que l'Afrique subsaharienne reste en phase 2.
\textbf{Conclusion} : la population mondiale continue d'augmenter, mais de moins en moins vite. Les projections de l'ONU (variante médiane) prévoient une stabilisation autour de \textbf{10,4 milliards vers 2080-2100}.
\end{exoresolu}

\courssub{Deux problèmes démographiques opposés : explosion au Sud, vieillissement au Nord}

\begin{aretenir}[Diagnostics contrastés au Probatoire]
\begin{itemize}
\item \textbf{Pays du Sud (Afrique, partie de l'Asie) :} \textbf{Croissance rapide} (TAN $> 2\,\%$, IF $> 3$). \emph{Conséquences} : pression sur l'éducation (scolarisation massive), la santé, l'emploi (insertion des jeunes), l'environnement, l'urbanisation non maîtrisée. \emph{Dividende démographique possible} si investissement dans capital humain.
\item \textbf{Pays du Nord (Europe, Japon, Amérique du Nord, Chine) :} \textbf{Vieillissement et/ou déclin} (IF $< 2,1$, part $> 65$ ans $> 20\,\%$). \emph{Conséquences} : pénurie de main-d'œuvre, financement des retraites et de la santé, perte de dynamisme économique, « piège à faible fécondité ».
\item \textbf{Déséquilibres spatiaux internes} : Métropolisation extrême (mégapoles $> 10$ M hab) vs désertification rurale ; littoraux surpeuplés vs arrière-pays vides.
\end{itemize}
\end{aretenir}

\begin{savaistu}[Malthus, Boserup et la prospective]
\emph{Thomas Malthus} (1798) : la population croît géométriquement, les ressources arithmétiquement $\Rightarrow$ catastrophes (famine, guerre) ou frein moral. \emph{Esther Boserup} (1965) : la pression démographique \emph{stimule} l'innovation technique et l'intensification agricole. Le débat reste ouvert : la transition démographique (baisse choisie de la fécondité) a jusqu'ici évité le scénario malthusien global, mais les défis climatiques et ressources (eau, sols, biodiversité) réactualisent la question des limites planétaires.
\end{savaistu}

\coursec{Dossier 1 : Les politiques démographiques – Types, acteurs et bilans}

\courssub{Typologie des politiques démographiques}

\begin{definition}[Types de politiques démographiques]
\begin{itemize}
\item \textbf{Politique nataliste (pro-nataliste)} : objectif d'\emph{augmenter} la natalité/fécondité. Outils : allocations familiales, quotient familial, congés parentaux rémunérés, crèches publiques, aides au logement, primes à la naissance. Cible : pays à fécondité très basse (France, Suède, Hongrie, Japon, Corée du Sud, Russie).
\item \textbf{Politique antinataliste / de limitation des naissances (néo-malthusienne)} : objectif de \emph{réduire} la fécondité. Outils : planning familial (accès contraception), éducation/scolarisation des filles (levier n°1), relèvement âge légal mariage, campagnes d'information, incitations financières (ex. prime à la stérilisation), contrainte (ex. politique de l'enfant unique en Chine 1979-2015). Cible : pays à croissance rapide (Inde années 70, Chine, pays africains via programmes UNFPA).
\item \textbf{Politique migratoire} : régulation des flux (immigration choisie pour combler pénuries, ou restriction). Agit sur la structure par âge et la main-d'œuvre à court terme.
\item \textbf{Politique de redistribution spatiale / d'aménagement du territoire} : agir sur la répartition (nouvelles capitales, pôles d'équilibre, développement rural) pour désengorger les métropoles et occuper l'espace.
\end{itemize}
\end{definition}

\begin{propriete}[Acteurs des politiques démographiques]
\begin{itemize}
\item \textbf{État} : législateur, financeur principal, fixateur d'objectifs (Plan national de développement).
\item \textbf{Collectivités territoriales} : mise en œuvre locale (crèches, centres de santé, planning familial).
\item \textbf{ONG et associations} : plaidoyer, services de proximité, éducation communautaire (ex. CAMNAFAW au Cameroun).
\item \textbf{Organisations internationales} : \textbf{UNFPA} (Fonds des Nations Unies pour la population) — chef de file technique/financier ; \textbf{OMS}, \textbf{Banque mondiale}, \textbf{USAID}, \textbf{Union africaine} (Agenda 2063, Dividende démographique).
\end{itemize}
\end{propriete}

\courssub{Bilans contrastés et leçons}

\begin{exoresolu}[Justifier une politique démographique : deux cas contrastés]
\textbf{Énoncé.} Le pays A a un indice de fécondité de 5,6 enfants par femme et un TAN de 3\,\%. Le pays B a un indice de fécondité de 1,3 enfant par femme et 28\,\% de plus de 60 ans. Pour chaque pays, proposez et justifiez le type de politique démographique approprié, en citant un exemple réel.
\tcblower
\textbf{Solution détaillée.}

\textbf{Pays A — Diagnostic} : fécondité de 5,6 (très supérieure au seuil de renouvellement de 2,1) et TAN de 3\,\%, soit un doublement de la population en $\dfrac{70}{3} \approx 23$ ans. Structure par âge : base large (jeunesse). Problème : croissance trop rapide, pression sur services sociaux.
\textbf{Politique justifiée} : \textbf{politique de limitation des naissances (antinataliste)} fondée sur le planning familial (accès contraception), l'éducation et la scolarisation des filles — car la scolarisation secondaire retarde l'âge au premier mariage et réduit la fécondité de manière durable — et l'autonomisation économique des femmes.
\emph{Exemple réel} : la \textbf{Chine} avec la politique de l'enfant unique (1979-2015), qui a fait passer la fécondité d'environ 6 à moins de 1,7 enfant par femme (baisse rapide) mais a provoqué un vieillissement accéléré, un déséquilibre des sexes à la naissance (sélection prénatale) et une contraction de la main-d'œuvre actuelle. \emph{Autre exemple} : les programmes de planning familial soutenus par l'UNFPA en Afrique (ex. Rwanda, Éthiopie) avec des résultats plus progressifs et respectueux des droits.

\textbf{Pays B — Diagnostic} : fécondité de 1,3 (très sous le seuil de renouvellement) et 28\,\% de plus de 60 ans : le problème est le déclin démographique, le vieillissement structurel, la pénurie de main-d'œuvre et l'insoutenabilité du système de retraite par répartition.
\textbf{Politique justifiée} : \textbf{politique nataliste} (allocations familiales universelles, congés parentaux longs et bien rémunérés, réseau dense de crèches publiques, aides au logement pour les familles nombreuses) \emph{éventuellement complétée par une politique migratoire d'accueil de travailleurs qualifiés}.
\emph{Exemple réel} : la \textbf{France}, dont le système d'allocations familiales (quotient familial), le congé parental et le réseau de crèches/écoles maternelles gratuites maintiennent une fécondité proche de 1,8 (among the highest in EU), bien que toujours sous le seuil de renouvellement. La \textbf{Suède} (congé parental partagé, garde d'enfants subventionnée) obtient des résultats similaires.

\textbf{Conclusion} : le choix d'une politique démographique dépend du diagnostic établi par les indicateurs : limitation des naissances là où la croissance est trop rapide (Pays A), soutien à la natalité là où la population décline et vieillit (Pays B). L'efficacité repose sur la cohérence du paquet de mesures (éducation + santé + emploi + logement) et la durée (20-30 ans pour voir les effets structurels).
\end{exoresolu}

\begin{savaistu}[La politique de l'enfant unique en Chine : un cas d'école]
Instaurée en 1979 par Deng Xiaoping pour freiner la croissance (population passée de 540 M en 1949 à 970 M en 1979). \textbf{Résultats} : 400 millions de naissances « évitées » (selon Pékin), fécondité chute à 1,5-1,6. \textbf{Coûts sociaux} : vieillissement rapide (dépendance personnes âgées), « 4-2-1 » (un enfant supporte 2 parents + 4 grands-parents), déséquilibre sexes (118 garçons pour 100 filles en 2010), « femmes surnuméraires » non enregistrées. \textbf{Abandon progressif} : 2 enfants (2016), 3 enfants (2021), suppression des limites (2023) + mesures natalistes. \emph{Leçon} : la contrainte coercitive a des effets rapides mais des externalités négatives lourdes ; les politiques incitatives respectueuses des droits sont plus lentes mais plus durables.
\end{savaistu}

\coursec{Focus Cameroun : Dynamiques internes, défis et prospective 2035}

\courssub{Une population jeune, urbaine et inégalement répartie}

\begin{aretenir}[Portrait démographique du Cameroun (estimations 2023-2024)]
\begin{itemize}
\item \textbf{Population} : $\approx 29$ millions (5\textsuperscript{e} Afrique, 1\textsuperscript{er} Afrique centrale).
\item \textbf{Croissance} : TAN $\approx 2,7\,\%$ (doublement $\approx 26$ ans). Projection 2035 : $\approx 35-37$ millions ; 2050 : $\approx 50$ millions.
\item \textbf{Structure par âge} : $\approx 42\,\%$ de moins de 15 ans ; $\approx 3,5\,\%$ de plus de 65 ans. \textbf{Indice de dépendance des jeunes} très fort ($\approx 80\,\%$).
\item \textbf{Urbanisation} : $\approx 58\,\%$ (2023), +3,5\,\%/an. \textbf{Douala} ($\approx 4$ M) et \textbf{Yaoundé} ($\approx 4$ M) concentrent $\approx 30\,\%$ de la population urbaine. Exode rural massif.
\item \textbf{Répartition} : Ouest (Grassfields) et Littoral (Douala) : densités $> 150-300$ hab/km\textsuperscript{2}. Est, Adamaoua, Nord, Extrême-Nord (hors Maroua) : $< 20$ hab/km\textsuperscript{2}.
\item \textbf{Fécondité} : IF $\approx 4,6$ (2022), en baisse lente (5,1 en 2004). Fort différentiel : 3,5 en milieu urbain vs 5,5 en milieu rural ; 2,8 pour femmes secondaire+ vs 6,0 sans instruction.
\end{itemize}
\end{aretenir}

\begin{popfigure}
\begin{tikzpicture}[scale=0.65, font=\small]
% Contour Cameroun très simplifié (polygone)
\draw[popdark, thick] (0,0) -- (8,0) -- (9,2) -- (9,5) -- (7,7) -- (5,8) -- (4,7) -- (2,6) -- (0,4) -- cycle;
% Villes
\fill[poporange] (2.5,1.5) circle (0.12) node[below right] {Douala};
\fill[poporange] (4.5,3.5) circle (0.12) node[above right] {Yaoundé};
\fill[popblue] (6.5,5.5) circle (0.08) node[above] {Garoua};
\fill[popblue] (7.5,2.5) circle (0.08) node[right] {Maroua};
\fill[popblue] (1.5,5.5) circle (0.08) node[left] {Bertoua};
\fill[popblue] (3.5,6.5) circle (0.08) node[above] {Ngaoundéré};
% Zones de densité (hachures / couleur)
% Ouest : haute densité
\fill[popgreen!30, pattern=north east lines, pattern color=popgreen] (1.5,1) -- (4,1) -- (4.5,3) -- (2,3.5) -- cycle;
% Nord/Extrême-Nord : moyenne
\fill[popgold!30, pattern=north west lines, pattern color=popgold] (5.5,2) -- (8.5,2) -- (9,5) -- (6,5.5) -- cycle;
% Est/Adamaoua : faible
\fill[popblue!15] (0,0) -- (2,3.5) -- (4.5,3) -- (4,7) -- (2,6) -- (0,4) -- cycle;
% Légende manuelle
\node[align=left, anchor=south west, fill=popcream, draw=popline, inner sep=3pt] at (9.5,0) {
\textbf{Légende densité}\\
\rule{5pt}{5pt} \textcolor{popgreen}{> 150 hab/km² (Ouest, Littoral)}\\
\rule{5pt}{5pt} \textcolor{popgold}{20-150 hab/km² (Nord, Ouest périph.)}\\
\rule{5pt}{5pt} \textcolor{popblue!50}{< 20 hab/km² (Est, Adamaoua, Sud)}\\
$\bullet$ \textcolor{poporange}{Métropoles > 3 M hab}\\
$\bullet$ \textcolor{popblue}{Villes moyennes}
};
% Flèches exode rural
\draw[poporange, thick, ->, >=Stealth] (1,5) .. controls (2,3) .. (3,3.5) node[midway, above, sloped, font=\tiny] {Exode rural};
\draw[poporange, thick, ->, >=Stealth] (6,4) .. controls (5,3.5) .. (4.5,3.5) node[midway, above, sloped, font=\tiny] {Exode rural};
% Nord
\draw[->, >=Stealth, popmuted] (9.5,6.5) -- (10,7) node[right, font=\tiny] {Nord};
\end{tikzpicture}
\legende{Croquis schématique : répartition de la population et flux migratoires internes au Cameroun.}
\end{popfigure}

\courssub{Défis et prospective « Vision 2035 »}

\begin{exemplebox}[Les 4 défis majeurs pour 2035]
\begin{enumerate}
\item \textbf{Dividende démographique ou bombe à retardement ?} 300\,000 à 400\,000 jeunes arrivent sur le marché du travail chaque an. L'économie formelle en absorbe $< 10\,\%$. Risque : chômage de masse, instabilité. \emph{Réponse} : formation professionnelle technique, appui PME, économie numérique, investissement massif dans l'éducation qualité.
\item \textbf{Transition sanitaire inachevée} : mortalité maternelle (406 pour 100\,000 naissances vivantes, 2018), mortalité infantile (48‰), paludisme, VIH/SIDA, montée des maladies non transmissibles. \emph{Réponse} : Couverture Santé Universelle (CSU), renforcement districts de santé.
\item \textbf{Urbanisation non maîtrisée} : habitats précaires (bidonvilles), inondations, insécurité alimentaire urbaine, gestion déchets. \emph{Réponse} : Plan Directeur d'Aménagement (PDA) Douala/Yaoundé, villes secondaires (Bafoussam, Bamenda, Garoua, Maroua, Bertoua, Ebolowa) comme pôles d'équilibre.
\item \textbf{Pression sur ressources et environnement} : déforestation (bois énergie, agriculture itinérante), érosion sols, conflits agro-pasteurs (Nord), surexploitation halieutique. \emph{Réponse} : reboisement, agroforesterie, énergies renouvelables, code forestier.
\end{enumerate}
\end{exemplebox}

\begin{methode}[Analyser une situation démographique concrète et proposer une réponse adaptée]
Face à une situation de la vie courante (surcharge d'une école, vieillissement d'un village, exode rural) :
\begin{enumerate}
\item \textbf{Reformuler la situation} avec le vocabulaire du démographe (natalité, fécondité, migration, structure par âge).
\item \textbf{Identifier le problème démographique} : croissance trop rapide, vieillissement, sous-peuplement, surpeuplement relatif.
\item \textbf{Relier le problème à un indicateur mesurable} (TAN élevé, proportion de plus de 60 ans, densité, solde migratoire).
\item \textbf{Proposer une réponse cohérente} avec le type de politique démographique adapté : politique nataliste, politique de limitation des naissances (planning familial), politique migratoire, politique de redistribution spatiale.
\item \textbf{Justifier} la réponse par ses effets attendus et ses limites éventuelles.
\end{enumerate}
\end{methode}

\begin{exoresolu}[Application à une situation concrète : la surcharge scolaire à Douala]
\textbf{Énoncé.} Dans certains quartiers de Douala, les classes de primaire accueillent plus de 100 élèves. La ville grandit de plus de 4\,\% par an, à la fois par accroissement naturel élevé et par afflux de migrants ruraux. Identifiez les deux causes démographiques de cette surcharge et proposez deux réponses adaptées, en les justifiant.
\tcblower
\textbf{Solution détaillée.}

\textbf{Étape 1 — Reformulation.} La population de Douala croît de 4\,\% par an, soit un doublement en $\dfrac{70}{4} \approx 17,5$ ans. La demande de places scolaires augmente au même rythme.

\textbf{Étape 2 — Les deux causes.}
\begin{itemize}
\item \emph{Cause naturelle} : une natalité élevée (fécondité urbaine encore forte, IF $\approx 3,5$) produit chaque année de nombreux enfants en âge d'être scolarisés.
\item \emph{Cause migratoire} : l'exode rural alimente l'arrivée de familles entières venues des régions intérieures (Ouest, Nord-Ouest, Sud-Ouest, Nord) à la recherche d'emplois, ce qui concentre les enfants dans les quartiers périphériques non lotis.
\end{itemize}

\textbf{Étape 3 — Réponses proposées et justifiées.}
\begin{itemize}
\item \emph{Réponse 1 : renforcer le planning familial} (politique de limitation des naissances) pour réduire progressivement la fécondité et donc, à terme, le nombre d'enfants à scolariser. Justification : une baisse de l'indice de fécondité diminue la pression sur les écoles après quelques années ; l'éducation des filles est le levier le plus efficace.
\item \emph{Réponse 2 : construire des écoles dans les quartiers d'accueil ET développer les services (écoles, santé, emplois) dans les zones rurales de départ} (politique d'aménagement et de redistribution spatiale). Justification : offrir des emplois et des services en milieu rural freine l'exode rural à la source et rééquilibre la répartition de la population.
\end{itemize}
\textbf{Conclusion} : la surcharge scolaire de Douala résulte d'une croissance naturelle rapide combinée à un fort solde migratoire positif ; seule une réponse double, à la fois démographique (planification familiale) et spatiale (aménagement du territoire), peut la résoudre durablement.
\end{exoresolu}

\coursec{Atelier méthodologique : Cartographie, graphiques et analyse de documents (TP)}

\courssub{Méthodes transversales pour le Probatoire / Bac Tech}

\begin{methode}[Calculer et interpréter les indicateurs démographiques]
Pour traiter un exercice de démographie à partir de données chiffrées :
\begin{enumerate}
\item \textbf{Identifier l'indicateur demandé} : taux de natalité, taux de mortalité, taux d'accroissement naturel, indice de fécondité, densité, temps de doublement.
\item \textbf{Appliquer la formule exacte} :
\begin{itemize}
\item Taux de natalité : $TN = \dfrac{\text{naissances vivantes}}{\text{population totale}} \times 1000$ (en pour mille, \textperthousand)
\item Taux de mortalité : $TM = \dfrac{\text{décès}}{\text{population totale}} \times 1000$
\item Taux d'accroissement naturel : $TAN = TN - TM$ (en \textperthousand) ; pour l'exprimer en \% : $TAN(\%) = \dfrac{TN - TM}{10}$
\item Densité : $D = \dfrac{\text{population}}{\text{superficie}}$ (hab/km\textsuperscript{2})
\item Temps de doublement : $t = \dfrac{70}{TAN(\%)}$ (approximation utile)
\end{itemize}
\item \textbf{Présenter le calcul complet} : formule littérale, application numérique, résultat avec l'unité.
\item \textbf{Interpréter le résultat} : le comparer à un seuil ou à une moyenne (ex. un TAN supérieur à 2\,\% traduit une croissance rapide ; un indice de fécondité inférieur à 2,1 enfants par femme est sous le seuil de renouvellement des générations).
\item \textbf{Conclure} par une phrase reliant le chiffre à la réalité du territoire étudié.
\end{enumerate}
\end{methode}

\begin{methode}[Construire une réponse argumentée en géographie (méthode du paragraphe rédigé)]
Pour rédiger une réponse ou une conclusion au Probatoire / Bac technique :
\begin{enumerate}
\item \textbf{Annoncer l'idée} (phrase d'introduction claire qui répond directement à la consigne).
\item \textbf{Développer} avec au moins deux arguments, chacun appuyé par un fait précis : un chiffre, une date, un exemple localisé (Cameroun, Afrique, monde).
\item \textbf{Expliquer les mécanismes} : ne pas se contenter d'affirmer, montrer la relation de cause à effet (ex. baisse de la mortalité $\Rightarrow$ accroissement naturel rapide).
\item \textbf{Nuancer} si nécessaire (contrastes entre régions, limites d'une politique).
\item \textbf{Conclure} en revenant à la question posée, sans introduire d'idée nouvelle.
\end{enumerate}
Réflexes : employer le vocabulaire du cours (transition démographique, solde naturel, seuil de renouvellement, dividende démographique), éviter les opinions personnelles non étayées, soigner l'orthographe des termes techniques.
\end{methode}

\begin{methode}[Justifier le choix d'une politique démographique (démarche complète)]
Pour rédiger et justifier une réponse portant sur les politiques démographiques :
\begin{enumerate}
\item \textbf{Diagnostiquer} : établir le problème à partir d'indicateurs (fécondité, TAN, structure par âge, pyramide des âges).
\item \textbf{Choisir le type de politique} : nataliste (encourager les naissances : allocations, congés parentaux, crèches) ou antinataliste / de limitation (planning familial, éducation des filles, contraception, relèvement âge mariage), éventuellement politique migratoire ou d'aménagement.
\item \textbf{Nommer les acteurs} : État, collectivités, ONG, agences internationales (ONU, UNFPA, OMS, Banque mondiale), services de santé.
\item \textbf{Justifier} par les effets attendus sur les indicateurs et par un exemple de bilan réel (ex. politique de l'enfant unique en Chine 1979-2015 : baisse rapide fécondité mais vieillissement accéléré ; planning familial Rwanda/Éthiopie : baisse progressive ; France/Suède : maintien fécondité relativement haute).
\item \textbf{Rédiger la conclusion} en reliant la politique au diagnostic initial.
\end{enumerate}
\end{methode}

\begin{methode}[Lire et exploiter un document statistique ou graphique démographique]
\begin{enumerate}
\item \textbf{Lire le titre, la source, la date et les unités} avant tout calcul.
\item \textbf{Dégager la tendance générale} (hausse, baisse, stagnation) et les ruptures éventuelles (changement de pente).
\item \textbf{Extraire deux ou trois données précises} (valeurs extrêmes, dates de rupture) et les citer avec leurs unités.
\item \textbf{Calculer une évolution} si utile : taux de variation $= \dfrac{V_f - V_i}{V_i} \times 100$ (en \%).
\item \textbf{Interpréter} en reliant les chiffres aux notions du cours (transition démographique, phases de la natalité et de la mortalité, politiques menées).
\end{enumerate}
\end{methode}

\begin{methode}[Réaliser un croquis de géographie démographique (carte ou graphique)]
\begin{enumerate}
\item \textbf{Analyser le sujet} : identifier les phénomènes à représenter (densité, flux, évolution, structure).
\item \textbf{Choisir le fond de carte} (Cameroun, Afrique, Monde) et l'échelle.
\item \textbf{Construire la légende} \emph{avant} de dessiner : classes de valeurs (seuils ronds), symboles (points proportionnels, flèches de flux, hachures/couleurs pour densités), figés (villes, axes).
\item \textbf{Dessiner proprement} : traits fins, coloriage régulier, flèches nettes (épaisseur $\propto$ importance flux).
\item \textbf{Compléter} : titre précis (phénomène + lieu + date), orientation (flèche Nord), échelle graphique, source, légende organisée.
\end{enumerate}
\end{methode}

\courssub{Exercices d'entraînement (non corrigés)}

\begin{exercice}[Calculs et interprétation : Le Niger en 2023]
Données 2023 (estimations) : Population $\approx 27$ millions. Superficie $\SI{1267000}{km^2}$. Naissances $\approx 1\,300\,000$. Décès $\approx 250\,000$.
\begin{enumerate}
\item Calculez le TN, le TM, le TAN (en \textperthousand et en \%), la densité et le temps de doublement.
\item Situez le Niger dans le modèle de la transition démographique. Justifiez.
\item Quels défis ce diagnostic pose-t-il pour la scolarisation et la sécurité alimentaire ?
\end{enumerate}
\end{exercice}

\begin{exercice}[Analyse de pyramide des âges]
On vous donne deux pyramides des âges (non reproduites ici, à imaginer) :
\begin{itemize}
\item \textbf{Pyramide A} : Base très large, sommet très étroit, côtés concaves.
\item \textbf{Pyramide B} : Base étroite, sommet large, côtés convexes (forme « urne » ou « pagode »).
\end{itemize}
\begin{enumerate}
\item Pour chaque pyramide, identifiez la phase de transition démographique correspondante et donnez un exemple de pays/région.
\item Comparez les indices de dépendance (jeunes / vieux) et les besoins en infrastructures (écoles vs hôpitaux/maisons de retraite).
\end{enumerate}
\end{exercice}

\begin{exercice}[Cartographie : Répartition de la population au Cameroun]
À partir des données ci-dessous, réalisez un croquis de la répartition de la population au Cameroun.
\begin{center}
\begin{tabular}{|l|c|c|}\hline
\textbf{Région} & \textbf{Population (2020 est.)} & \textbf{Superficie (km\textsuperscript{2})} \\ \hline
Extrême-Nord & 4\,500\,000 & 34\,263 \\
Nord & 3\,000\,000 & 66\,090 \\
Adamaoua & 1\,200\,000 & 63\,691 \\
Est & 1\,000\,000 & 109\,002 \\
Centre (hors Yaoundé) & 2\,000\,000 & 68\,953 \\
Littoral (hors Douala) & 1\,500\,000 & 20\,248 \\
Ouest & 2\,500\,000 & 13\,892 \\
Nord-Ouest & 2\,200\,000 & 17\,300 \\
Sud-Ouest & 2\,000\,000 & 25\,410 \\
Sud & 800\,000 & 47\,191 \\
Yaoundé (ville) & 4\,000\,000 & 180 \\
Douala (ville) & 4\,000\,000 & 210 \\ \hline
\end{tabular}
\end{center}
\textit{Consignes :} Calculez les densités régionales. Définissez 4 classes de densité. Représentez-les par hachures/couleurs. Placez Douala et Yaoundé par points proportionnels. Tracez les principaux flux d'exode rural (flèches). Titre, légende, Nord, échelle, source.
\end{exercice}

\begin{exercice}[Rédaction argumentée : Sujet type Bac]
\textbf{Sujet :} « Les politiques démographiques mises en œuvre dans le monde depuis 50 ans ont-elles été efficaces ? » Vous répondrez en un paragraphe structuré de 12 à 15 lignes, en vous appuyant sur au moins deux exemples contrastés (un pays du Sud, un pays du Nord) et en nuancant le bilan.
\end{exercice}

\courssub{Pièges à éviter (Rappel pour l'examen)}

\begin{attention}[Les 5 erreurs qui coûtent des points au Bac]
\begin{enumerate}
\item \textbf{Confondre \textperthousand\ et \%} : les taux de natalité et de mortalité s'expriment en pour mille (\textperthousand) ; oublier de diviser par 10 pour obtenir le TAN en \% fausse toute l'interprétation (27\,\textperthousand\ $=$ 2,7\,\%, pas 27\,\% !).
\item \textbf{Oublier l'unité dans le résultat} : une densité sans hab/km\textsuperscript{2}, un taux sans \textperthousand\ ou \% est une réponse incomplète, même si le calcul est juste.
\item \textbf{Citer un chiffre sans l'interpréter} : écrire \og la fécondité est de 4,6 \fg{} sans la comparer au seuil de renouvellement (2,1) ne rapporte pas tous les points ; chaque donnée doit être mise en relation avec une notion du cours.
\item \textbf{Mélanger les types de politiques} : attribuer le planning familial à une politique nataliste (ou l'inverse) est une erreur de notion rédhibitoire ; vérifier la cohérence entre le diagnostic (croissance rapide ou déclin) et la politique proposée.
\item \textbf{Se tromper dans le taux de variation} : diviser par la valeur finale au lieu de la valeur initiale, ou oublier de multiplier par 100 ; toujours écrire la formule $\dfrac{V_f - V_i}{V_i} \times 100$ avant l'application numérique.
\end{enumerate}
\end{attention}

\begin{corrige}[Exercice n°1 : Le Niger en 2023]
\textbf{1. Calculs.}
$TN = \frac{1\,300\,000}{27\,000\,000} \times 1000 \approx 48,1\,\textperthousand$.
$TM = \frac{250\,000}{27\,000\,000} \times 1000 \approx 9,3\,\textperthousand$.
$TAN = 48,1 - 9,3 = 38,8\,\textperthousand = 3,88\,\%$.
$D = \frac{27\,000\,000}{1\,267\,000} \approx 21,3\ \text{hab/km}^2$.
$t = \frac{70}{3,88} \approx 18\ \text{ans}$.

\textbf{2. Transition.} Phase 2 (début de transition) : natalité très forte, mortalité en baisse forte mais encore élevée, croissance très rapide (TAN $\approx 3,9\,\%$). Exemple type : Niger, Tchad, Somalie, Mali.

\textbf{3. Défis.} Scolarisation : il faut construire l'équivalent du parc scolaire actuel tous les ~18 ans. Sécurité alimentaire : production agricole doit croître $> 3,9\,\%$/an pour suivre la demande, or croissance rendements $\approx 1-2\,\%$. Risque de dépendance aide alimentaire, conflits fonciers, migration.
\end{corrige}

\begin{corrige}[Exercice n°2 : Analyse de pyramides]
\textbf{Pyramide A (base large, sommet étroit, concave)} : Phase 2 (transition début). Natalité forte, mortalité encore assez forte (côtés concaves = mortalité infantile/juvénile), espérance de vie faible. Ex. : Niger, Nigeria, RDC, Angola. Dépendance jeunes très forte ($> 80\,\%$). Besoins : écoles, pédiatrie, vaccins, emploi futur.

\textbf{Pyramide B (base étroite, sommet large, convexe)} : Phase 4/5 (post-transition). Natalité très faible (base étroite), mortalité très faible (sommet large = longévité), espérance de vie élevée. Ex. : Japon, Italie, Allemagne, Corée du Sud. Dépendance vieux forte ($> 30\,\%$). Besoins : gériatrie, Ehpad, financement retraites, adaptation logement/transport, immigration de travail.
\end{corrige}

\newpage

\coursec{Exercices}

\courssub{Je vérifie mes connaissances}

\begin{exercice}[QCM : Indicateurs et transition démographique]
\textbf{Pour chaque question, une seule réponse est exacte. Recopiez le numéro et la lettre choisie.}
\begin{enumerate}
    \item L'Indice Conjoncturel de Fécondité (ICF) représente :
    \begin{itemize}
        \item[A.] Le nombre moyen d'enfants qu'une femme aurait au cours de sa vie féconde si elle était soumise aux taux de fécondité par âge de l'année considérée.
        \item[B.] Le nombre de naissances vivantes pour 1000 habitants en une année.
        \item[C.] La différence entre le taux de natalité et le taux de mortalité.
        \item[D.] Le nombre d'enfants nés vivants par femme de 15 à 49 ans recensée.
    \end{itemize}
    \item Dans le modèle de la transition démographique, la phase 2 (transition précoce) se caractérise par :
    \begin{itemize}
        \item[A.] Natalité forte, mortalité forte, croissance nulle.
        \item[B.] Natalité forte, mortalité en forte baisse, croissance forte.
        \item[C.] Natalité en baisse, mortalité faible, croissance modérée.
        \item[D.] Natalité faible, mortalité faible, croissance nulle ou négative.
    \end{itemize}
    \item Le « dividende démographique » désigne :
    \begin{itemize}
        \item[A.] L'augmentation rapide de la population totale.
        \item[B.] La part des personnes âgées dans la population.
        \item[C.] Le gain économique potentiel lié à la baisse de la charge de dépendance (ratio population en âge de travailler / population inactive) lors de la transition.
        \item[D.] Les transferts financiers des migrants vers leur pays d'origine.
    \end{itemize}
    \item La politique de l'enfant unique en Chine (1979-2015) est un exemple type de politique :
    \begin{itemize}
        \item[A.] Pro-nataliste incitative.
        \item[B.] Anti-nataliste coercitive.
        \item[C.] Migratoire de peuplement.
        \item[D.] De planning familial non directif.
    \end{itemize}
\end{enumerate}
\end{exercice}

\begin{exercice}[Vrai ou Faux ? Justifiez brièvement chaque réponse par une phrase ou un chiffre clé.]
\begin{enumerate}
    \item La population mondiale a dépassé les 8 milliards d'habitants en novembre 2022.
    \item Le Cameroun se trouve actuellement dans la phase 1 de la transition démographique (natalité forte, mortalité forte).
    \item Une pyramide des âges à base large et sommet étroit indique une population vieillissante et une croissance démographique faible.
    \item Le Taux d'Accroissement Naturel (TAN) se calcule par la formule : TAN = (Taux de natalité -- Taux de mortalité) / 10.
\end{enumerate}
\end{exercice}

\begin{exercice}[Définitions attendues au Probatoire]
\textbf{Donnez une définition précise (2 à 3 lignes) pour chacun des termes suivants :}
\begin{enumerate}
    \item Transition démographique
    \item Indice de jeunesse (ou part des moins de 15 ans)
    \item Politique pro-nataliste
    \item Urbanisation primate (ou macrocéphalie urbaine)
\end{enumerate}
\end{exercice}

\begin{exercice}[Calculs directs et interprétation]
\textbf{Données 2023 (estimations) :}
\begin{center}
\begin{tabularx}{\linewidth}{|l|X|X|X|}\hline
\textbf{Indicateur} & \textbf{Cameroun} & \textbf{France} & \textbf{Niger} \\ \hline
Population (millions) & 28,6 & 68,1 & 26,2 \\ \hline
Taux de natalité (‰) & 34,5 & 10,2 & 46,8 \\ \hline
Taux de mortalité (‰) & 8,4 & 9,8 & 10,5 \\ \hline
ICF (enfants/femme) & 4,5 & 1,7 & 6,7 \\ \hline
Espérance de vie à la naissance (ans) & 60,3 & 82,5 & 62,9 \\ \hline
Part des 0-14 ans (\%) & 41,5 & 17,2 & 49,8 \\ \hline
\end{tabularx}
\end{center}
\begin{enumerate}
    \item Calculez le Taux d'Accroissement Naturel (TAN en \%) pour les trois pays.
    \item Calculez le doublement approximatif de la population (règle du 70) pour le Cameroun et le Niger.
    \item Comparez l'Indice de Jeunesse du Cameroun et de la France. Quel défi majeur cela pose-t-il pour le système éducatif camerounais ?
    \item La France a un taux de mortalité (9,8‰) supérieur à celui du Cameroun (8,4‰). Comment expliquer ce paradoxe apparent ?
\end{enumerate}
\end{exercice}

\courssub{Je m'entraîne}

\begin{exercice}[Analyse comparative de pyramides des âges : Niger vs Japon]
\textbf{Documents :} Deux pyramides des âges (schématiques) sont fournies ci-dessous.
\begin{popfigure}
\begin{tikzpicture}[scale=0.65]
% Niger Pyramid (Left)
\def\maxpop{12}
\foreach \y/\male/\female in {
    0/5.8/6.0, 5/5.2/5.4, 10/4.5/4.7, 15/3.8/4.0, 20/3.0/3.3,
    25/2.4/2.7, 30/1.9/2.1, 35/1.5/1.7, 40/1.1/1.3, 45/0.8/1.0,
    50/0.6/0.8, 55/0.4/0.6, 60/0.3/0.4, 65/0.2/0.3, 70/0.1/0.2, 75/0.05/0.1} {
    \fill[popblue!60] (-\male,\y) rectangle (0,\y+5);
    \fill[poppink!60] (0,\y) rectangle (\female,\y+5);
}
\draw[<->] (-\maxpop,0) -- (\maxpop,0) node[midway,below] {Population (millions)};
\foreach \y in {0,20,40,60,80} \draw (-\maxpop,\y) -- (\maxpop,\y) node[right] {\y ans};
\node[align=center,font=\bfseries] at (0,-1) {Niger (2023)\\Base large, sommet pointu};
\end{tikzpicture}
\quad
\begin{tikzpicture}[scale=0.65]
% Japan Pyramid (Right)
\def\maxpop{6}
\foreach \y/\male/\female in {
    0/2.5/2.4, 5/2.8/2.7, 10/3.0/2.9, 15/3.2/3.1, 20/3.5/3.4,
    25/3.8/3.7, 30/4.0/3.9, 35/4.2/4.1, 40/4.5/4.4, 45/4.8/4.7,
    50/5.0/5.0, 55/4.8/4.9, 60/4.5/4.7, 65/4.0/4.3, 70/3.5/3.9, 75/2.5/3.2, 80/1.5/2.2, 85/0.8/1.5, 90/0.3/0.8, 95/0.1/0.3} {
    \fill[popblue!60] (-\male,\y) rectangle (0,\y+5);
    \fill[poppink!60] (0,\y) rectangle (\female,\y+5);
}
\draw[<->] (-\maxpop,0) -- (\maxpop,0) node[midway,below] {Population (millions)};
\foreach \y in {0,20,40,60,80,100} \draw (-\maxpop,\y) -- (\maxpop,\y) node[right] {\y ans};
\node[align=center,font=\bfseries] at (0,-1) {Japon (2023)\\Base étroite, sommet large};
\end{tikzpicture}
\legende{Pyramides des âges schématiques (largeur proportionnelle aux effectifs par classe d'âge 5 ans). Hommes à gauche (bleu), Femmes à droite (rose).}
\end{popfigure}
\begin{enumerate}
    \item Identifiez la phase de transition démographique correspondante pour chaque pays. Justifiez par deux caractéristiques visuelles de la pyramide.
    \item Calculez approximativement l'indice de jeunesse (part des 0-14 ans) pour le Niger en lisant la base de la pyramide (somme des 3 premières barres / total). Même chose pour le Japon.
    \item Quels sont les deux défis majeurs pour le système de santé au Niger et pour le système de retraite au Japon ?
    \item Représentez schématiquement l'évolution prévisible de la pyramide du Niger d'ici 2050 si la fécondité baisse rapidement (transition accélérée).
\end{enumerate}
\end{exercice}

\begin{exercice}[Lecture de tableaux statistiques : Cameroun vs France (2023)]
\textbf{À partir du tableau de l'exercice 4 (partie « Je vérifie mes connaissances »), répondez aux questions suivantes :}
\begin{enumerate}
    \item Classez les trois pays (Cameroun, France, Niger) par ordre croissant de : a) ICF, b) Espérance de vie, c) TAN.
    \item Le Cameroun a un ICF de 4,5 et la France de 1,7. Exprimez en pourcentage de combien la fécondité camerounaise est supérieure à la française.
    \item En vous basant sur la part des 0-14 ans et l'espérance de vie, calculez le ratio de dépendance approximatif (jeunes + vieux / 15-64 ans) pour le Cameroun et la France. (Hypothèse : Part 65+ = 100 - Part 0-14 - Part 15-64 ; on approxime Part 15-64 = 100 - Part 0-14 - Part 65+ ; pour simplifier : Part 65+ Cameroun $\approx$ 3,5\%, France $\approx$ 21\%).
    \item Rédigez un paragraphe de synthèse (8 lignes max) contrastant les « problèmes démographiques » de ces deux pays (explosion scolaire vs vieillissement/pensions).
\end{enumerate}
\end{exercice}

\begin{exercice}[Typologie des politiques démographiques : Études de cas]
\textbf{Associez chaque description de politique à son type (Pro-nataliste / Anti-nataliste / Migratoire / Planning familial) et citez un pays exemple (autre que la Chine pour l'anti-nataliste).}
\begin{enumerate}
    \item « Allocation de naissance de 923€, congé parental de 3 ans rémunéré, quotient familial avantageux, crèches subventionnées. » (France, 2020s)
    \item « Stérilisation forcée de millions d'hommes (vasectomies) dans les années 1970, slogan "Deux enfants c'est assez". » (Inde, 1975-77)
    \item « Distribution gratuite de contraceptifs, éducation sexuelle à l'école, avortement légal, sans contrainte ni incitation financière forte. » (Tunisie, depuis 1960s / Iran, 1990s)
    \item « Programme "Transmigration" : déplacement de millions de personnes de Java/Bali vers Sumatra/Kalimantan/Papouasie pour désengorger les îles denses. » (Indonésie, 1970s-2000s)
    \item « Politique de l'enfant unique (1979-2015) : amendes, pression sociale, avortements sélectifs, déséquilibre sexuels à la naissance (115 garçons pour 100 filles). »
\end{enumerate}
\textbf{Question de synthèse :} Pourquoi les politiques pro-natalistes des pays développés (France, Suède, Japon, Singapour) peinent-elles à remonter durablement l'ICF au-dessus du seuil de remplacement (2,1) ?
\end{exercice}

\begin{exercice}[Croquis commenté : Répartition et dynamique de la population au Cameroun]
\textbf{Consigne :} Sur le fond de carte simplifié ci-dessous (fourni à l'examen ou à tracer), réalisez un croquis légendé et titré montrant :
\begin{itemize}
    \item Les 5 grandes zones de densité (Très forte $>$ 100 hab/km² ; Forte 50-100 ; Moyenne 20-50 ; Faible 5-20 ; Très faible $<$ 5).
    \item Les 5 principales villes (Yaoundé, Douala, Garoua, Bamenda, Maroua) avec un symbole proportionnel à la population.
    \item Les deux grands flux de migration interne (Nord $\rightarrow$ Sud ; Ouest $\rightarrow$ Littoral/Centre).
    \item Les zones de pression foncière / conflits agro-pastoraux (Adamaoua, Nord-Ouest, Extrême-Nord).
\end{itemize}
\begin{popfigure}
\begin{tikzpicture}[scale=1.2]
% Simplified Cameroon Outline
\draw[popdark, thick] (0,0) -- (1,0.5) -- (2,1.5) -- (3,2) -- (4,1.8) -- (5,1) -- (5.5,0) -- (5,-0.5) -- (4.5,-1.5) -- (3.5,-2) -- (2.5,-1.8) -- (1.5,-1) -- (0.5,-0.2) -- cycle;
% Regions approximate centroids for labels
\node at (2.5,0.5) {Cameroun};
% Grid for student reference
\draw[step=1cm, popmuted, very thin] (0,-2) grid (5,2);
\end{tikzpicture}
\legende{Fond de carte simplifié du Cameroun (grille 1cm $\approx$ 100 km). À compléter par l'élève.}
\end{popfigure}
\textbf{Attendus :} Légende organisée (Densités / Villes / Flux / Pressions), Titre précis (ex: « Répartition et dynamique de la population au Cameroun, 2023 »), Flèche Nord, Échelle, Analyse commentée (10 lignes) sous le croquis.
\end{exercice}

\courssub{Je me prépare au Bac}

\begin{exercice}[Sujet type Probatoire : Les problèmes démographiques mondiaux sont contrastés]
\textbf{Durée : 1h30 -- Coefficient 3}
\textbf{Documents fournis (à imaginer / fournis le jour de l'examen) :}
\begin{itemize}
    \item \textbf{Doc 1 :} Carte mondiale du taux d'accroissement naturel (2023) : Afrique $>$ 2,5\% ; Europe/Amérique du Nord $<$ 0,3\% ; Asie/LatAm intermédiaires.
    \item \textbf{Doc 2 :} Deux pyramides des âges : Niger (base très large, ICF 6,7) et Japon (base étroite, sommet large, ICF 1,3).
    \item \textbf{Doc 3 :} Tableau comparatif d'indicateurs (PIB/hab, Taux alphabétisation femmes, ICF, Part 0-14 ans, Part 65+ ans) pour Niger, Cameroun, Brésil, France, Japon.
    \item \textbf{Doc 4 :} Extrait rapport ONU 2024 : « La croissance démographique se concentre dans les pays les plus pauvres... Le vieillissement s'accélère dans les pays riches... Les politiques migratoires deviennent un enjeu géopolitique majeur. »
\end{itemize}
\textbf{Travail demandé :}
\begin{enumerate}
    \item \textbf{Analyse des documents (6 pts) :}
    \begin{enumerate}
        \item À l'aide du Doc 1, citez deux régions du monde à forte croissance et deux à croissance faible/négative.
        \item Comparez la structure par âge du Niger et du Japon (Doc 2). En déduire le type de problèmes démographiques propres à chacun (mots-clés : « explosion scolaire », « vieillissement », « charge de dépendance »).
        \item Montrez, grâce au Doc 3, le lien entre niveau de développement (PIB/hab, scolarisation femmes) et niveau de fécondité (ICF).
    \end{enumerate}
    \item \textbf{Mise en relation et connaissances personnelles (9 pts) :}
    \begin{enumerate}
        \item En vous appuyant sur la transition démographique, expliquez pourquoi les problèmes démographiques sont « contrastés » à l'échelle mondiale (phases 2/3 vs phases 4/5).
        \item Présentez deux types de politiques démographiques mises en œuvre pour répondre à ces problèmes contrastés (un exemple anti-nataliste, un exemple pro-nataliste). Évaluez brièvement leur efficacité.
        \item Discutez l'affirmation du Doc 4 : « La croissance démographique se concentre dans les pays les plus pauvres ». Est-ce un frein ou une opportunité (dividende démographique) pour le développement ? Argumentez.
    \end{enumerate}
    \item \textbf{Synthèse cartographique (5 pts) :} Réalisez un croquis mondial simplifié (planisphère) intitulé « Les grands problèmes démographiques mondiaux (2024) » avec une légende organisée en 3 items : 1. Croissance forte / Jeunesse (hachures), 2. Vieillissement / Déclin (pointillés), 3. Flux migratoires majeurs (flèches).
\end{enumerate}
\end{exercice}

\begin{exercice}[Sujet type Baccalauréat Technique : Le Cameroun et le dividende démographique 2035]
\textbf{Durée : 2h -- Coefficient 4}
\textbf{Sujet :} « Le Cameroun peut-il capter son dividende démographique d'ici 2035 ? Appuyez votre démonstration sur l'analyse de la structure par âge, les défis éducation/emploi/santé, les disparités régionales et les politiques publiques actuelles. »
\textbf{Données clés à mobiliser (connaissances du cours) :}
\begin{itemize}
    \item Population 2023 : $\approx$ 28,6 millions ; ICF : 4,5 (baisse lente) ; TAN : 2,6\% ; Part 0-14 ans : 41,5\% ; Part 15-64 ans : 55\% ; Part 65+ : 3,5\%.
    \item Objectif PND30 / Vision 2035 : Pays à revenu intermédiaire, émergent.
    \item Indicateurs sociaux : Taux de scolarisation primaire $\approx$ 95\% (mais achèvement faible), Secondaire $\approx$ 50\% ; Chômage jeunes (15-24 ans) $\approx$ 13\% (sous-emploi massif 70\%) ; Dépense santé $\approx$ 4\% PIB ; Mortalité maternelle 438/100 000 NV.
    \item Disparités : ICF Nord/Extrême-Nord $>$ 6,0 vs Sud/Centre $\approx$ 3,5 ; Urbanisation 58\% (Douala/Yaoundé $\approx$ 6-7 millions) ; Pressions foncières (Adamaoua, NO, SO).
    \item Politiques : Politique de population 1991 (révisée 2007), Plan National de Développement 2020-2030 (PND30), Stratégie Nationale de Développement (SND30), gratuité enseignement primaire, assurance maladie universelle (en cours).
\end{itemize}
\textbf{Plan de dissertation imposé :}
\begin{enumerate}
    \item \textbf{Introduction :} Accroche (chiffre clé), Définition du dividende démographique, Définition du contexte Cameroun 2035, Annonce du plan en 3 parties.
    \item \textbf{Partie 1 : Un potentiel démographique indéniable mais une fenêtre d'opportunité étroite.}
    \begin{itemize}
        \item Structure par âge actuelle : « jeunesse » massive (base large), charge de dépendance jeune élevée (75\%).
        \item Projection 2035 : Entrée massive sur le marché du travail (300 000+ jeunes/an), baisse espérée de l'ICF vers 3,5-4,0.
        \item Condition \textit{sine qua non} : Baisse accélérée de la fécondité (éducation filles, planning familial).
    \end{itemize}
    \item \textbf{Partie 2 : Les goulots d'étranglement structurels : Éducation, Emploi, Santé, Disparités.}
    \begin{itemize}
        \item Qualité de l'éducation (surpopulation classes, inadéquation formation/emploi).
        \item Marché du travail : prédominance secteur informel, création d'emplois formels insuffisante ($<$ 50 000/an).
        \item Santé reproductive : mortalité maternelle/infantile encore haute, accès contraception inégal (rural/urbain, Nord/Sud).
        \item Disparités régionales : « Deux Camerouns » (Nord jeune/pauvre/fécond vs Sud plus âgé/urbain/éduqué).
    \end{itemize}
    \item \textbf{Partie 3 : Les réponses publiques : ambitions vs mise en œuvre.}
    \begin{itemize}
        \item Cadre stratégique : Vision 2035, PND30, Politique de population 2007.
        \item Mesures concrètes : Gratuité école, formation professionnelle (MINFOPRA), Projet « Emploi Jeunes », Couverture Santé Universelle (CSU).
        \item Limites : Budgétisation insuffisante (secteur social $<$ 20\% budget), gouvernance, corruption, insécurité (NO/SO, Extrême-Nord) freinant l'investissement.
    \end{itemize}
    \item \textbf{Conclusion :} Bilan nuancé (fenêtre ouverte jusqu'à 2040-2050 mais risque de « malédiction démographique » si investissements humano-productifs ratés), Ouverture (rôle de la diaspora, migration, technologie).
\end{enumerate}
\end{exercice}

\begin{exercice}[Cas pratique / Note de synthèse : Conseiller au Ministère de l'Économie]
\textbf{Contexte :} Le Ministre de l'Économie, de la Planification et de l'Aménagement du Territoire (MINEPAT) vous demande une note d'argumentation (2 pages max, 600-800 mots) sur la proposition suivante : \textbf{« Instaurer une allocation familiale universelle majorée au 3e enfant (ex: 10 000 FCFA/mois/enfant 1-2 ; 25 000 FCFA/mois à partir du 3e) pour relancer la natalité et soutenir les familles. »}
\textbf{Données d'appui :}
\begin{itemize}
    \item Budget État 2024 : $\approx$ 6 000 milliards FCFA. Dépenses sociales (éducation+santé+protection sociale) $\approx$ 1 800 milliards.
    \item Nombre d'enfants 0-14 ans : $\approx$ 12 millions. Allocataires potentiels (ménages déclarés) : estimation 2 millions de foyers.
    \item Coût estimé mesure (scénario majoration 3e+) : $\approx$ 150-200 milliards FCFA/an (0,3\% PIB).
    \item Élasticité de la fécondité aux allocations (littérature internationale) : +0,05 à +0,15 enfant par femme pour 1\% PIB dépensé (effet faible, surtout sur le calendrier pas le quanta final).
    \item Contexte camerounais : ICF 4,5 (encore au-dessus remplacement), forte fécondité non désirée (besoin non satisfait planning familial $\approx$ 18\%), pauvreté monétaire 37\%, inégalités filles/garçons scolarisation secondaire.
\end{itemize}
\textbf{Consignes de rédaction :}
\begin{enumerate}
    \item Structurez la note : Objet, Contexte, Analyse coûts-bénéfices (budgétaire, démique, social), Analyse d'équité (ciblage vs universalisme, genre), Recommandation finale (Pour / Contre / Amendement) argumentée.
    \item Utilisez les concepts du cours : Dividende démographique, Transition, Charge de dépendance, Planning familial, Équité horizontale/verticale.
    \item Ton administratif : objectif, chiffré, synthétique.
\end{enumerate}
\end{exercice}

\newpage

\coursec{Corrigés des exercices}

\coursec{La population mondiale : problèmes démographiques et politiques}
\courssub{Exercices corrigés et analyses types — Première Technique}

\begin{corrige}[Exercice 1 — QCM : Indicateurs et transition démographique]
\begin{enumerate}
    \item \textbf{Réponse A.} L'ICF est un indicateur \cle{conjoncturel} : il mesure le nombre moyen d'enfants qu'aurait une femme soumise, toute sa vie féconde, aux taux de fécondité par âge observés l'année considérée. B décrit le taux de natalité ; C décrit l'accroissement naturel ; D confond effectif recensé et indice synthétique.
    \item \textbf{Réponse B.} La phase 2 (transition précoce) : la mortalité chute fortement (progrès médicaux, hygiène) tandis que la natalité reste élevée, d'où une \cle{croissance démographique forte}. A = phase 1 ; C = phase 3 ; D = phases 4/5.
    \item \textbf{Réponse C.} Le \cle{dividende démographique} est le gain de croissance potentiel qui apparaît quand la part de la population en âge de travailler augmente relativement aux inactifs (jeunes et vieux), grâce à la baisse de la fécondité. Il n'est pas automatique : il exige emplois, éducation et santé.
    \item \textbf{Réponse B.} La politique de l'enfant unique (amendes, stérilisations et avortements contraints) est une politique \cle{anti-nataliste coercitive}, visant à réduire la natalité par la contrainte, contrairement au planning familial non directif.
\end{enumerate}
\end{corrige}

\begin{attention}[Points de vigilance — Exercice 1]
Ne pas confondre taux de natalité (pour 1000 habitants) et ICF (enfants par femme). Au Probatoire, la définition de l'ICF doit mentionner le caractère hypothétique (« si elle était soumise aux taux de l'année »).
\end{attention}

\begin{corrige}[Exercice 2 — Vrai ou Faux ?]
\begin{enumerate}
    \item \textbf{Vrai.} L'ONU a officialisé le franchissement des 8 milliards d'habitants le 15 novembre 2022 (contre 1 milliard vers 1804 et 2 milliards en 1927).
    \item \textbf{Faux.} Le Cameroun est en \cle{phase 2/3} (transition engagée) : natalité encore forte (34,5\,\textperthousand) mais mortalité déjà fortement abaissée (8,4\,\textperthousand), ICF en baisse lente (4,5). La phase 1 (natalité et mortalité fortes) correspond à l'époque précoloniale.
    \item \textbf{Faux.} C'est l'inverse : une base large et un sommet étroit signalent une population \cle{jeune} à forte natalité et forte croissance (ex. Niger). Une population vieillissante a une base étroite et un sommet élargi (ex. Japon).
    \item \textbf{Vrai.} Les taux étant exprimés en \textperthousand, on obtient le TAN en \% par : TAN $= \dfrac{\text{natalité} - \text{mortalité}}{10}$. Exemple Cameroun : $(34,5 - 8,4)/10 = 2,61\,\%$.
\end{enumerate}
\end{corrige}

\begin{attention}[Piège fréquent — Exercice 2]
Oublier de diviser par 10 pour passer des \textperthousand aux \% est l'erreur de calcul la plus fréquente au Probatoire. Toujours vérifier l'ordre de grandeur : un TAN mondial moyen est d'environ 1\,\%.
\end{attention}

\begin{corrige}[Exercice 3 — Définitions attendues au Probatoire]
\begin{enumerate}
    \item \textbf{Transition démographique :} passage historique d'un régime démographique traditionnel (natalité et mortalité élevées, croissance faible) à un régime moderne (natalité et mortalité faibles, croissance faible ou négative), en plusieurs phases, la baisse de la mortalité précédant celle de la natalité.
    \item \textbf{Indice de jeunesse :} part (en \%) des moins de 15 ans (ou parfois moins de 20 ans) dans la population totale. Il mesure le poids de la jeunesse ; au-dessus de 40\,\%, on parle de population très jeune (Cameroun : 41,5\,\%).
    \item \textbf{Politique pro-nataliste :} ensemble de mesures publiques (allocations familiales, congés parentaux rémunérés, crèches, avantages fiscaux) visant à encourager les naissances pour relever la fécondité et ralentir le vieillissement (ex. France, Suède).
    \item \textbf{Urbanisation primate (macrocéphalie urbaine) :} situation où une ou deux villes concentrent une part disproportionnée de la population urbaine et des fonctions économiques du pays, au détriment des autres villes (ex. Douala et Yaoundé au Cameroun).
\end{enumerate}
\end{corrige}

\begin{corrige}[Exercice 4 — Calculs directs et interprétation]
\begin{enumerate}
    \item \textbf{TAN} $=$ (natalité $-$ mortalité)$/10$ :
    \begin{itemize}
        \item Cameroun : $(34,5 - 8,4)/10 = 26,1/10 = \mathbf{2,61\,\%}$
        \item France : $(10,2 - 9,8)/10 = 0,4/10 = \mathbf{0,04\,\%}$
        \item Niger : $(46,8 - 10,5)/10 = 36,3/10 = \mathbf{3,63\,\%}$
    \end{itemize}
    \item \textbf{Règle du 70} (temps de doublement $\approx 70/\text{TAN en \%}$) :
    \begin{itemize}
        \item Cameroun : $70/2,61 \approx \mathbf{27}$ \textbf{ans} (doublement vers 2050).
        \item Niger : $70/3,63 \approx \mathbf{19}$ \textbf{ans} (doublement vers 2042).
    \end{itemize}
    \item Indice de jeunesse : Cameroun 41,5\,\% contre France 17,2\,\%, soit environ \textbf{2,4 fois plus} ($41,5/17,2 \approx 2,41$). Défi majeur : il faut scolariser chaque année des cohortes croissantes — construction de salles de classe, recrutement d'enseignants, manuels — d'où le risque de classes surchargées et de baisse de qualité (\cle{explosion scolaire}).
    \item \textbf{Paradoxe de la mortalité :} le taux brut de mortalité dépend de la \cle{structure par âge}. La France compte beaucoup de personnes âgées (environ 21\,\% de 65 ans et plus), âges où les décès sont nombreux ; le Cameroun a une population très jeune, donc mécaniquement moins de décès rapportés à la population totale, malgré une espérance de vie bien plus faible (60,3 ans contre 82,5 ans). Pour comparer les conditions sanitaires, on utilise l'espérance de vie ou la mortalité infantile, pas le taux brut.
\end{enumerate}
\end{corrige}

\begin{attention}[Point de vigilance — Exercice 4]
Question 4 : citer explicitement le rôle de la structure par âge est la justification attendue ; une réponse du type « les Français meurent plus » est fausse et sanctionnée.
\end{attention}

\begin{corrige}[Exercice 5 — Analyse comparative de pyramides : Niger vs Japon]
\begin{enumerate}
    \item \textbf{Niger : phase 2} (transition précoce) — base très large (forte natalité, ICF 6,7) et rétrécissement brutal des barres avec l'âge (mortalité encore marquée, sommet pointu). \textbf{Japon : phase 4/5} (transition achevée, post-transition) — base étroite et rétrécie (natalité faible, ICF 1,3), classes d'âge moyennes et âgées renflées (vieillissement, sommet large).
    \item \textbf{Niger} (somme des 3 premières classes, hommes + femmes) : $(5,8+6,0)+(5,2+5,4)+(4,5+4,7) = 11,8+10,6+9,2 = 31,6$ millions. Total approximatif : somme de toutes les barres $\approx 62$ millions d'équivalents-barres ; en réalité le Niger compte $\approx 26$ millions d'habitants, mais en \emph{proportions lues sur la pyramide} : $31,6/62 \approx \mathbf{51\,\%}$, cohérent avec les 49,8\,\% officiels. \textbf{Japon} : 3 premières classes $= (2,5+2,4)+(2,8+2,7)+(3,0+2,9) = 4,9+5,5+5,9 = 16,3$ ; total $\approx 126$ équivalents ; $16,3/126 \approx \mathbf{13\,\%}$, cohérent avec la réalité ($\approx 12\,\%$).
    \item \textbf{Niger :} (1) pédiatrie et maternité débordées (mortalité maternelle et infantile) ; (2) besoins massifs en vaccination, nutrition et centres de santé de proximité. \textbf{Japon :} (1) financement des retraites avec un ratio actifs/retraités qui s'effondre ; (2) explosion des dépenses de santé et de dépendance liées au grand âge.
    \item \textbf{Pyramide du Niger en 2050 (transition accélérée) :} base resserrée (fécondité en baisse), renflement des classes 15-49 ans (cohortes nombreuses nées avant la baisse), sommet encore étroit mais s'épaississant : forme en « fuseau » ou « silo », typique de la phase 3 — c'est la configuration du \cle{dividende démographique}.
\end{enumerate}
\end{corrige}

\begin{corrige}[Exercice 6 — Lecture de tableaux statistiques]
\begin{enumerate}
    \item Classements croissants :
    \begin{itemize}
        \item a) ICF : France (1,7) $<$ Cameroun (4,5) $<$ Niger (6,7).
        \item b) Espérance de vie : Cameroun (60,3) $<$ Niger (62,9) $<$ France (82,5).
        \item c) TAN : France (0,04\,\%) $<$ Cameroun (2,61\,\%) $<$ Niger (3,63\,\%).
    \end{itemize}
    \item Écart de fécondité : $\dfrac{4,5 - 1,7}{1,7} \times 100 = \dfrac{2,8}{1,7}\times 100 \approx \mathbf{164,7\,\%}$. La fécondité camerounaise est supérieure d'environ 165\,\% à la française (soit $\approx 2,6$ fois plus élevée).
    \item \textbf{Ratio de dépendance} $= \dfrac{\text{0--14} + \text{65+}}{\text{15--64}} \times 100$ :
    \begin{itemize}
        \item Cameroun : part 15--64 $= 100 - 41,5 - 3,5 = 55\,\%$ ; ratio $= \dfrac{41,5+3,5}{55}\times 100 = \dfrac{45}{55}\times 100 \approx \mathbf{82\,\%}$ (82 inactifs pour 100 actifs).
        \item France : part 15--64 $= 100 - 17,2 - 21 = 61,8\,\%$ ; ratio $= \dfrac{17,2+21}{61,8}\times 100 = \dfrac{38,2}{61,8}\times 100 \approx \mathbf{62\,\%}$.
    \end{itemize}
    \item \textbf{Synthèse attendue (corrigé type) :} Le Cameroun et la France illustrent deux problèmes démographiques opposés. Le Cameroun, en pleine transition (ICF 4,5, TAN 2,6\,\%), fait face à une explosion de la demande sociale : 41,5\,\% de moins de 15 ans exigent écoles, enseignants, puis emplois pour les 300 000 jeunes entrant chaque année sur le marché du travail ; la charge de dépendance (82\,\%) pèse sur les actifs et freine l'épargne. La France, transition achevée (ICF 1,7 sous le seuil de remplacement), connaît un quasi-arrêt de la croissance naturelle (TAN 0,04\,\%) et un vieillissement massif (21\,\% de 65 ans et plus) : financement des retraites, dépenses de santé et pénurie de main-d'œuvre deviennent critiques. L'un doit investir dans la jeunesse, l'autre dans la prise en charge des aînés : deux diagnostics contrastés appelant des politiques démographiques inverses.
\end{enumerate}
\end{corrige}

\begin{corrige}[Exercice 7 — Typologie des politiques démographiques]
\begin{enumerate}
    \item \textbf{Pro-nataliste incitative} (France) : aides financières et services sans contrainte, pour encourager les naissances.
    \item \textbf{Anti-nataliste coercitive} (Inde, 1975-1977) : stérilisations forcées sous l'État d'urgence — exemple type de coercition, qui a d'ailleurs provoqué un rejet durable du planning en Inde.
    \item \textbf{Planning familial non directif} (Tunisie, Iran des années 1990) : mise à disposition libre des moyens contraceptifs et éducation, sans obligation ni forte incitation ; l'Iran a ainsi divisé son ICF par deux en une décennie.
    \item \textbf{Politique migratoire de peuplement} (Indonésie, « Transmigration ») : redistribution spatiale de la population des îles surpeuplées vers les îles périphériques.
    \item \textbf{Anti-nataliste coercitive} (Chine, enfant unique) : amendes, pressions, avortements sélectifs ; bilan démographique « efficace » mais lourds effets pervers (déséquilibre des sexes, vieillissement accéléré, « 4-2-1 »).
\end{enumerate}
\textbf{Question de synthèse :} Les politiques pro-natalistes peinent à ramener l'ICF au-dessus de 2,1 car les causes de la baisse de fécondité sont structurelles : allongement des études et entrée tardive dans la vie adulte, coût du logement et de l'enfant, massification du travail féminin sans conciliation famille/emploi suffisante, évolutions culturelles (individualisme, choix de vie). Les allocations agissent surtout sur le \emph{calendrier} des naissances (on avance une naissance prévue), peu sur le nombre final d'enfants ; l'élasticité mesurée est faible (+0,05 à +0,15 enfant par point de PIB dépensé). Seuls les pays combinant crèches, congés parentaux partagés et égalité professionnelle (Suède, France) stabilisent l'ICF vers 1,7-1,9, sans atteindre durablement 2,1.
\end{corrige}

\begin{corrige}[Exercice 8 — Croquis commenté : population du Cameroun]
\textbf{Corrigé type du croquis} (l'élève reproduit ces éléments sur le fond fourni) :
\begin{itemize}
    \item \textbf{Densités} : Très forte ($>$100 hab/km²) : région de l'Ouest (hautes terres bamiléké) et Littoral (Douala) ; Forte (50--100) : Nord-Ouest, Sud-Ouest, Centre (Yaoundé) ; Moyenne (20--50) : Sud, Nord (axe Garoua) ; Faible (5--20) : Adamaoua, Extrême-Nord hors axes ; Très faible ($<$5) : Est (forêt équatoriale).
    \item \textbf{Villes} (ronds proportionnels) : Douala ($\approx$ 3,9 M) et Yaoundé ($\approx$ 4,3 M) en très gros ronds ; Garoua, Bamenda, Maroua en ronds moyens.
    \item \textbf{Flux} : flèche large Nord $\rightarrow$ Sud (exode rural et recherche d'emplois) ; flèche Ouest $\rightarrow$ Littoral/Centre (migrations vers Douala et Yaoundé).
    \item \textbf{Pressions} : hachures ou symboles « éclair » sur l'Adamaoua (conflits agriculteurs-éleveurs), le Nord-Ouest et l'Extrême-Nord (pression foncière, insécurité).
\end{itemize}
\textbf{Analyse commentée (corrigé type, 10 lignes) :} La population camerounaise est inégalement répartie : les deux tiers vivent sur moins de la moitié du territoire. Les densités maximales se concentrent dans l'Ouest (terres volcaniques fertiles, tradition d'habitat dense) et dans le binôme Douala-Yaoundé, qui polarise l'économie et attire les flux migratoires internes (macrocéphalie urbaine). À l'inverse, l'Est forestier reste quasi vide (enclavement, milieu équatorial). Les migrations Nord $\rightarrow$ Sud et Ouest $\rightarrow$ Littoral traduisent la recherche d'emplois et de services, mais alimentent une urbanisation rapide et mal maîtrisée (habitat précaire, chômage). Dans l'Adamaoua et l'Extrême-Nord, la croissance démographique rapide (ICF $>$ 6) accentue la pression sur les terres et les ressources en eau, source de conflits agro-pastoraux. Cette répartition contrastée pose un défi majeur d'aménagement du territoire pour 2035.
\end{corrige}

\begin{attention}[Barème indicatif du croquis — Exercice 8 (5 pts)]
Titre précis (0,5) ; légende organisée en items (1) ; 5 zones de densité correctement placées (1,5) ; villes + symboles proportionnels (0,5) ; 2 flux fléchés (0,5) ; nord + échelle (0,5) ; analyse commentée cohérente (0,5). Un croquis sans titre ni légende est plafonné à la moitié des points.
\end{attention}

\begin{corrige}[Exercice 9 — Sujet type Probatoire : problèmes démographiques contrastés]
\textbf{1. Analyse des documents (6 pts)}
\begin{enumerate}
    \item Forte croissance : Afrique subsaharienne (TAN $>$ 2,5\,\%, ex. Niger, RDC) et Asie du Sud (Inde, Pakistan). Croissance faible/négative : Europe (Allemagne, Italie ; TAN parfois négatif) et Asie de l'Est (Japon, Chine).
    \item Niger : base très large (49,8\,\% de moins de 15 ans), sommet pointu $\Rightarrow$ population jeune, explosion scolaire, forte charge de dépendance des jeunes, besoins en emplois. Japon : base étroite, sommet large ($\approx$ 29\,\% de 65 ans et plus) $\Rightarrow$ vieillissement, déficit de main-d'œuvre, financement des retraites et de la dépendance.
    \item Le Doc 3 montre une corrélation inverse : plus le PIB/hab et la scolarisation des filles sont élevés, plus l'ICF est bas (France : ICF 1,7 ; Niger : ICF 6,7 avec faible alphabétisation féminine). L'éducation des femmes retarde l'âge au mariage, diffuse la contraception et réduit la fécondité.
\end{enumerate}
\textbf{2. Mise en relation (9 pts)}
\begin{enumerate}
    \item Les pays n'ont pas entamé ni achevé la transition démographique au même moment : l'Europe a commencé au XIX\textsuperscript{e} siècle et est en phase 4/5 (natalité et mortalité faibles, croissance nulle ou négative, vieillissement) ; l'Afrique subsaharienne est en phase 2 (mortalité en baisse, natalité encore forte, croissance rapide). D'où des problèmes opposés : surcharge des systèmes éducatifs et du marché du travail au Sud ; pénurie d'actifs et coût des aînés au Nord.
    \item \textbf{Anti-nataliste} : politique de l'enfant unique en Chine (1979-2015) — efficace quantitativement (environ 400 millions de naissances évitées selon Pékin) mais coercitive, avec des effets pervers (déséquilibre des sexes, vieillissement rapide) qui ont conduit à son abandon. \textbf{Pro-nataliste} : politique familiale française (allocations, crèches, congé parental) — elle maintient l'ICF à 1,7-1,8, plus élevé que la moyenne européenne, mais sans atteindre le seuil de remplacement (2,1) : efficacité réelle mais limitée.
    \item Discussion : la croissance concentrée dans les pays pauvres est un \textbf{frein} à court terme (charge de dépendance élevée, pression sur écoles, santé, emplois, ressources). Mais elle peut devenir une \textbf{opportunité} : si la fécondité baisse et que les investissements en éducation, santé et emplois suivent, le pays capte un \cle{dividende démographique} (ex. « tigres » asiatiques). Sans ces investissements, le risque est une « malédiction démographique » (chômage des jeunes, migrations forcées, instabilité).
\end{enumerate}
\textbf{3. Croquis mondial (5 pts) :} planisphère simplifié ; hachures sur l'Afrique subsaharienne et l'Asie du Sud (croissance forte/jeunesse) ; pointillés sur l'Europe, le Japon, la Chine (vieillissement/déclin) ; flèches migratoires majeures : Afrique $\rightarrow$ Europe, Amérique latine $\rightarrow$ Amérique du Nord, Asie du Sud $\rightarrow$ Golfe. Titre, légende en 3 items, orientation.
\end{corrige}

\begin{attention}[Barème indicatif et vigilance — Exercice 9]
Analyse : 6 pts (2 par question) ; mise en relation : 9 pts (3 par question) ; croquis : 5 pts. Vigilance : citer les \emph{phases} de la transition pour justifier le contraste ; nommer au moins un pays par politique ; le croquis sans légende organisée perd la moitié des points.
\end{attention}

\begin{corrige}[Exercice 10 — Sujet type Baccalauréat Technique : dividende démographique 2035]
\textbf{Corrigé structuré (plan imposé respecté) :}

\textbf{Introduction.} Accroche : avec 41,5\,\% de moins de 15 ans et 300 000 jeunes entrant chaque année sur le marché du travail, le Cameroun est un des pays les plus jeunes du monde. Définition : le dividende démographique est le gain de croissance potentiel lié à l'augmentation de la part des actifs par rapport aux inactifs lors de la transition. Problématique : ce potentiel peut-il être transformé en moteur de l'émergence visée en 2035 ? Plan : potentiel réel mais fenêtre étroite (I) ; goulots d'étranglement (II) ; réponses publiques insuffisantes (III).

\textbf{I. Un potentiel indéniable, une fenêtre étroite.} La pyramide camerounaise à base large (ICF 4,5, TAN 2,6\,\%) annonce une entrée massive d'actifs d'ici 2035 ; si l'ICF descend vers 3,5-4,0, la charge de dépendance jeune (actuellement $\approx$ 82\,\%) diminuera, libérant épargne et investissement. Mais la condition \emph{sine qua non} est l'accélération de la baisse de la fécondité via l'éducation des filles et le planning familial (besoins non satisfaits $\approx$ 18\,\%).

\textbf{II. Les goulots d'étranglement.} Éducation : scolarisation primaire élevée ($\approx$ 95\,\%) mais achèvement faible et secondaire à $\approx$ 50\,\%, classes surchargées, inadéquation formation/emploi. Emploi : moins de 50 000 emplois formels créés par an face à 300 000 entrants ; sous-emploi $\approx$ 70\,\%, informel dominant. Santé : mortalité maternelle de 438/100 000 naissances vivantes, dépenses de santé $\approx$ 4\,\% du PIB, contraception inégalement accessible. Disparités : « deux Camerouns » — Nord/Extrême-Nord (ICF $>$ 6, pauvreté, insécurité) contre Sud/Centre urbanisé et éduqué (ICF $\approx$ 3,5) ; urbanisation à 58\,\% concentrée sur Douala-Yaoundé (macrocéphalie).

\textbf{III. Les réponses publiques : ambitions vs mise en œuvre.} Cadres : Vision 2035, PND30/SND30, Politique de population (1991, révisée 2007). Mesures : gratuité du primaire, formation professionnelle, projets « Emploi Jeunes », Couverture Santé Universelle en déploiement. Limites : part du social $<$ 20\,\% du budget, gouvernance et corruption, crises sécuritaires (Nord-Ouest/Sud-Ouest, Extrême-Nord) qui détournent investissements et détruisent le capital humain.

\textbf{Conclusion.} Le Cameroun peut capter une partie de son dividende, la fenêtre restant ouverte jusqu'à 2040-2050 ; mais sans accélération de la transition et investissements massifs dans le capital humain, le risque est la « malédiction démographique ». Ouverture : rôle de la diaspora et des technologies numériques.
\end{corrige}

\begin{attention}[Barème indicatif — Exercice 10 (20 pts)]
Introduction (3) ; Partie I (4) ; Partie II (6) ; Partie III (4) ; Conclusion (2) ; qualité rédactionnelle et chiffres clés (1). Vigilance : définir le dividende en introduction ; mobiliser au moins 5 données chiffrées ; respecter le plan imposé ; nuancer la conclusion (ni optimisme béat ni catastrophisme).
\end{attention}

\begin{corrige}[Exercice 11 — Note de synthèse : allocation familiale majorée au 3e enfant]
\textbf{Corrigé type (structure et arguments attendus) :}

\textbf{Objet :} opportunité d'une allocation familiale universelle majorée à partir du 3\textsuperscript{e} enfant.

\textbf{Contexte :} le Cameroun compte $\approx$ 12 millions d'enfants de 0-14 ans (41,5\,\% de la population) ; l'ICF (4,5) reste très au-dessus du seuil de remplacement (2,1) et la fécondité non désirée est élevée (besoins de planning familial non satisfaits $\approx$ 18\,\%). Le pays est en phase 2/3 de transition : l'enjeu n'est pas de relancer la natalité mais d'accompagner sa baisse pour capter le dividende démographique.

\textbf{Analyse coûts-bénéfices :}
\begin{itemize}
    \item \emph{Budgétaire} : coût estimé 150-200 milliards FCFA/an, soit $\approx$ 10\,\% des dépenses sociales actuelles (1 800 milliards) et 0,3\,\% du PIB — soutenable en volume, mais en concurrence directe avec l'éducation et la santé, déjà sous-financées (santé $\approx$ 4\,\% du PIB).
    \item \emph{Démographique} : l'élasticité internationale (+0,05 à +0,15 enfant par point de PIB) indique un effet faible mais ici \emph{contre-productif} : majorer au 3\textsuperscript{e} enfant incite à des naissances supplémentaires alors que la priorité est la réduction de la fécondité non désirée. Cela retarderait la baisse de la charge de dépendance ($\approx$ 82\,\%) et donc le dividende.
    \item \emph{Social} : effet positif sur la pauvreté monétaire (37\,\%) si le versement est effectif ; mais l'universalisme verse aussi aux ménages aisés (équité verticale faible).
\end{itemize}

\textbf{Analyse d'équité :} risque d'exclusion des ménages ruraux non déclarés (état civil incomplet) ; la majoration au 3\textsuperscript{e} enfant défavorise les petites familles et peut renforcer les inégalités filles/garçons si l'allocation n'est pas conditionnée à la scolarisation.

\textbf{Recommandation (Contre, avec amendement) :} rejeter la majoration pro-nataliste au 3\textsuperscript{e} enfant ; lui substituer (1) une allocation \emph{plafonnée aux deux premiers enfants} et \emph{conditionnée} à la scolarisation (notamment des filles) et aux visites sanitaires ; (2) un renforcement du planning familial (coût bien moindre, effet démographique direct) ; (3) un ciblage progressif sur les ménages pauvres via le registre social. Ce paquet, à budget comparable, soutient les familles sans freiner la transition, et rapproche le Cameroun des conditions du dividende démographique à l'horizon 2035.
\end{corrige}

\begin{attention}[Points de vigilance — Exercice 11]
Ton administratif exigé (pas de dissertation) ; chaque argument doit être chiffré ; mobiliser explicitement les concepts : dividende, charge de dépendance, transition, équité horizontale/verticale ; la recommandation doit être tranchée et cohérente avec l'analyse.
\end{attention}

\newpage

\coursec{Fiche bilan}

\begin{popfigure}
\begin{tikzpicture}[mindmap, grow cyclic, every node/.style={concept, font=\small},
  root concept/.append style={concept color=popblue, font=\bfseries},
  level 1/.append style={level distance=3.4cm, sibling angle=60},
  level 2/.append style={level distance=2.2cm, font=\scriptsize}]
\node[root concept]{La population mondiale}
  child[concept color=poporange]{ node {Indicateurs du démographe}
    child { node {TN, TM, TGN} }
    child { node {Taux d'accroissement} }
  }
  child[concept color=popgreen]{ node {Transition démographique}
    child { node {4 phases} }
    child { node {Modèles contrastés} }
  }
  child[concept color=poppurple]{ node {Problèmes mondiaux}
    child { node {Surpopulation} }
    child { node {Vieillissement} }
    child { node {Migrations} }
  }
  child[concept color=poppink]{ node {Politiques démographiques}
    child { node {Natalistes} }
    child { node {Malthusiennes} }
    child { node {Acteurs : État, ONG} }
  }
  child[concept color=popgold]{ node {Focus Cameroun}
    child { node {Population jeune} }
    child { node {Vision 2035} }
  }
  child[concept color=popblueL]{ node {Méthodes}
    child { node {Cartes et croquis} }
    child { node {Graphiques} }
  };
\end{tikzpicture}
\legende{Carte mentale de la leçon : les six blocs de savoirs à maîtriser pour le Probatoire / Baccalauréat technique.}
\end{popfigure}

\begin{bilan}
\textbf{Définitions clés}
\begin{itemize}
  \item \cle{Taux de natalité} (TN) : nombre de naissances vivantes pour 1000 habitants en un an.
  \item \cle{Taux de mortalité} (TM) : nombre de décès pour 1000 habitants en un an.
  \item \cle{Taux d'accroissement naturel} : différence entre natalité et mortalité.
  \item \cle{Transition démographique} : passage d'un régime traditionnel (natalité et mortalité élevées) à un régime moderne (natalité et mortalité faibles).
  \item \cle{Politique nataliste} : ensemble de mesures visant à augmenter les naissances ; \cle{politique malthusienne} : mesures visant à limiter les naissances.
  \item \cle{Densité de population} : nombre d'habitants par km\textsuperscript{2}.
\end{itemize}

\textbf{Formules à connaître par c\oe ur}
\begin{center}
\begin{tabularx}{\linewidth}{|l|X|}\hline
\rowcolor{popblueL} \textbf{Indicateur} & \textbf{Formule} \\ \hline
Taux de natalité & TN $= \dfrac{\text{naissances}}{\text{population totale}} \times 1000$ \\ \hline
Taux de mortalité & TM $= \dfrac{\text{décès}}{\text{population totale}} \times 1000$ \\ \hline
Accroissement naturel & TGN $=$ TN $-$ TM (en pour mille) \\ \hline
Accroissement total & TGN $+$ solde migratoire \\ \hline
Densité & $d = \dfrac{\text{population}}{\text{superficie en km}^2}$ \\ \hline
\end{tabularx}
\end{center}

\textbf{Méthodes express}
\begin{itemize}
  \item \textbf{Lire une pyramide des âges} : base large = population jeune (Sud) ; sommet large = vieillissement (Nord).
  \item \textbf{Tracer un croquis} : titre, légende numérotée, flèche du nord, échelle indicative.
  \item \textbf{Analyser un document} : 1) identifier nature, auteur, date ; 2) dégager les données chiffrées ; 3) calculer un taux ou une part en \% ; 4) conclure en reliant au problème démographique.
  \item \textbf{Comparer deux pays} : toujours utiliser les mêmes indicateurs (TN, TM, TGN, espérance de vie) et justifier par la phase de transition.
\end{itemize}
\end{bilan}

\begin{aretenir}[Checklist avant le Bac]
\begin{itemize}
  \item $\square$ Je sais définir TN, TM, TGN, densité et transition démographique sans hésiter.
  \item $\square$ Je sais calculer un taux d'accroissement naturel à partir de données brutes.
  \item $\square$ Je connais les 4 phases de la transition démographique et un pays exemple pour chacune.
  \item $\square$ Je sais distinguer politique nataliste (France) et politique malthusienne (Chine, Inde).
  \item $\square$ Je connais les acteurs des politiques démographiques : État, ONG, ONU (UNFPA).
  \item $\square$ Je sais citer les grands problèmes démographiques : surpopulation, vieillissement, migrations, urbanisation.
  \item $\square$ Je connais les traits de la population camerounaise : jeunesse, croissance rapide, répartition inégale.
  \item $\square$ Je sais relier les défis démographiques du Cameroun à la Vision 2035.
  \item $\square$ Je sais lire et commenter une pyramide des âges.
  \item $\square$ Je sais construire un croquis avec titre, légende, nord et échelle.
  \item $\square$ Je sais rédiger une conclusion justifiée à partir d'un document statistique.
  \item $\square$ Je vérifie toujours mes unités (pour mille, \%, habitants/km\textsuperscript{2}).
\end{itemize}
\end{aretenir}

\findecours

\end{document}
